%!TEX root = ./../../springer-mono-format/book.tex
%!TEX root = ./../../thesis-trier.tex

In~\cref{constr_sync:sec:prelim_results} we show that \constrsync{L} is decidable and in \PSPACE.
Furthermore, we give example constraint languages
yielding constrained synchronization 
problems with computational complexities 
in $\AC^0$, or \LOGSPACE-, \NL-, \NP-, or \PSPACE-complete.

In~\cref{constr_sync:sec:manip} we present results on altering and simplifying 
the constraint language and how certain operations as forming factor sets, unions, intersections, morphisms, and concatenation
affect the computational complexity.
This also yields
general methods to deduce complexity lower bounds between different instances of constrained synchronization problems.

General hardness conditions for \NL, \NP, and \PSPACE\  are given in \cref{constr_sync:sec:hardness}.
This section also develops a general framework to deduce these hardness results introducing
the notion of auxiliary state avoiding words and the introduction of relations derived from the Myhill-Nerode relation.

In~\cref{constr_sync:sec:aut} we discuss the representation of the constraint by PDFAs.
We show that we have some freedom in the choice of a start state and a set of final states.

\section{Basic Complexity Results and Examples}
\label{constr_sync:sec:prelim_results}

%!TEX root = ./../../springer-mono-format/book.tex
%!TEX root = ./../../thesis-trier.tex

Here, \cref{constr_sync:thm:PSPACE} shows that \constrsync{L} 
is decidable and solvable in polynomial space for regular $L \subseteq \Sigma^*$.
Then, \cref{constr_sync:thm:introductory_examples} 
presents various example problems and their complexities.

\begin{theorem}
	\label{constr_sync:thm:PSPACE}
	Let $L \subseteq \Sigma^*$ be regular.
	Then \constrsync{L} is decidable and in \PSPACE.
\end{theorem}
\begin{proof}
	An input semiautomaton $\semiautomat{A}{Q}{\Sigma}{\delta}$
	admits a synchronizing word in~$L$, i.e. $L \cap \Syn(\mathcal A) \ne \emptyset$,
	if and only if there exists $t \in Q$ such that
	\[  
	L \cap \bigcap_{q \in Q} L(\mathcal A_{q, \{t\}}) \ne \emptyset.
	\]
	Let $\mathcal B$ be a complete automaton 
	with $L = L(\mathcal B)$, which can be chosen in advance
	since~$L$ is fixed by the definition of the problem.
	
	An algorithm solving the problem
	can guess the state~$t$
	and then call a procedure
	to solve \DFAINT for $\mathcal B$
	and the automata 
	$\mathcal A_{q, \{t\}}$ for $q \in Q$
	as input (so that it is called with $|Q|+1$ automata as arguments).
	By \cref{prelim:thm:complexities},
	\DFAINT is in \PSPACE.
	So this algorithm can run in nondeterministic polynomial space.
	Since $\PSPACE = \NPSPACE$ 
	by \cref{prelim:eqn:cc_inclusions}
	we obtain the claim.
\end{proof}

The following introductory examples (one for each complexity class that will come up in this part)
are presented here to motivate the later results and to have some easy and concrete cases
with self-contained and independent proofs. Each example will be generalized in the coming sections.

\begin{theorem}
	\label{constr_sync:thm:introductory_examples}
	Let $a,b,c \in \Sigma$ be three distinct letters.
	\begin{enumerate}
		
		\item \label{constrsync:ex:item:a} $\constrsync{\{a\}}$ is in $\AC^0$.
		
		\item \label{constrsync:ex:item:bstar} $\constrsync{b^*}$ and $\constrsync{ab^*}$ are \LOGSPACE-complete
	    even for trap-state input automata.
		
		\item \label{constrsync:ex:item:ab} $\constrsync{(a+b)^*}$ and $\constrsync{\Sigma^*}$ are \NL-complete. 
	
		\item \label{constrsync:ex:item:abstara} $\constrsync{ab^*a}$ is \NP-complete even for input automata with a trap state~\cite{FernauGHHVW19}.
		
		\item \label{constrsync:ex:item:abcL} $\constrsync{(a+b)^*c^lL}$ is \PSPACE-complete for 
		non-empty and regular $L \subseteq \Sigma^*$ such that $|\fact{L} \cap \{a^n, b^m, c^l \}| \le 2$
		for some $n,m, l \ge 1$ even for trap-state automata as input.
		
		\item \label{constrsync:ex:item:cab} $\constrsync{c(a+b)^*}$ is \PSPACE-complete even for strongly connected automata as input~\cite{FernauGHHVW19}.	
	\end{enumerate}
\end{theorem}
\begin{proof}
	\ref{constrsync:ex:item:a}: 
	The letter $a$ synchronizes a given input automaton
	if and only if the transition matrix $A = (a_{ij})_{1 \le i,j \le n}$ 
	of $a$ contains a column of ones.
	Equivalently, there exists $J \in \{1,\ldots,n\}$
	such that for all $i \in \{1,\ldots,n\}$ we have 
	$a_{iJ} = 1$.
	This can be checked by an~$\AC^0$ circuit implementing the formula 
	\[
	\bigvee_{J \in \{1,\ldots,n\}} \left( \bigwedge_{i \in \{1,\ldots,n\}} a_{iJ} \right).
	\]
	In \cref{sec:constr_sync:finite} (\cref{constr_sync:thm:finite}) 
	this is generalized to finite constraint languages.
		
	\ref{constrsync:ex:item:bstar}:  
	In \cref{thm:SYNC-solv} we have shown that \DFAST for unary
	automata is \LOGSPACE-complete. This problem is equivalent
	to $\constrsync{b^*}$.			
	By adding self-loops to each state for the letter $a$
	in an input automaton $\semiautomat{A}{Q}{\{b\}}{\delta}$
	over~$\{b\}$, we obtain 
	an $\leACZEROm$-reduction from this problem. So, $\constrsync{b^*} \le_m^{\text{AC}^0} \constrsync{ab^*}$.	
	Next, we get an $\leACZEROm$-reduction
	from $\constrsync{ab^*}$
	to \SFS over the unary alphabet $\{b\}$ 
	in the following way: for $\semiautomat{A}{Q}{\{a,b\}}{\delta}$
	construct the \SFS-instance $(\mathcal A_{|\{b\}}, S)$
	with $S = \applytrans{\delta}{Q}{a}$
	where w.r.t.~$Q = \{q_1, \ldots, q_n\}$
	a characteristic bit-vector of $S$ can 
	be computed by an~$\AC^0$ circuit 
	as, for $[a] = (m_{ij}^{(a)})_{i,j\in \{1,\ldots,n\}}$,
	we have $q_j \in S$ if and only if $m_{ij}^{(a)} = 1$ for some $i \in \{1,\ldots,n\}$.	
	Then the claim follows as \SFS over $\{b\}$ is in \LOGSPACE
	by \cref{thm:SF-ST-comm}. 
		
	\ref{constrsync:ex:item:ab}: $\constrsync{(a+b)^*}$ and $\constrsync{\Sigma^*}$ are
	equivalent to the unconstrained synchronization problem
	over $\{a,b\}$ and~$\Sigma$, which are \NL-complete by Theorem~\ref{prelim:complexity:general_aut_sync}
	since $|\Sigma| \ge 2$.
	
	\ref{constrsync:ex:item:abstara}:
	We give a reduction from \DST
	over a unary alphabet, which is \NP-complete by Theorem~\ref{compl:thm:comm_ST}.
	Let $(\mathcal A, S, T)$ be an instance
	with $\semiautomat{A}{Q}{\{b\}}{\delta}$ and $S, T \subseteq Q$
	such that $S \cap T = \emptyset$.
	Pick some $\hat s \in S$
	and let $\overline{S} \colon ( Q \setminus T ) \to S$
	be given by $\overline{S}(q) = \hat{s}$ if $q \in Q \setminus (S \cup T)$
	and $\overline{S}(q) = q$ otherwise. Construct
	$\semiautomat{A'}{Q \cup \{t\}}{\{a,b\}}{\delta'}$ with a new state $t \notin Q$ 
	where, for $q \in Q$ and $x \in \{a,b\}$,
	\[
	\applytrans{\delta'}{q}{x} = \begin{cases}
		\applytrans{\delta}{q}{b} & \text{if } x = b, \\
		\overline{S}(q)           & \text{if } x = a \text{ and } q \in Q \setminus T, \\
		t                         & \text{if } x = a \text{ and } q \in T
	\end{cases}
	\]
	and $\applytrans{\delta'}{t}{x} = t$ for all $x \in \{a,b\}$. 
	The state $t$ must be the unique synchronizing state since it is a trap state.
	Furthermore, for $m \ge 0$,     
	we have $\applytrans{\delta}{S}{b^m} \subseteq T$
	if and only if $\applytrans{\delta'}{Q\cup \{t\}}{ab^ma} = \{t\}$,
	which yields that this is a valid reduction.
	
	In \cref{constr_sync:sec:motivation} we have already established that
	$\constrsync{ab^*a}$ is in \NP\ 
	by using the matrix orbit problem.
	For a different argument, given $\semiautomat{A}{Q}{\Sigma}{\delta}$,
	by \cref{prelim:word_in_PA} applied to $S = \delta(Q, a)$
	there exists a word in $ab^*a$ synchronizing~$\mathcal A$
	if and only if there exists $0 \le i < 2^{|Q|} - 1$
	such that $|\applytrans{\delta}{Q}{ab^ia}| = 1$.
	Using a binary encoding, $i$ can be guessed
	in polynomial time. By a square-and-multiply
	algorithm the transition matrix of $b^i$ can be computed in polynomial time. Then,
	the condition $|\applytrans{\delta}{Q}{ab^ia}| = 1$  
	can be checked in polynomial time.
		
	\ref{constrsync:ex:item:abcL}:
	By \cref{constr_sync:thm:PSPACE} we only have to establish \PSPACE-hardness.
	We give a reduction from \SIS over a binary alphabet, which
	is \PSPACE-complete by Theorem~\ref{prelim:thm:binary_pspace_problems}. 
	Let $(\mathcal A, S)$ be an instance of \SIS
	where $\semiautomat{A}{Q}{\{a_1, a_2\}}{\delta}$ 
	and $S \subseteq Q$. 		
			
	Let $u_1 = a^{n}b^{m}$, $u_2 = a^{n+1}b^{m}$.
	Intuitively, we re-encode the letters $a_1$ and $a_2$
	into words $u_1$ and~$u_2$, respectively, and build a new automaton~$\mathcal A'$
	containing a copy of~$Q$ in its state set 
	that branches along these two words in the same way as~$\mathcal A$
	branches on $a_1$ and~$a_2$.
	Additionally, the new automaton uses $c^l$
	to map states from~$S$ into a trap state.
	A word from the constraint language $(a+b)^*cL$
	synchronizes $\mathcal A'$
	if and only if, before reading $c^l$, we land
	inside~$S$. This follows since the suffix of this word following $c^l$ only uses letters
	from $\{u_1,u_2\}$, $\{u_1,c^l\}$, or $\{u_2,c^l\}$,
	but the construction is done in a way such that 
	we need all three words to make any changes to the current state set
	yielding a new synchronizing word.
	
	Set $P = \pref{\{a^n b^{m-1}, a^{n+1}b^{m-1}, c^{l-1} \}}$.
	Let $\semiautomat{A'}{ ( Q \times P ) \cup \{t\}}{\{a,b,c\}}{\delta'}$
	where, for $q \in Q$, $u \in P$,
	and $x \in \{a,b,c\}$, we set
	\[
	\delta'((q, u), x) =
	\begin{cases}
		(q, ux) & \text{if } ux \in P, \\
		(\delta(q, a_i), \varepsilon) & \text{if $ux = u_i$ for some $i \in \{1,2\}$,} \\	
		(q, \varepsilon) & \text{if $ux = c^l$ and $q \notin S$}, \\
		t                & \text{if $ux = c^l$ and $q \in S$},\\
		(q, \sigma_P(ux)) & \text{otherwise, i.e., if $ux \notin P \cup \{u_1, u_2, c^l\}$,}
	\end{cases}
	\]
	and $\delta'(t, x) = t$ for all $x \in \{a,b,c\}$.\footnote{Note that 
		for $u \in P$ and $x \in \{a,b,c\}$
		with $ux \notin P$
		we have $\sigma_P(ux) = x$ expect in the case $u = a^{n+1}$ and $x = a$.}

	\Cref{constrsync:fig:reduction} shows a part of the reduction
	for a single state~$q \in Q$ of the input automaton
	and $n = m = 2$.
	 The additional ``auxiliary'' states, used to track the current prefix and to branch when $abb$ or $aabb$ has been read, are indicated by rectangles. Transitions labeled~$c$ are omitted for clarity: they either return to $(q,\varepsilon)$ if $q \notin S$ or lead to the trap state~$t$. 
	 Along such a $c$-path $(q, c^i)$ with $1 \le i \le l-2$, 
	 reading $a$ or $b$ leads to $(q,\varepsilon)$.
	
		\begin{figure}[htb]
			\centering
			\begin{tikzpicture}				
				\node (q) {$(q, \varepsilon)$};
				\node[draw,rectangle, below = of q] (qa) {$(q, a)$};
				\node[draw,rectangle, below = of qa] (qaa) {$(q, aa)$};
				\node[draw,rectangle, below left = of qa] (qab) {$(q, ab)$};
				\node[below = of qab] (qabb) {$(\delta(q, a_1),\varepsilon)$};
				\node[draw, rectangle, below right = of qaa] (qaab) {$(q, aab)$};
				\node[below = of qaab] (qaabb) {$(\delta(q, a_2), \varepsilon)$};
				
				\node[right = of q, yshift=1em] (q-dummy) {$\ldots$};
				\node[right = of qa, yshift=1em] (qa-dummy) {$\ldots$};
				\node[right = of qaa, yshift=1.7em] (qaa-dummy) {$\ldots$};
				\node[right = of qaab, yshift=1em] (qaab-dummy) {$\ldots$};
				\node[left = of qab, yshift=1em] (qab-dummy) {$\ldots$};
				\node[left = of qabb, yshift=1em] (qabb-dummy) {$\ldots$};
				\node[right = of qaabb, yshift=1em] (qaabb-dummy) {$\ldots$};
				
				\path[->] (q) edge [right] node {$a$} (qa)
				(qa) edge [right] node {$a$} (qaa)
				(qaa) edge [below] node {$b$} (qaab);
				
				\path[->] (qa) edge [above] node {$b$} (qab)
				(qab) edge [left] node {$b$} (qabb);
				
				\path[->] (qaab) edge [right] node {$b$} (qaabb);
				
				\path[->] (qaab) edge [bend right=68,right] node {$a$} (q)
				(qab) edge [bend left,left] node {$a$} (q);
				\path[->] (qaa) edge [loop right] node {$a$} (qaa);
				\path[->] (q) edge [loop left] node {$b$} (q);
				
				\path[->] (q) edge [above] node {$c$} (q-dummy);
				\path[->] (qa) edge [above] node {$c$} (qa-dummy);
				\path[->] (qaa) edge [above] node {$c$} (qaa-dummy);
				\path[->] (qab) edge [above] node {$c$} (qab-dummy);
				\path[->] (qaab) edge [above] node {$c$} (qaab-dummy);
				\path[->] (qaabb) edge [above] node {$c$} (qaabb-dummy);
				\path[->] (qabb) edge [above] node {$c$} (qabb-dummy);				
			\end{tikzpicture}
			\caption[Construction for hardness of constraint language $\{a,b\}^*c^lL$]
			{Some states for an example of the construction in \cref{constrsync:ex:item:abcL} with $n = m = 2$,
				see the proof of \cref{constrsync:ex:item:abcL} for an explanation.}
			\label{constrsync:fig:reduction}
		\end{figure}
	
	Next, we show that $(\mathcal A, S)$ is a \rYES-instance of \SIS 
	if and only if $\Syn(\mathcal A')$ contains a word from the constraint language.
			
	First, suppose $\delta(Q, w) \subseteq S$ for some $w \in \{a_1,a_2\}^*$.
	We have 
	\[
	\delta'(( Q \times P ) \cup \{t\}, b^{m+1}) = Q \times \{\varepsilon\} \cup \{t\}.
	\]
	Let $\varphi \colon \{a_1,a_2\}^* \to \{a,b,c\}^*$
	be the morphism given by $\varphi(a_i) = u_i$ for $i \in \{1,2\}$.
	Since $\delta'((q, \varepsilon), u_i) = (\delta(q, a_i), \varepsilon)$
	for $i \in \{1,2\}$ and $q \in Q$,
	we find $\delta'((q, \varepsilon), \varphi(w)) = (\delta(q, w), \varepsilon)$
	for all $q \in Q$
	and thus $\delta'(Q \times \{\varepsilon\}, \varphi(w)) \subseteq S \times \{\varepsilon\}$.
	For $q \in S$
	we have $\delta'((q, \varepsilon), c^l) = t$.
	Since~$L$ is non-empty, we can pick some word $v \in L$.
	Combining everything and since~$t$ is a trap state,
	we find that 	
	$\delta'((Q \times P) \cup \{t\}, b^{m+1} \varphi(w) c^l v) = \{t\}$.
		
	Conversely, suppose we have $|\delta'(Q \times P \cup \{t\}, wc^lv)| = 1$
	for some $w \in \{a,b\}^*$ and $v \in L$.
	Since~$t$ is a trap state, we find $\delta'(Q \times P \cup \{t\}, wc^lv) = \{t\}$.
	Let $q \in Q$.
	We have $\delta'((q, \varepsilon), w) = (p, u)$
	for some $u \in \pref{\{u_1, u_2\}}$ and $p \in Q$. Assume $p \notin S$.
	Then $\delta'((p, u), c^l) = (p, \varepsilon)$ (this holds for $u = \varepsilon$ and 
	for $u \ne \varepsilon$).	
	Next, to map $(p, \varepsilon)$ to~$t$, we must first map it to a state
	in $S \times P$ and then apply $c^l$.	
	Now, for $w' \in \{a,b,c\}^*$
	to have $\delta'((p, \varepsilon), w') \in S \times P$
	it must have $u_1$ or $u_2$ as a factor,
	and to map a state from~$S \times P$
	to~$t$ it must have $c^l$ as a factor.
	Since by assumption $|\fact{L} \cap \{a^n, b^m, c^l\}| \le 2$,
	we have $c^l \notin \fact{L}$
	or $|\fact{L} \cap \{a^n, b^m\}| \le 1$.
	where the latter case yields $u_1, u_1 \notin \fact{L}$.
	Thus there is no way to map $(p, \varepsilon)$
	to~$t$ with a word from~$L$	and so we must have $p \in S$.
	Now, we have $\delta'((q, \varepsilon), w) = (p, u)$
	if and only if
	there exist words $v_0, v_1, \cdots, v_r$ 
	and $w_1, w_2, \cdots w_r \in \{u_1, u_2\}$
	such that
	\begin{enumerate}
	\item $w = v_0 w_1 v_1 w_2 \cdots w_r v_r$,
		
	\item $v_r$ has $u$ as a suffix, and
	
	\item $\delta'((q, \varepsilon), w_1 \cdots w_r) = (p, \varepsilon)$. 	
	\end{enumerate}
	Let $x_i \in \{a_1, a_2\}$ for $i \in\{1,\ldots,r\}$
	with $x_i = a_1$ if $w_i = u_1$
	and $x_i = a_2$ otherwise.
	Then $\delta(q, x_1 \cdots x_r) = p$
	and thus $\delta(Q, x_1 \cdots x_r) \subseteq S$.
	
	
	\ref{constrsync:ex:item:cab}: By \cref{constr_sync:thm:PSPACE} we only have to establish \PSPACE-hardness.
	We give a reduction from \SFS over a binary alphabet, which is \PSPACE-complete
	by \cref{prelim:thm:binary_pspace_problems} even for strongly connected automata.
	Let $(\mathcal A, S)$ be
	an \SFS-instance with $\semiautomat{A}{Q}{\{a,b\}}{\delta}$
	and $S \subseteq Q$. Pick some $\hat s \in S$.
	By extending~$\delta$ with $\applytrans{\delta}{q}{c} = q$ for each $q \in S$
	and $\applytrans{\delta}{q}{c} = \hat s$ for $q \notin S$
	we obtain an automaton over $\{a,b,c\}$ with 
	$\applytrans{\delta}{Q}{c} = S$
	that has a synchronizing word of the form~$cu$ for $u \in \{a,b\}^*$
	if and only if $|\applytrans{\delta}{S}{u}| = 1$.	
\end{proof}

\section{Manipulating Constraint Languages}
\label{constr_sync:sec:manip}

%!TEX root = ./../../springer-mono-format/book.tex
%!TEX root = ../../thesis-trier.tex

Here, we investigate how 
applying operations to 
constraint languages
affects the computational complexities of resulting constrained problems. Specifically, we investigate the following operations: 
the prefix/suffix/factor operations (\SEC~\ref{constr_sync:subsec:factor}),
set union (\cref{constr_sync:subsec:union}),
intersection with a regular
ideal language (Section~\ref{constr_sync:subsec:ideal}), applying morphism (Section~\ref{constr_sync:subsec:morphisms}),
and concatenating specific languages to an arbitrary constraint language (Section~\ref{constr_sync:subsec:concat}).
The operations from \cref{constr_sync:subsec:factor} 
preserve the complexity,
those from \cref{constr_sync:subsec:union,constr_sync:subsec:ideal,constr_sync:subsec:morphisms}
mostly decrease the complexity,
and those from \cref{constr_sync:subsec:concat}
increase the complexity.

\subsection{Factor Sets and Constraint Languages}
\label{constr_sync:subsec:factor}

%!TEX root = ./../../../thesis-trier.tex

Our first result is a consequence of Theorem~\ref{prelim:thm:sync_words_is_ideal}. Although it has a simple proof,
it is quite important for many of the following results.

\begin{theorem}
	\label{constr_sync:thm:L_in_FactLbar}
	Let $L, L' \subseteq \Sigma^*$ with $L \subseteq \factor{L'}$.
	If $\semiautomat{A}{Q}{\Sigma}{\delta}$ 
	has a synchronizing word in $L$, then it has one in $L'$.
\end{theorem}
\begin{proof}
	Suppose $u \in L$ synchronizes $\mathcal A$.
	By assumption there exist $v, w \in \Sigma^*$
	such that $vuw \in L'$.
	By Theorem~\ref{prelim:thm:sync_words_is_ideal} the
	latter word synchronizes $\mathcal A$ as well.
\end{proof}

Note that $L \subseteq \factor{L'}$ is equivalent to $\factor{L} \subseteq \factor{L'}$.
%As $L \subseteq L'$ implies $\factor{L} \subseteq \factor{L'}$
%and $\factor{\factor{L}} = \factor{L}$, the next is implied.
This gives the following consequence of Theorem~\ref{constr_sync:thm:L_in_FactLbar}.

\begin{corollary}
	\label{constr_sync:cor:factors_sets_included}
	Let $L, L' \subseteq \Sigma^*$ with $\factor{L} = \factor{L'}$.
	Then every automaton $\mathcal A$ has a synchronizing word in $L$
	if and only if it has one in~$L'$.
\end{corollary}

For any $L \subseteq \Sigma^*$ we have
$\factor{L} = \factor{\pref{L}} = \factor{\suff{L}}$. So with Corollary~\ref{constr_sync:cor:factors_sets_included},
we conclude that a given automaton has a synchronizing word in one of
the languages $L$, $\factor{L}$, $\pref{L}$, or $\suff{L}$ if and only if it has one in all of these languages.
As a consequence, we obtain that the corresponding constrained synchronization problems
are equivalent for $\leACZEROm$-reductions (in fact,
the input can be passed through without any modifications
and so these problems are also equivalent
for reductions by projections as considered in~\cite{valiant1982rfducibility}).
%, but $\leACZEROm$-reductions is sufficient for our applications). 

\begin{corollary}
	\label{gen_results:cor:factors_etc_complexity}
	Let $L \subseteq \Sigma^*$ be a regular language. 
	The following problems are equivalent for~$\leACZEROm$-reductions.
	\begin{enumerate}
		\item $\constrsync{L}$
		\item $\constrsync{\factor{L}}$
		\item $\constrsync{\pref{L}}$
		\item $\constrsync{\suff{L}}$
	\end{enumerate}
	Furthermore, if for two languages $L, L'\subseteq \Sigma^*$ the factor, prefix, or suffix sets %closures
	are equal, then $\constrsync{L} \equiv_m^{\AC^0} \constrsync{L'}$.
\end{corollary}

\begin{corollary} \label{constrsync:genresults:cor:union-fact}
	Let $L, L' \subseteq \Sigma^*$ be regular languages.
	Then we have  $\constrsync{L \cup L'} \eqACZEROm \constrsync{\fact{L}\cup L'}$
\end{corollary}
\begin{proof}
	This follows from \cref{gen_results:cor:factors_etc_complexity}
	since $\fact{L \cup L'} = \fact{\fact{L}\cup L'}$.
\end{proof}

\begin{corollary} \label{constrsync:genresults:cor:left-right-ideal}
	Let $L \subseteq \Sigma^*$ be a regular language
	containing a non-empty left or right ideal language.
	Then we have $\constrsync{L} \eqACZEROm \constrsync{\Sigma^*}$.	
\end{corollary}
\begin{proof}
	Without loss of generality, 
	we have $I \subseteq L$
	for some non-empty left ideal language~$I \subseteq \Sigma^*$ (which might be non-regular).
	Pick some $u \in I$ and consider a word $v \in \Sigma^*$.
	Then $vu \in I$
	so that $v \in \fact{I}$.
	This gives $\Sigma^* \subseteq \fact{I} \subseteq \fact{L}$.
	Thus $\fact{L} = \Sigma^*$
	and the claim follows from \cref{gen_results:cor:factors_etc_complexity}.
\end{proof}

\begin{corollary}
	Let $I, J \subseteq \Sigma^*$ be a left and a right ideal, respectively,
	which are both non-empty.
	Then $\constrsync{I\cap J} \eqACZEROm \constrsync{\Sigma^*}$.
\end{corollary}
\begin{proof}
	By assumption, we have $u \in I$ and $v \in J$ for some $u,v \in \Sigma^*$.
	Let $w \in \Sigma^*$.
	Then we have $vwu \in I$ since~$I$ is a left ideal 
	and $vwu \in J$ since~$J$ is a right ideal.
	Thus $vwu \in I \cap J$.
	This gives $\Sigma^* \subseteq \fact{I \cap J}$
	and the claim follows from \cref{gen_results:cor:factors_etc_complexity}. 
\end{proof}

\subsection{The Union of Constraint Languages}
\label{constr_sync:subsec:union}

%!TEX root = ../../../thesis-trier.tex

Here, in~\cref{constr_sync:thm:union}, we state that certain complexity classes
are preserved when considering the union
of two constraint languages. \cref{constrsync:genresults:cor:subset}
gives further conditions
for complexity preservation w.r.t.~union
and concatenation.
% when the complexity is preserved 
%for unions and concatenation of languages.

\begin{theorem}
	\label{constr_sync:thm:union}
	Let $\mathcal C$ be one of the complexity classes $\AC^0$, $\NC^1$, \LOGSPACE, \NL, \PTIME or \NP.
	Suppose $\constrsync{L}$ and $\constrsync{L'}$ are both in $\mathcal C$ for regular languages $L,L' \subseteq \Sigma^*$.
	Then  the problem $\constrsync{L \cup L'}$ is in $C$.
\end{theorem}
\begin{proof}
%Given two computing mechanism (circuits or Turing machines)
%$M_1$ and $M_2$ for \constrsync
For \LOGSPACE, \NL, \PTIME, and \NP\ suppose
$M_1$ and $M_2$ are Turing machines deciding $\constrsync{L}$
and $\constrsync{L'}$, respectively, for an input automaton
by halting and writing~$1$ for a \rYES-instance
or~$0$ for a \rNO-instance.
Let $M$ be a Turing machine that, for an input automaton $\mathcal A$,
first simulates $M_1$ with input~$\mathcal A$ yielding output~$r_1$,
then $M_2$ with input $\mathcal A$ yielding output $r_2$, and finally outputs $\max\{r_1, r_2\}$,
which equals one if and only if $\mathcal A$ has a synchronizing word in $L \cup L'$.  
We can combine the results $r_1$ and $r_2$  (repositioning reading head, erasing the working tape while storing the output of $M_1$, compute
$\max\{r_1, r_2\}$) 
by operating deterministically and using only constant space;
and if~$M_1$ and~$M_2$ obey the convention to
position the working tape head next to the output, erase
everything else and  move the reading head back to the start,
the overhead
needed to simulate both machines and combine  their results is actually constant in time and space.
Hence, if $M_1$ and $M_2$ both work in deterministic logspace, nondeterministic logspace,
deterministic polynomial time, or nondeterministic polynomial time,
also $M$ works in deterministic logspace, nondeterministic logspace,
deterministic polynomial time, or nondeterministic polynomial time, respectively.

For $\AC^0$ or $\NC^1$, we can combine the outputs
of two circuits from respective circuit families in $\AC^0$ or $\NC^1$ for $\constrsync{L}$
and $\constrsync{L'}$ by an OR-gate
to obtain an \LOGSPACE-uniform circuit family in $\AC^0$ or $\NC^1$, respectively, 
for $\constrsync{L \cup L'}$.
\end{proof}

Next, we state a condition such that the complexity is preserved when taking the union.

\begin{corollary} \label{constrsync:genresults:cor:subset}
	Let $L, L' \subseteq \Sigma^*$ be regular languages
	such that $L' \subseteq \fact{L}$.
	Then we have $\constrsync{L} \eqACZEROm \constrsync{L\cup L'}$.
\end{corollary}
\begin{proof}
	By \cref{constrsync:genresults:cor:union-fact}
	we have $\constrsync{L\cup L'} \eqACZEROm \constrsync{\fact{L} \cup L'}$
	and the assumption gives $\fact{L} \cup L' = \fact{L}$.
	By  \cref{gen_results:cor:factors_etc_complexity}
	we have $\constrsync{\fact{L} \cup L'} \eqACZEROm \constrsync{L}$ (the claim
	also follows directly from \cref{gen_results:cor:factors_etc_complexity}
	by noting that $\fact{L \cup L'} = \fact{L} \cup \fact{L'} = \fact{L}$).
\end{proof}

By \cref{constrsync:genresults:cor:subset}, for a right or left ideal~$I$ that is regular
and any regular $L \subseteq \Sigma^*$
we have $\constrsync{I \cup L} \eqACZEROm \constrsync{\Sigma^*}$
since $\factor{I} = \Sigma^*$.

\begin{remark}
	Let $L = \{a,b\}^* c \{a,b\}^*$ for distinct $a,b,c\in \Sigma$. 
	By \crefitem{constr_sync:thm:introductory_examples}{constrsync:ex:item:abcL}, $\constrsync{L}$
	is \PSPACE-complete. We have $\Sigma^* \setminus L = \{ u \in \Sigma^* : |u|_c = 0 \text{ or } |u|_c \ge 0 \}$
	and $\fact{\Sigma^* \setminus L} = \{a,b,c\}^*$.
	By \crefitem{constr_sync:thm:introductory_examples}{constrsync:ex:item:ab}, the problem
	$\constrsync{\Sigma^* \setminus L}$ is \NL-complete.
	Thus, in general, the computational complexity of the constrained synchronizing problem
	is not preserved when taking the complement.
\end{remark} 

\subsection{The Ideal Intersection Theorem}
\label{constr_sync:subsec:ideal}

%!TEX root = ./../../../thesis-trier.tex

Here, we show the ideal intersection theorem (\cref{constr_sync:gen_results:thm:ideal_separation}).
Later, we mostly use this result
to deduce \PSPACE-hardness from known \PSPACE-hard problems, especially
for constraint languages written as unions (see Remark~\ref{constrsync:compl:ideal_intersection_for_unions} for details).

\begin{theorem}[Ideal Intersection Theorem]
	\label{constr_sync:gen_results:thm:ideal_separation}
	Let $I \subseteq \Sigma^*$ be a fixed ideal language that is regular.
	Suppose $L \subseteq \Sigma^*$ is any regular language,
	then
	$$
	\constrsync{L \cap I} \leACZEROm \constrsync{L}.
	$$
\end{theorem}
\begin{proof}
	By \cref{prelim:thm:sync_words_is_ideal}
	there exists an automaton $\mathcal A_I$
	such that $I = \Syn(\mathcal A_I)$.
	Let $\semiautomat{A}{Q}{\Sigma}{\delta}$ be an input to $\constrsync{L\cap I}$.
	By~\cref{prelim:eqn:product_aut_sync_words}, 
	it is $\Syn(\mathcal A) \cap I = \Syn(\mathcal A \times \mathcal A_I)$,
	which gives	
	\begin{equation}\label{constrsync:eqn:ideal-int}
	\Syn(\mathcal A) \cap (L \cap I) \ne \emptyset
	\Leftrightarrow 
    \Syn(\mathcal A \times \mathcal A_I) \cap L \ne \emptyset,
	\end{equation}
	i.e., $\mathcal A \times \mathcal A_I$ has a synchronizing
	word in $L$ if and only if $\mathcal A$ has a synchronizing word in $L \cap I$.
	As $I$ is fixed,
	we can construct the product automaton $\mathcal A \times \mathcal A_I$
	%for a fixed automaton $\mathcal A_I$ 
	by a circuit family in $\AC^0$ 
	 as we only have to add transitions
	and states to~$\mathcal A$ determined by a fixed automaton $\mathcal A_I$.
	So, by~\cref{constrsync:eqn:ideal-int} 
	constructing $\mathcal A \times \mathcal A_I$ gives a valid reduction to \constrsync{L}.	
\end{proof}

\begin{example}
	Let $\Sigma = \{a,b,c\}$ and consider the ideal language $I = \Sigma^* c^{n} \Sigma^*$
	that is regular.
	Let $L_1 = \{ u \in \Sigma^* : |u|_c \le n \}$,
	$L_2 = \{ u \in \Sigma^* : |u|_c = n \}$,
	and $L_3 = \{ u \in \Sigma^* : m \le |u|_c \le n \}$
	for $n,m \ge 1$ with $m \le n$.
	Then we have
	$L_1 \cap I = L_2 \cap I = L_3 \cap I = \{a,b\}^* c^n \{a,b\}^*$.
	From \crefitem{constr_sync:thm:introductory_examples}{constrsync:ex:item:abcL} 
	and \cref{constr_sync:gen_results:thm:ideal_separation} 
	we get that $\constrsync{L_i}$
	for $i\in \{1,2,3\}$
	are all \PSPACE-complete problems.
\end{example}

\begin{corollary}
	Let $L \subseteq \Sigma^*$ be regular
	and $u_1, \ldots, u_n \in \Sigma$.
	Then we have 
	\begin{align*}
	\constrsync{L\cap \Sigma^*u_1\cdots u_n\Sigma^*} & \leACZEROm \constrsync{L}, \\
	\constrsync{L\cap \Sigma^*u_1 \Sigma^* \cdots \Sigma^* u_n \Sigma^*} & \leACZEROm \constrsync{L}.
	\end{align*}
\end{corollary}
\begin{proof}
	The languages
	 $\Sigma^*u_1\cdots u_n\Sigma^*$
	 and $\Sigma^*u_1 \Sigma^* \cdots \Sigma^* u_n \Sigma^*$
	 are regular and ideals.
	 So, the claim follows from \cref{constr_sync:gen_results:thm:ideal_separation}.		
\end{proof}

\IN{principal ideal}
Mostly, we apply Theorem~\ref{constr_sync:gen_results:thm:ideal_separation}
with \emph{principal ideals} $\Sigma^* u \Sigma^*$
for $u \in \Sigma^*$.

As pointed out in the next remark, for a union the problem either has the same complexity as each part, or its complexity is bounded from below by the intersection of one part with a principal ideal.

\begin{remark}\label{constrsync:compl:ideal_intersection_for_unions}
	Consider two regular languages $L_1, L_2 \subseteq \Sigma^*$.
	There exists $u \in L_2$
	such that $u \notin \fact{L_1}$
	or $L_2 \subseteq \fact{L_1}$.
	By \cref{constrsync:genresults:cor:subset},
	in the former case, 
	we have
	\[
	\constrsync{L_1} \eqACZEROm \constrsync{L_1\cup L_2}
	\]
	and in the latter case, since then
	\[
	( L_1 \cup L_2 )  \cap \Sigma^* u \Sigma^* = L_2 \cap \Sigma^* u \Sigma^*,
	\]
	from \cref{constr_sync:gen_results:thm:ideal_separation}
	we get 
	$\constrsync{L_2 \cap \Sigma^*u \Sigma^*} \leACZEROm \constrsync{L_1\cup L_2}$.
	%	This scheme is often use in hardness proofs.
	%	However, finding an appropriate word $u \in \fact{L_2}$ 
	%	such that a provable 
	%	hard problem $\constrsync{L_2 \cap \Sigma^*u \Sigma^*}$
	%	results is the tricky part.
\end{remark}

Because we often use \cref{constr_sync:gen_results:thm:ideal_separation} to ``separate
a hard part out'' of a constraint language we called it the Ideal Separation Theorem
in~\cite{HoffmannCocoonExtended}, but we changed the name here because it involves an intersection
and we thought the new name to be more appropriate.

Lastly, we show how the Ideal Intersection Theorem  can be of help
in a heuristic way by employing it to exclude problematic
cases from a construction. This usage might not always be possible (in fact, in most cases it is not 
possible),
but we give the following example to also build some intuition
on various ways in which the Ideal Intersection Theorem 
can be applied (more examples of its application 
will follow in the rest of this work).

\begin{remark} \label{constrsync:rem:construction-ideal-intersection}
	Consider $L = \{bb,ba\}^*a\Gamma^*$ for $|\Gamma|\le 1$.
	We use a reduction from \SIS
	to show \PSPACE-hardness of $\constrsync{L}$.
	However, the construction we give
	does not work straightforward way for~$L$,
	but we then argue that it can be
	used for a modified problem obtained by intersection with an appropriate ideal.
	We then obtain \PSPACE-hardness of $\constrsync{L}$
	from \cref{constr_sync:gen_results:thm:ideal_separation}.
	
	Let $\semiautomat{A}{Q}{\{a,b\}}{\delta}$ and $S \subseteq Q$.
	Set $\semiautomat{A'}{Q \times \{\varepsilon,b\} \cup \{t\}}{\{a,b\}}{\delta'}$
	where~$\delta'(t, x) = t$ and
	\[
	\delta'((q, u), x) =
	\begin{cases}
		t                           & \text{if } u = \varepsilon \text{ and } x = a, \\
		(q, b) & \text{if } u = \varepsilon \text{ and } x = b, \\
		(\delta(q, a), \varepsilon) & \text{if } ux = bb, \\
		(\delta(q, b), \varepsilon) & \text{if } ux = ba, \\
		(q, \varepsilon)            & \text{if } ux \notin \{bb,ba\} \text{ and } q \notin S, \\
		t							& \text{if } ux \notin \{bb,ba\} \text{ and } q \in S, 		
	\end{cases}
	\]
	for $u \in \{\varepsilon, b\}$, $q \in Q$ and $x \in \{a,b\}$.
	
	First, suppose $u_1 \cdots u_n a v$
	with $u_1, \ldots, u_n \in \{bb,ba\}$ and $v \in \Gamma^*$ 
	synchronizes~$\mathcal A'$.
	Then it must map all states to~$t$.
	Consider the state $\delta'((q, \varepsilon), u_1\cdots u_n)$
	for some $q \in Q$.
	Then it must have the form $(p, \varepsilon)$
	for some $p \in Q$.
	Assume $p \notin S$.
	Then $\delta'((p, \varepsilon), a) = (p, \varepsilon)$
	and since $|\Gamma|\le 1$
	the word~$v$ cannot map $(p, \varepsilon)$ to~$t$.
	Thus we must have $p \in S$
	and then the word $v_1 \cdots v_n$
	with $v_i = a$ if $u_i = bb$
	and $v_i = b$ if $u_i = ba$
	maps~$Q$ into~$S$.	
	
	If $u_1 \cdots u_n$ for $u_1,\ldots,u_n \in \{a,b\}$
	maps~$Q$ into~$S$,
	then $v_1 \dots v_n a$ with $v_i = bb$ if $u_i = a$
	and $v_i = ba$ if $u_i = b$ for all $i \in \{1,\ldots,n\}$
	maps all states of~$Q \times \{\varepsilon\}$ to~$t$.
	However, if $u_1 = \cdots = u_n = bb$,
	then $\delta'(Q \times \{b\}, u_1\cdots u_n) \subseteq Q \times \{b\}$
	and thus $\delta'(Q \times \{b\}, u_1 \cdots u_n a) \cap \delta(Q, a) \times \{\varepsilon\} \ne \emptyset$.
	Now no word from~$\Gamma^*$ can map this state to~$t$.
	But this problem does not arise if we read~$ba$
	at least once, since this word maps all states in $Q \times \{b\}$
	to~$t$.
	Now, note that $L \cap \Sigma^* baa \Sigma^* = \{bb,ba\}baa\Gamma^*$.
	Thus, by \cref{constr_sync:gen_results:thm:ideal_separation} we obtain \PSPACE-hardness for $\constrsync{L}$
	if we can show \PSPACE-hardness
	for $\constrsync{\{bb,ba\}baa\Gamma^*}$
	and to establish \PSPACE-hardness of this problem
	we can use the above construction since the problematic case discussed above
	does not arise for $\{bb,ba\}baa\Gamma^*$.
\end{remark}

%\subsection{Homomorphic Mappings}
\subsection{Applying Morphisms and Solid Codes}
\label{constr_sync:subsec:morphisms}

%!TEX root = ../../../thesis-trier.tex

Here, in \cref{thm:hom_lower_bound_complexity},
we show that transforming a constraint language
yields a problem whose complexity is
at most that of the original problem w.r.t.~$\leACZEROm$-reductions.
In \cref{thm:constr_sync_hom_strongly_self_sync}
we show that for coding morphisms of solid codes
the complexity is preserved.

\begin{theorem}
	\label{thm:hom_lower_bound_complexity}
	Let $\varphi \colon \Sigma^* \to \Gamma^*$ be a fixed morphism
	and $L \subseteq \Sigma^*$.
	Then \[ \constrsync{\varphi(L)} \le_m^{\AC^0} \constrsync{L}. \]
\end{theorem}
\begin{proof}
	For the $\leACZEROm$-reduction,
	let $\semiautomat{A}{Q}{\Gamma}{\delta}$ be an input automaton. 
	We want to know if $\mathcal A$ is synchronizing with respect to~$\varphi(L)$.
	
	\medskip
	\noindent\underline{Construction:} Let $\semiautomat{A'}{Q}{\Sigma}{\delta'}$
	with $\applytrans{\delta'}{q}{x} = \applytrans{\delta}{q}{\applyfunc{x}{\varphi}}$ for $q \in Q$ and $x \in \Sigma$.
	We obtain a representation of $\mathcal A'$
	by transition matrices by multiplying
	for each $x \in \Sigma$
	the respective transition matrices
	of the letters in~$\applyfunc{x}{\varphi}$,
	which is doable by a \LOGSPACE-uniform circuit family in $\AC^0$ as $\varphi$ is fixed and matrix multiplication for a fixed number of matrices is in $\AC^0$ by \cref{prelim:circuit:ex:AC0mtx}.
		
	\medskip 
	\noindent\underline{Verification:} We have
	\begin{align*}
		                  \exists u \in \varphi(L) \colon |\delta(Q, u)| = 1 
		& \Leftrightarrow \exists v \in L \colon |\delta(Q, \varphi(v))| = 1 \tag{$\varphi(v) = u$} \\
		& \Leftrightarrow \exists v \in L \colon |\delta'(Q, v)| = 1  \tag{$\delta'(Q, v) = \delta(Q, \varphi(v))$}.  
	\end{align*}
So $\mathcal A'$ has a synchronizing word in $L$ if and only if
$\mathcal A$ has a synchronizing word in $\applyfunc{L}{\varphi}$.
\end{proof}

\begin{example}
	Let $L = \{a,b\}^* \cup \{a,c\}^* \cup \{b,c\}^*$.
	Then $\pi_{\{a,b\}}(L) = \{a,b\}^*$,
	so that $\constrsync{L}$ is \NL-hard
	with \cref{thm:hom_lower_bound_complexity}
	and \crefitem{constr_sync:thm:introductory_examples}{constrsync:ex:item:ab}.
	With \cref{constr_sync:thm:union} 
	we also get membership in \NL.
\end{example}	

For coding morphisms of solid codes,
we have a stronger statement than 
Theorem~\ref{thm:hom_lower_bound_complexity}.

\begin{theorem}
	\label{thm:constr_sync_hom_strongly_self_sync}
	Let $\varphi \colon \Sigma^* \to \Gamma^*$
	be a coding morphism such that\footnote{Recall that $\Sigma$ is finite, and so $C$ is finite. Also coding morphism
		are injective by definition.}
	$C = \varphi(\Sigma)$ 
	is a solid code.
	Then, for each regular $L \subseteq \Sigma^*$ we have 
	$
	\constrsync{L} \eqACZEROm \constrsync{\varphi(L)}.
	$
\end{theorem}
\begin{proof}
	By Theorem~\ref{thm:hom_lower_bound_complexity}, we have 
	$\constrsync{\varphi(L)} \leACZEROm \constrsync{L}$.
	Next, we show $\constrsync{L} \leACZEROm \constrsync{\applyfunc{L}{\varphi}}$
	by utilizing the Levenshtein-Romanov mappping from \cref{constr_sync:subsec:levenshtein-romanov-map}
    w.r.t.~the language~$C$.
	Now, let $\semiautomat{A}{Q}{\Sigma}{\delta}$ be an input automaton.
	Suppose $\Sigma = \{a_1, \ldots, a_k\}$
	and set $u_i = \applyfunc{a_i}{\varphi}$ for $i \in \{1,\ldots,k\}$.
	Then $C = \{u_1, \ldots, u_k\}$.	
	
	\begin{figure}[htbp]
		\centering		
		\begin{subfigure}{0.33\textwidth}
			\centering
	\begin{tikzpicture}[>=stealth,shorten >=1pt,auto]
	\node[state] (q0) at (90:2.5cm) {$q_0$};
	\node[state] (q1) at (210:2.5cm) {$q_1$};
	\node[state] (q2) at (330:2.5cm) {$q_2$};
	
	\path[->]
	(q0) edge[bend left=10] node[right] {$a_1$} (q2)
	(q2) edge[bend left=10] node[below] {$a_1$} (q1)
	(q1) edge[bend left=5]  node[left] {$a_1$}  (q0);
	
	\path[->] (q1) edge[bend left=10] node {$a_2$} (q2)
	(q2) edge[bend left=10] node {$a_2$} (q0)
	(q0) edge [loop above] node {$a_2$} (q0);
\end{tikzpicture}
			\caption{A base automaton~$\mathcal A$}
		\end{subfigure}
		\hfill
		\begin{subfigure}{0.55\textwidth}
		\centering
	\begin{tikzpicture}[>=stealth,shorten >=1pt, node distance=0.97cm, auto]
	
	%\tikzset{state/.style={}}
	\tikzset{state/.style={inner sep=2pt}}
	
	%%%%%%%%%%%%%%%%%%%%%%%%%%%%%%%%%%%%
	% TOP LAYER (epsilon)
	%%%%%%%%%%%%%%%%%%%%%%%%%%%%%%%%%%%%
	\node[state] (q0e) at (90:2.7cm) {$(q_0,\varepsilon)$};
	\node[state] (q1e) at (210:2.7cm) {$(q_1,\varepsilon)$};
	\node[state] (q2e) at (330:2.7cm) {$(q_2,\varepsilon)$};
	
	%%%%%%%%%%%%%%%%%%%%%%%%%%%%%%%%%%%%
	% MIDDLE LAYER (a)
	%%%%%%%%%%%%%%%%%%%%%%%%%%%%%%%%%%%%
	\node[state] (q0a) [below=of q0e] {$(q_0,a)$};
	\node[state] (q1a) [below=of q1e] {$(q_1,a)$};
	\node[state] (q2a) [below=of q2e] {$(q_2,a)$};
	
	%%%%%%%%%%%%%%%%%%%%%%%%%%%%%%%%%%%%
	% LOWER LAYER (aa)
	%%%%%%%%%%%%%%%%%%%%%%%%%%%%%%%%%%%%
	\node[state] (q0aa) [below=of q0a] {$(q_0,aa)$};
	\node[state] (q1aa) [below=of q1a] {$(q_1,aa)$};
	\node[state] (q2aa) [below=of q2a] {$(q_2,aa)$};
	
	\path[->]
	(q0e) edge (q0a)
	(q0a) edge (q0aa)
	(q0aa) edge[in=230,out=260, looseness=9] (q0aa)
	(q0aa) edge[dotted, in=280, out=310, looseness=9] (q0aa)
	
	(q1e) edge  (q1a)
	(q1a) edge  (q1aa)
	(q1aa) edge[loop below] ()
	
	(q2e) edge  (q2a)
	(q2a) edge  (q2aa)
	(q2aa) edge[loop below] ();		
	
	\path[->]
	(q0e)  edge[dashed, in=80, out=50, looseness=9] (q0e)
	(q0e)  edge[dotted, in=130, out=100, looseness=9] (q0e)
	(q0a)  edge[dotted, bend left=20] (q0e)
	(q0aa) edge[dashed, bend right=10] (q2e);
	
	\path[->]
	(q1e)  edge[dashed, in=80, out=50, looseness=9] (q1e)
	(q1e)  edge[dotted, in=130, out=100, looseness=9] (q1e)
	(q1a)  edge[dotted, bend left=20] (q1e)
	(q1aa) edge[dashed, bend left=65] (q0e);
	
	\path[->]
	(q2e)  edge[dashed, in=80, out=50, looseness=9] (q2e)
	(q2e)  edge[dotted, in=130, out=100, looseness=9] (q2e)
	(q2a)  edge[dotted, bend right=20] (q2e)
	(q2aa) edge[dashed, bend left=30] (q1e);		
	
	\path[->] 
	(q0a) edge[dashed, bend right=10] (q2e)
	(q2a) edge[dashed,bend left=30] (q1e)
	(q1a) edge[dashed,bend left=55] (q0e);	
	
	\path[->]
	(q1aa) edge[dotted, bend right=10] (q2e)
	(q2aa) edge[dotted, bend right=50] (q0e);		
\end{tikzpicture}
			\caption{Constructed automaton as in the proof.}
		\end{subfigure}		
		\caption[Construction from proof of \cref{thm:constr_sync_hom_strongly_self_sync}.]{Construction from the proof of \cref{thm:constr_sync_hom_strongly_self_sync}
		    for the solid code $\{ab, aac\}$
			and base automaton shown on the left side. Solid, dashed, and dotted lines represent transitions
			labeled with $a$, $b$, and~$c$, respectively.}
	\label{constrsync:fig:construction-solid-code}
	\end{figure}	
	
	\medskip
	\noindent\underline{Construction:} Define $\semiautomat{A'}{Q \times (\pref{C} \setminus C)}{\Gamma}{\delta'}$
	where, for $(q, u) \in Q \times \pref{C} \setminus C$
	and $x \in \Sigma$,
	\begin{equation}\label{constr_sync:eqn:aut_solid_code}
		\applytrans{\delta'}{(q,u)}{x} = \begin{cases}
			(q, \sigma_C(ux)) & \text{if } \sigma_C(ux) \notin C, \\
			(\applytrans{\delta}{q}{a_i}, \varepsilon) & \text{if } \exists i \in \{1,\ldots,k\} : \sigma_C(ux) = u_i. 
		\end{cases}
	\end{equation}
	By encoding the proper prefixes in the second component
	we keep track of the prefix scanned in the input so far.	
	See \cref{constrsync:fig:construction-solid-code} for an example construction
	for the solid code $C = \{ab,aac\}$.
	Note that in the example $\delta'((q, a), b) = \delta'((q,aa), b) = (\delta(q, a_1), \varepsilon)$
	since $\sigma_C(ab) = \sigma_C(aab) = ab$.
		
	
	\noindent\underline{Properties:} We give the proofs for Properties~\ref{constr_sync:solid_code_morphism:apply_codeword_proper_prefix} and~\ref{constr_sync:solid_code_morphism:apply_codeword},
	Property~\ref{constr_sync:solid_code_morphism:wordsvarphi} follows from Property~\ref{constr_sync:solid_code_morphism:apply_codeword} by induction
	as $\applyfunc{v}{\varphi} \in C^+$ for $v \in \Sigma^+$.
	\begin{enumerate}
		\item \label{constr_sync:solid_code_morphism:apply_codeword_proper_prefix}
		For $(q, u) \in Q \times ( \pref{C} \setminus C )$ and $v \in \pref{C} \setminus C$
		we have $\sigma_C(uv) \in \pref{C} \setminus C$ and $\applytrans{\delta'}{(q,u)}{v}  = (q, \sigma_C(uv))$.		
		
		\begin{proofitem}{\ref{constr_sync:solid_code_morphism:apply_codeword_proper_prefix}} 
		We show the claim by induction on $v$.
		As $\sigma_C(u) = u$, the claim is true for $v = \varepsilon$.
		Now, suppose $v = wx$ for $w \in \Sigma^*$ and $x \in \Sigma$.		
		As $C$ is overlap-free, 
		$wx \in \pref{C} \setminus C$ implies 
		$\sigma_C(uwx) \notin C$ and so $\sigma_C(uwx) \in \pref{C} \setminus C$;
		this implication is noted in \crefitem{constrsync:prop:sigma-solid-code-char}{constrsync:prop:sigma-solid-code-char:2}.
		By~\crefitem{constr_sync:prop:sigma_L_properties}{constr_sync:sigme_L:item:concatenation}
		we have $\sigma_C(\sigma_C(uw)x) = \sigma_C(uwx)$.
		So, for $q \in Q$, we obtain		
		\begin{align*}
			& \delta'((q,u), wx) \\
			& = \delta'(\delta'((q,u), w), x) & \tag{Def.~of $\delta'$} \\
			& = \delta'((q, \sigma_C(uw)), x) & \tag{Induction hypothesis} \\
			& = (q, \sigma_C(\sigma_C(uw)x)). & \tag{Eqn.~\eqref{constr_sync:eqn:aut_solid_code} as $\sigma_C(\sigma_C(uw)x) \notin C$} 
		\end{align*}
		Hence, we get $ \delta((q,u), wx) = (q, \sigma_C(uwx))$
        and the claim follows by induction.
     	\end{proofitem}
		
		\item \label{constr_sync:solid_code_morphism:apply_codeword}
		For $(q, u) \in Q \times ( \pref{C} \setminus C )$ and $i \in \{1,\ldots,k\}$ we have
		\[
		\applytrans{\delta'}{(q,u)}{u_i} = (\applytrans{\delta}{q}{a_i}, \varepsilon) 
		\]
		
		\begin{proofitem}{\ref{constr_sync:solid_code_morphism:apply_codeword}}
		Let $u_i = vx$ for $v \in \Sigma^*$ and $x \in \Sigma$.
		As $C$ is a prefix code, we have $v \in \pref{C} \setminus C$.		
		By \crefitem{constr_sync:prop:sigma_L_properties}{constr_sync:sigme_L:item:concatenation}
		we have $\sigma_C(\sigma_C(uv)x) = \sigma_C(uvx)$.
		By \crefitem{constrsync:prop:sigma-solid-code-char}{constrsync:prop:sigma-solid-code-char:3}
		as $vx \in C$, we find $\sigma_C(uvx) = vx = u_i$. Then
		\begin{align*}
			\delta'((q, u), vx) & = \delta'(\delta'((q,u), v), x)  & \tag{Def.~of $\delta'$} \\
			& = \delta'((q, \sigma_C(uv)), x) & \tag{Property~\ref{constr_sync:solid_code_morphism:apply_codeword_proper_prefix}} \\
			& = (\applytrans{\delta}{q}{a_i}, \varepsilon). & \tag{Eqn.~\eqref{constr_sync:eqn:aut_solid_code} as $\sigma_C(\sigma_C(uv)x)  = u_i$}
		\end{align*}
	    So, $\applytrans{\delta'}{(q,u)}{u_i} = (\applytrans{\delta}{q}{a_i}, \varepsilon)$.
		\end{proofitem}
		
		\item  \label{constr_sync:solid_code_morphism:wordsvarphi} For every $v \in \Sigma^+$ and $(q,u) \in Q \times (\pref{C} \setminus C)$,
		\begin{equation}\label{constr_sync:eqn:solid_code3}
			\applytrans{\delta'}{(q,u)}{\applyfunc{v}{\varphi}} = (\applytrans{\delta}{q}{v}, \varepsilon).
		\end{equation}
	\end{enumerate}	
	
	\medskip 
	\noindent\underline{Verification:} That the reduction is correct follows by the following
	equivalences:
	\begin{align*}
		& \exists w \in \varphi(L) \colon |\applytrans{\delta'}{Q \times (\pref{C}\setminus C)}{w}|  = 1 \\
		& \Leftrightarrow \exists v \in L \colon |\applytrans{\delta'}{Q \times (\pref{C}\setminus C)}{\varphi(v)}|  = 1 \\
		& \Leftrightarrow \exists v \in L \colon |\applytrans{\delta}{Q}{v} \times \{\varepsilon\}| = 1 \tag{Eqn.~\eqref{constr_sync:eqn:solid_code3}} \\
		& \Leftrightarrow \exists v \in L \colon |\applytrans{\delta}{Q}{v}| = 1. \qedhere
	\end{align*}
\end{proof}

\begin{remark} \label{constrsync:rem:alphabet_solid_special_case}
	The construction from \cref{constrsync:rem:alphabet} is a special case 
	of \cref{thm:constr_sync_hom_strongly_self_sync}
	as for $\Gamma, \Sigma$ with $\Gamma \subseteq \Sigma$
	the inclusion map $i \colon \Gamma \to \Sigma$
	given by $i(x) = x$ for $x \in \Gamma$
	is a coding morphism where $i(\Gamma) \subseteq \Sigma$ is a solid code.
\end{remark} 

\begin{remark}
	The construction from the proof of \cref{thm:constr_sync_hom_strongly_self_sync}
	with the transition function given by \cref{constr_sync:eqn:aut_solid_code}
	can be applied to any prefix code~$C$
	and yields a well-defined automaton~$\semiautomat{A'}{Q \times \pref{C}\setminus C}{\Sigma}{\delta'}$.	
	Can the assumption for~$C$ to be a solid code
	be weakened? 
	The crucial property for the reduction
	to work is that $\delta'(Q \times \pref{C} \setminus C, u_i) = \delta(Q, a_i) \times \{\varepsilon\}$
	for $i \in \{1,\ldots,k\}$ and $C = \{u_1, \ldots, u_k\}$.
	However, this is not guaranteed to hold if~$C$
	is merely a comma-free code.
	For example, consider $C = \{ u_1, u_2 \}$ with $u_1 = aab$ and $u_2 = bba$,
	which is a comma-free code, but not a solid code since $u_1$ and $u_2$ have a
	non-empty proper overlap.
	Then we have
	$
	\delta'(Q \times \{aa\}, u_2) = \delta(Q, a_1) \times \{a\}  	
	$
	and so $\delta(Q, a_1) \times \{a\} \subseteq \delta'(Q \times \pref{C} \setminus C, u_2)$.
	However, there are codes that are not solid codes for which it does hold.
	For example, for $C = \{ a_1, \ldots, a_{k-1}, a_k^m a_{k-1}\}$ for distinct $a_1, \ldots, a_k \in \Sigma$
	and $m \ge 1$.
\end{remark}

\begin{remark}
	Solid codes and the construction from the proof of  \cref{thm:constr_sync_hom_strongly_self_sync}
	can also be used to give an alternative proof of \cref{prelim:prop:sync_NL_complete}
	by first giving an automaton derived from the digraph
    by labeling the edges with letters and using an additional letter
    to realize the transitions going back to the start state.
    Then using a binary solid code
    and the construction from the proof of \cref{thm:constr_sync_hom_strongly_self_sync}
    we can reduce the alphabet to a binary alphabet.
    However, note that the words $bb,a,ba$ actually used in the proof of \cref{prelim:prop:sync_NL_complete}
    do not form a solid code.
\end{remark}

\begin{remark}
	We actually have $\Syn(\mathcal A') \cap \varphi(\Sigma^*) = \varphi(\Syn(\mathcal A))$
	in the construction from the proof of \cref{thm:constr_sync_hom_strongly_self_sync}.
\end{remark}

Coding morphism of solid codes are not the only morphisms
that preserve the complexity, as shown in the following remark.

\begin{remark}
	Consider the coding morphism $\varphi \colon \{a,b\}^* \to \{c,d, e\}^*$
	given by $\varphi(a) = cd$ and $\varphi(b) = ec$.
	Now $\varphi(\{a,b\}) = \{cd, ec\}$ is not a solid code.
	For $\semiautomat{A}{Q}{\{a,b\}}{\delta}$
	let 
	$\semiautomat{A'}{Q'}{\{c,d\}}{\delta'}$
	with $Q' = Q \times \{\varepsilon, c, e\}$
	where, for $q \in Q$, $u \in \{\varepsilon,c,d\}$, and $x \in \Sigma$,
	\[
	\delta'((q, u), x)
	= 
	\begin{cases}
		(\delta(q, a), \varepsilon)  & \text{if } u \in \{\varepsilon,c\} \text{ and } x = d, \\
		(\delta(q, b), \varepsilon)  & \text{if } u = e \text{ and } x = c, \\
		(q,x)                        & \text{if } u \in \{\varepsilon,c,e\} \text{ and } x \in \{c,e\} \text{ and } ux \ne ec \\
		%(q,e)                       & \text{if } u \in \{\varepsilon,e\} \text{ and } x = e,
		(q,u)                        & \text{if } u = e \text{ and } x = d.
	\end{cases}
	\]
	Then we have $\delta'((q, u), \varphi(v)) = (\delta(q, v), \varepsilon)$
	for $v \in \{a,b\}^+$ with $u \in \{\varepsilon, c, e\}$
	and $q \in Q$.
	This yields 
	\begin{align*}
		\delta'(Q', \varphi(a)) = \delta(Q, a) \times \{\varepsilon\}, \\
		\delta'(Q', \varphi(b)) = \delta(Q, b) \times \{\varepsilon\},	
	\end{align*}
	and	that for $|Q| \ge 2$ and regular $L \subseteq \{a,b\}^*$,
	the automaton~$\mathcal A$ has a synchronizing word in~$L$
	if and only if~$\mathcal A'$ has a synchronizing word in~$\varphi(L)$.
\end{remark}

\subsection{Concatenation and Complexity Lower Bounds}
\label{constr_sync:subsec:concat}

%!TEX root = ./../../../springer-mono-format/book.tex
%!TEX root = ./../../../thesis-trier.tex

In general, the computational complexity might increase or decrease
when concatenating languages.
For example, 
by \cref{constr_sync:thm:introductory_examples},
\cref{gen_results:cor:factors_etc_complexity}
and as $\fact{(a+b+c)^*c(a+b)^*} = \fact{(a+b)^*}$,
$\constrsync{(a+b)^*}$, $\constrsync{c(a+b)^*}$,
and $\constrsync{(a+b+c)^*c(a+b)^*}$ are 
complete for \NL, \PSPACE, and \NL, respectively.
Here, we identify conditions on concatenating languages
to obtain complexity lower bounds.
First, we give two results
for concatenating a language
to the front of a constraint language such that hardness is preserved, where
\cref{constrsync:prop:UanL} is related to languages contained in a subalphabet
and \cref{constr_sync:thm:prefix_before} to prefix codes.
Then, \cref{constr_sync:prop:Gamma_suffix}
is concerned with concatenation to the back.
The conditions presented here are used to deduce hardness results in later sections.

The reductions of the present section involve operations
as mentioned in \cref{prelim:rem:reductionsAC0}
and can be realized by $\AC^0$ circuits.

%\subsubsection{Concatenating to the Front}

\begin{proposition}
	\label{constrsync:prop:UanL}
	Let $L \subseteq \Sigma^*$ 
	and $U \subseteq \Gamma^*$ for some $\Gamma \subseteq \Sigma$
	be regular languages and $w \in \Sigma^*$.
	Suppose one of the following holds:
	\begin{enumerate}
	\item $L \cap \Gamma\Sigma^* = \emptyset$ with $w \in \Gamma^*$, or
	
	\item $w \notin \Gamma^*$.
	\end{enumerate}
	Then $\constrsync{wL} \leACZEROm \constrsync{UwL}$.
\end{proposition}         
\begin{proof} 
	First, assume $L \cap \Gamma\Sigma^* = \emptyset$ and $w \in \Gamma^*$.	
	Let $\semiautomat{A}{Q}{\Sigma}{\delta}$. We construct a semiautomaton $\mathcal A'$ that has a synchronizing
	word in $UwL$ if and only if $\mathcal A$ has a synchronizing word in $wL$.
	
	\medskip
	\noindent\underline{Construction:} Let $\semiautomat{A'}{Q \times \{0,1\}}{\Sigma}{\delta'}$
	where, for $i \in \{0,1\}$, $q \in Q$ and $x \in \Sigma$,
	\[
	\applytrans{\delta'}{(q, i)}{x} = 
	\begin{cases}
		(q, 0) & \text{if } i = 0 \text{ and } x \in \Gamma, \\
		(\applytrans{\delta}{q}{wx}, 1) & \text{if } i = 0 \text{ and } x \notin \Gamma, \\
		(\applytrans{\delta}{q}{x}, 1) & \text{if } i = 1.
	\end{cases}
	\]
	
	\noindent\underline{Properties:} Let $q \in Q$, $u \in U$ and $v \in L$.
	Then, as $v \notin \Gamma\Sigma^*$ and $uw \in \Gamma^*$,
	we have 
	\begin{align}
		\applytrans{\delta'}{(q, 0)}{uwv} & = \applytrans{\delta'}{(q, 0)}{v}, \label{constrsync:eqn:UL1} \\
		\applytrans{\delta'}{(q, 0)}{v}     & = (\applytrans{\delta}{q}{wv}, 1), \label{constrsync:eqn:UL3} \\
		\applytrans{\delta'}{(q, 1)}{uwv} & = (\applytrans{\delta}{q}{uwv}, 1). \label{constrsync:eqn:UL2}
	\end{align} 
	
	\noindent\underline{Verification:} Let $u \in U$ and $v \in L$.
	Then we have
	\begin{align*}
		\applytrans{\delta'}{Q \times \{0,1\}}{uwv} 
		& = \applytrans{\delta'}{Q \times \{0\}}{uwv} \cup \applytrans{\delta'}{Q \times \{1\}}{uwv} \\
		& = \applytrans{\delta'}{Q \times \{0\}}{v} \cup \applytrans{\delta'}{Q \times \{1\}}{uwv} & \tag{Eqn.~\eqref{constrsync:eqn:UL1}} \\
		& = \delta(Q, wv) \times \{1\} \cup \delta'(Q \times \{1\}, uwv) & \tag{Eqn.~\eqref{constrsync:eqn:UL3}} \\
		& = \delta(Q, wv) \times \{1\} \cup \delta(Q, uwv) \times \{1\} & \tag{Eqn.~\eqref{constrsync:eqn:UL2}} \\
		& = \delta(Q, wv)  \times \{1\}.                                   & \tag{$\delta(\delta(Q, u), wv) \subseteq  \delta(Q, wv)$}
	\end{align*}
	So $|\applytrans{\delta'}{Q \times \{0,1\}}{uwv}| = 1$
	is equivalent to $|\delta(Q,wv)| = 1$
	and $\mathcal A'$ has a synchronizing word in $UwL$
	if and only if $\mathcal A$ has a synchronizing word in $wL$.
	
    \medskip 
    
    Lastly, for the remaining case, assume $w \notin \Gamma^*$.
    Then $w = ubv$  for some $u \in \Gamma^*$, $b \notin \Gamma$, and $v \in \Sigma^*$.
    Set $L' = bvL$. We have $L' \cap \Gamma\Sigma^* = \emptyset$. 
    So, applying the previous case to $L', u$ and $U$,
    we get 
    $\constrsync{uL'} \leACZEROm \constrsync{UuL'}$
    where $uL' = wL$.
\end{proof}             

\begin{example}\label{constrsync:ex:PSPACEhard}
	By~\crefitem{constr_sync:thm:introductory_examples}{constrsync:ex:item:abcL},
	$\constrsync{c(a+b)^*}$	is \PSPACE-complete.
	By \cref{constrsync:prop:UanL}, the problems
	$\constrsync{aabbc(a+b)^*}$ and $\constrsync{a^*c(a+b)^*}$
	are \PSPACE-complete as well.
\end{example}

For reference, we list two special cases  of~\cref{constrsync:prop:UanL}.

\begin{corollary}
	\label{constrsync:cor:UanL}
	Let $L \subseteq \Sigma^*$ be regular.
	
	\begin{enumerate}
		\item \label{constrsync:cor:UanL:1} If $L \cap \Gamma\Sigma^* = \emptyset$,
		$\Gamma \subseteq \Sigma$ and $U \subseteq \Gamma^*$ is regular, then $\constrsync{L} \leACZEROm \constrsync{UL}$. 
		
		\item \label{constrsync:cor:UanL:2} $\constrsync{wL} \leACZEROm \constrsync{a^*wL}$	for $w \notin a^*$ and $a \in \Sigma$.
	\end{enumerate}
\end{corollary}                    

Next, we present a condition related to prefix codes.

\begin{theorem}
	\label{constr_sync:thm:prefix_before}
	\label{constr_sync:thm:in_font}
	%Let $U, V \subseteq \Sigma^*$ be regular and $P = U \setminus U\Sigma^+$ (those words in $U$ 
	%without a proper prefix in $U$).
	%Then $\constrsync{V} \le_m^{\log} \constrsync{PV}$. 	
	Let $L \subseteq \Sigma^*$ be regular 
	and $P \subseteq \Sigma^*$ be a prefix code that is a regular language. 
	Then 
	$
	\constrsync{L} \leACZEROm \constrsync{PL}.
	$	
\end{theorem}
\begin{proof}    	
	By~\crefitem{constr_sync:prop:aut_prefix_free}{constr_sync:prop:aut_prefix_free:item:trim} 
	there exists a trim PDFA
	$\automat{C}{R}{\Sigma}{\rho}{r_0}{\{e\}}$
	such that $P = L(\mathcal C)$
	and where the state~$e$ has no outgoing transitions.
	Suppose $\semiautomat{A}{Q}{\Sigma}{\delta}$ is an input semiautomaton. 
	
	
	% bild oder beispiel?
	\medskip 	
	\noindent\underline{Construction:} Let $\semiautomat{A'}{R \times Q}{\Sigma}{\delta'}$
	where, for $(r,q) \in R \times Q$ and $x \in \Sigma$,
	\[
	\applytrans{\delta'}{(r, q)}{x} = \begin{cases}
		(\applytrans{\rho}{r}{x}, q)   & \text{if } \applytrans{\rho}{p}{x} \text{ is defined,} \\
		(e, \applytrans{\delta}{q}{x}) & \text{if } \rho(p, x) \text{ is undefined. }
	\end{cases} 
	\]
	
	
	
	 %Property~\ref{constrsync:prefix:prop:e} readily follows from the definition of $\delta'$.
	%For the other properties, we give the proofs.  
	\medskip 	
	\noindent\underline{Properties:} Property~\ref{constrsync:prefix:defdeltastar} follows readily by induction.
	We give the proof for Property~\ref{constrsync:prefix:prop:exactsync}.
	\begin{enumerate}    	
		
		\item \label{constrsync:prefix:defdeltastar} For $r \in R$, $q \in Q$ and $u \in \Sigma^*$ we have
		\begin{equation}\label{constrsync:prefix:eqn:transitionword}
			\delta'((r, q), u)
		=
		\begin{cases}
		(e, \delta(q, xu_2)) & \text{if } \exists u_1, u_2 \in \Sigma^* \ \exists x \in \Sigma \colon u = u_1 x u_2 \land \\ 
		                    & \qquad \rho(r, u_1) \text{ is defined } \land \\ 
		                    & \qquad \rho(r, u_1 x) \text{ is undefined}, \\
		(\rho(r, u), q)     & \text{if } \rho(r, u) \text{ is defined. } 
		\end{cases}
		\end{equation}
				
\begin{comment} 				
		\item \label{constrsync:prefix:defdeltastar}
		For $r \in R$, $q \in Q$ and $u \in \Sigma^*$ we have	
		\[
		\delta'((r, q), u)
		=
		\begin{cases}
		(e, \delta(q, u)) & \text{if } r = e, \\
		(\rho(r, u), q) & \text{if } r \ne e \text{ and } \rho(r, u) \text{ is defined. }
		\end{cases}
		\]
		
		\begin{proofitem}{\ref{constrsync:prefix:defdeltastar}}
		For $r = e$ the claim follows readily from the definition of $\delta'$
		and the fact that $\rho(r, x)$ is undefined for each $x \in \Sigma$.
		Next, suppose $r \ne e$ and $\rho(r,u)$ is defined. 
		We do induction on $u$. For $u = \varepsilon$
		the claim is clear. So, suppose we have $u = vx$
		with $v \in \Sigma^*$ and $x \in \Sigma$.
		Then $\rho(r, v)$ is defined and, by the definition of $\delta'$,
		we obtain $\delta'((\rho(r, v), q), x) = (\rho(r, vx), q)$
		and the claim follows by induction.		
		\end{proofitem} 
\end{comment}
\begin{comment} 
		\item \label{constrsync:prefix:prop:e} Let $u \in \Sigma^*$ and $q \in Q$. Then $\applytrans{\delta'}{(e, q)}{u} = (e, \applytrans{\delta}{q}{u})$.
		
		\item  \label{constrsync:prefix:prop:prefixP}  Let $u \in \pref{P} \setminus P$ and $q \in Q$.
		Then $\rho(r_0, u)$ is defined and $\applytrans{\delta'}{(r_0, q)}{u} = (\rho(r_0, u), q)$
		with $\rho(r_0, u) \ne e$.
		
		\begin{proofitem}{\ref{constrsync:prefix:prop:prefixP}}
		We do induction on $u$. For $u = \varepsilon$ 
		the claim is clear. So, suppose we have $u = u'x$ for some $u' \in \Sigma^*$ and $x \in \Sigma$.
		Then, $u' \in \pref{P}\setminus P$ and, by induction hypothesis, 
		$\rho(r_0, u')$
		is defined, $\delta'((r_0, q), u') = (\rho(r_0, u'), q)$ and $\rho(r_0, u') \ne e$.    	 
		So, as $P$ is a prefix code and $u$ a proper prefix,
		$\rho(r_0, u'x)$ is defined with $\rho(r_0, u'x) \ne e$
		and, by the definition of $\delta'$,
		we find 
		\[ 
		\delta'((r_0, q), u) = \delta'((\rho(r_0, u'), q), x) = (\rho(r_0, u'x), q),
		\]
		which gives the claim by induction on the structure of $u$.
		\end{proofitem}
		
		\item \label{constrsync:prefix:prop:wordP} Let $u \in P$ and $q \in Q$. Then $\applytrans{\delta'}{(r_0, q)}{u} = (e, q)$.
		
		\begin{proofitem}{\ref{constrsync:prefix:prop:wordP}}
		Let $u = u'x$ for $u' \in \pref{P}$ and $x \in \Sigma$.
		By Property~\ref{constrsync:prefix:prop:prefixP}
		we get $\applytrans{\delta'}{(r_0, q)}{u'} = (\rho(r_0,u'), q)$
		with $\rho(r_0, u') \ne e$.
		Then, by the definition of $\delta'$, this yields
		$\applytrans{\delta'}{(r_0, q)}{u} = \applytrans{\delta'}{(\applytrans{\rho}{r_0}{u'}, q)}{x}
		= (\applytrans{\rho}{\applytrans{\rho}{r_0}{u'}}{x}, q) = (e, q)$.
		\end{proofitem}
\end{comment}
\begin{comment} 
		\item \label{constrsync:prefix:prop:undefined} Let $w \in \Sigma^*$ and $r \in R$. If $\rho(r, w)$ is undefined,
		then $\delta'((r, q), w) \in \{e\} \times Q$ for each $q \in Q$.
		
		\begin{proofitem}{\ref{constrsync:prefix:prop:undefined}}
		Let $w = w_1 x w_2$ with $w_1, w_2 \in \Sigma^*$, $x \in \Sigma$ and 
		$w_1$ of maximal length
		such that $\rho(r, w_1)$ is defined.
		By Property~\ref{constrsync:prefix:defdeltastar} 
		we have $\delta'((r, q), w_1) = (\rho(r, w_1), q)$.
		Then, by the definition of $\delta'$, as $\rho(\rho(r, w_1), x)$ is undefined,
		we find 
		\[
		\delta'((\rho(r, w_1), q), x) = (e, \delta(q, x)).
		\]
		By Property~\ref{constrsync:prefix:defdeltastar} 
		we obtain $\delta'((e, \delta(q, x)), w_2) = (e, \delta(q, xw_2)) \in \{e\} \times Q$.
		\end{proofitem}
\end{comment}
		
		\item \label{constrsync:prefix:prop:exactsync} Let $w \in \Sigma^*$ be such that $\applytrans{\rho}{R}{w} = \{e\}$.
		Then $\applytrans{\delta'}{R \times Q}{w} = \{e\} \times Q$.
		
		
		\Cref{constrsync:prefix:eqn:transitionword} 
		gives $\delta'(R \times Q, w) \subseteq \{e\} \times Q$.
		Furthermore, there exists $r \in R$ such that $\rho(r, w) = e$.
		Then, for each $q \in Q$, we have $\delta'((r, q), w) = (e,q)$
		and this gives $\{e\} \times Q = \delta'(\{r\} \times Q, w) \subseteq \delta'(R \times Q, w)$. 	    	
	\end{enumerate}   
	
	\medskip 
	\noindent\underline{Verification:} 
	By Property~\ref{constrsync:prefix:defdeltastar}
	we have $\delta'(\{r_0\} \times Q, uv) = \{e\} \times \delta(Q, v)$
	for $u \in P$ and $v \in \Sigma^*$. (note that, since~$P$ is a 
	prefix code, $uv$ can only be factored into a word from~$P$
	and an arbitrary word in one way)
	So, $|\delta'(R \times Q, uv)| = 1$
	for $u \in P$ and $v \in L$
	implies $|\delta(Q, v)| = 1$.	
	
	Conversely, suppose $|\delta(Q, v)| = 1$
	for $v \in L$. 
	By Property~\ref{constrsync:prefix:defdeltastar} and as $\rho(e,x)$ is undefined for each $x \in \Sigma$
	we have $\delta'(\{e\} \times Q, v) = \{e\} \times \delta(Q, v)$ (here, we used that~$P$ is a prefix code).
	By \crefitem{constr_sync:prop:aut_prefix_free_exact_sync}{constr_sync:prop:aut_prefix_free:item:sync} 
	there exists $w \in P$ with
	$\rho(R, w) = \{e\}$.
	Then, by Property~\ref{constrsync:prefix:prop:exactsync} 
	we have $\applytrans{\delta'}{R \times Q}{w} \subseteq \{e\} \times Q$ (note that
	we might have that $\rho(r, w)$ is undefined for some $r \in R$ 
	and that, in fact, we have equality since $\delta'(\{r_0\}\times Q), w) = \{e\} \times Q$).
	So $|\delta'(R \times Q, wv)| = 1$.
\end{proof} 

\begin{remark}
	For $L = \{c\}$ and $P = \{a,b\}^*c$
	we have that $\constrsync{L}$ 
	and $\constrsync{PL}$
	are in $\AC^0$
	and \PSPACE-complete, respectively,
	by \cref{constr_sync:thm:introductory_examples}.
	However, the complexity does not have to increase when concatenating
	a prefix code that is regular to the front of a constraint
	language as can be seen with $L = \Sigma^*$.
	Then $\constrsync{L} \eqACZEROm \constrsync{PL}$
	by \cref{constrsync:genresults:cor:left-right-ideal}.		
\end{remark}

As $\{u\}$ for $u \in \Sigma^*$ is a prefix code, we obtain the following corollary.

\begin{corollary}
	Let $L \subseteq \Sigma^*$ be regular and $u \in \Sigma^*$.
	Then $\constrsync{L} \leACZEROm \constrsync{uL}$.
\end{corollary}

%\subsubsection{Concatenating to the Back} 

Next, we state a result on concatenating
languages to the back of the constraint language.

\begin{proposition}
	\label{constr_sync:prop:Gamma_suffix}
	Let $L \subseteq \Sigma^*$ be regular and $\Gamma \subseteq \Sigma$
	such that $L \cap \Sigma^* \Gamma = \emptyset$.
	If $U \subseteq \Gamma^*$ is a regular non-empty language,
	then $\constrsync{L} \leACZEROm \constrsync{LU}$.   
\end{proposition}
\begin{proof}
	We can assume $L \ne \emptyset$. 
	Let $\semiautomat{A}{Q}{\Sigma}{\delta}$. We construct a semiautomaton $\mathcal A'$ that has a synchronizing
	word in $LU$ if and only if $\mathcal A$ has a synchronizing word in $L$.
	\FIG~\ref{fig:irred-reduction} shows the basic idea of the construction graphically.
	\begin{comment} 
		For each $q \in Q$ and $a \in \Gamma$, set $S_{q,a} = \{ \applytrans{\delta}{q}{a^i} \mid i \ge 0 \}$.
		Define  
		$\semiautomat{A'}{Q'}{\Sigma}{\delta'}$ with
		\[
		Q' = \bigcup_{q \in Q} \bigcup_{a \in \Gamma} \{q\} \times S_{q,a}.
		\]
		and, for $(q, p) \in \{q\} \times S_{q,a}$, $x \in \Sigma$
		and $p = \applytrans{\delta}{q}{a^i}$, \todo{repräsentantenunabhgängig zeigen.}
		\[
		\applytrans{\delta'}{(q, p)}{x} = \begin{cases}
			(q, \applytrans{\delta}{p}{a})     & \text{if } x = a \\
			(\applytrans{\delta}{q}{a^ix}, \applytrans{\delta}{q}{a^ix})  & \text{if } x \ne a.
		\end{cases}
		\]
		
		Let $uv$ with $u \in L$ and $v \in U$ be a synchronizing word for $\mathcal A'$.
		Note that $u$ does not end with a letter in $\Gamma$.
		So, if $|u| > 0$,\todo{genauer?}
		we have $\applytrans{\delta'}{Q'}{u}\subseteq \{ (q,q) \mid q \in Q \}$.
		For $(q,q)$ we have $\applytrans{\delta'}{(q,q)}{a^i} = (q, \applytrans{\delta}{q}{a^i})$
		for all $i \ge 0$.
		Hence, $|\applytrans{\delta'}{Q'}{u}| = |\applytrans{\delta'}{(q,q)}{v}|$
	\end{comment} 
	
	\begin{figure}[t]
		\centering
		{			
			\begin{tikzpicture}[rounded corners=10pt]						
				% Main automaton
				\node[yshift=1cm] (A) at (0,0)
				{$\mathcal A'_{|\{(q,q)\mid q \in Q\},\,\Sigma\setminus\Gamma}
					\cong
					\mathcal A_{|\Sigma\setminus\Gamma}$};
				
				\node[draw, dashed, rectangle, fit={(-5,1.4) (5,-1.4)}] {};
				
				% Left column
				\node (q) at (-3.5,-1) {$(p,p)$};				
				\node[draw,circle,inner sep=2pt] (dummyq) at (-3.5,-2.5) {};
				\node at (-3.5,-2.75) {$\vdots$};				
				\node (labSq) at (-3.5,-3.7) {$\{p\}\times Q$};
				\node[draw,dashed,rectangle,inner sep=14pt, fit={(q) (labSq)}] (Sq) {};
				
				% Right column
				\node (p) at (3.5,-1) {$(q,q)$};				
				\node[draw,circle,inner sep=2pt] (dummyp) at (3.5,-2.5) {};
				\node at (3.5,-2.75) {$\vdots$};				
				\node (labSp) at (3.5,-3.7)	{$\{q\}\times Q$};
				\node[draw,dashed,rectangle,inner sep=14pt,	fit={(p) (labSp)}] (Sp) {};
				
				% Target states
				\node[draw,circle,inner sep=2pt] (r)  at (-1.3,-0.4) {};				
				\node[draw,circle,inner sep=2pt] (r2) at (1.3, -0.4) {};
				\node[above=0.15cm of r,  xshift=-0.4cm] {$(\delta(p,ua),\delta(p,ua))$};
				\node[above=0.15cm of r2, xshift=+0.4cm] {$(\delta(q,ua),\delta(q,ua))$};
				
				% Arrows
				\path[->] (q) edge[left,pos=.7] node {$u\in\Gamma^*$} (dummyq);				
				\path[->] (p) edge[right,pos=.7] node {$u\in\Gamma^*$} (dummyp);
				
				\path[->] (dummyq) edge[bend right,below,pos=.4] node [fill=white, anchor=center] 
				{$a \in \Sigma \setminus \Gamma$} (r);
				\path[->] (dummyp) edge[bend left,below,pos=.4]  node [fill=white, anchor=center] 
				{$a \in \Sigma \setminus \Gamma$} (r2);			
			\end{tikzpicture}
					
		}
		\caption[Schematic of reduction from the proof of Proposition~\ref{constr_sync:prop:Gamma_suffix}]{Schematic of the reduction from the proof of Proposition~\ref{constr_sync:prop:Gamma_suffix}.
			For a state $q$ and word $u \in \Gamma^*$ we ``branch out'' into a copy $\{q\} \times Q$.
			For a letter from $a \in \Sigma \setminus \Gamma$,
			we return to the state $(\applytrans{\delta}{q}{ua}, \applytrans{\delta}{q}{ua})$.
%			corresponding to the state in $\mathcal A$ when reading~$ua$.
			This is indicated 
			for two states $p$ and $q$.
		}
		\label{fig:irred-reduction}
	\end{figure}
	
	
	\medskip
	\noindent\underline{Construction:} Let 
	$\semiautomat{A'}{Q'}{\Sigma}{\delta'}$
	where\footnote{Actually, setting
		\[
		Q' = \bigcup_{q \in Q} \{q\} \times S_q
		\]
		where $S_{q} = \{ \applytrans{\delta}{q}{u} \mid u \in \Gamma^* \}$
		for $q \in Q$
		would be sufficient. However, it is not clear whether the
		sets $S_q$ can be computed by an \LOGSPACE-uniform $\AC^0$ circuit family
		whereas $Q \times Q$
		can be computed by such a circuit family.} 
	$Q' = Q\times Q$
	and, for $(q, p) \in \{q\} \times Q$, $x \in \Sigma$,
	\[
	\applytrans{\delta'}{(q, p)}{x} = \begin{cases}
		(q, \applytrans{\delta}{p}{x})     & \text{if } x \in \Gamma, \\
		(\applytrans{\delta}{p}{x}, \applytrans{\delta}{p}{x})  & \text{if } x \notin \Gamma.
	\end{cases}
	\]
		
	\noindent\underline{Properties:} Next, we collect some properties of the given construction.	
	\begin{enumerate} 
		
		\item \label{constr_sync:item:word_in_Gamma} For $w \in \Gamma^*$ and $q,p \in Q$, we have
		$
		%\begin{equation}\label{constr_sync:eqn:word_in_Gamma}
		\applytrans{\delta'}{(q,p)}{w}
		= (q, \applytrans{\delta}{p}{w}).
		%\end{equation} 
		$
		
		\begin{proofitem}{\ref{constr_sync:item:word_in_Gamma}} 
		This is clear for $w = \varepsilon$.
		Next, let $w = ux$ for some $x \in \Sigma$ and $u \in \Sigma^*$.
		Inductively, we have $\applytrans{\delta'}{(q,p)}{ux}
		= \applytrans{\delta'}{(q,\delta(p, u))}{x}$.
		Since $x \in \Gamma$,
		by the definition of the transition function
		we find $\applytrans{\delta'}{(q,\delta(p, u))}{x} = (q, \delta(p,w))$.			
		\end{proofitem}		
		
		\item \label{constr_sync:item:Gamma_suffix} If $w \in \Sigma^* (\Sigma \setminus \Gamma)$ and $q,p \in Q$, then
		$
		%\begin{equation}\label{constr_sync:eqn:Gamma_suffix}
		\applytrans{\delta'}{(q,p)}{w} = (\applytrans{\delta}{p}{w}, \applytrans{\delta}{p}{w}).
		%
		%\end{equation}
		$
				
		\begin{proofitem}{\ref{constr_sync:item:Gamma_suffix}} 
		By assumption, we can write $w = w_0 u_1 w_1 \cdots u_n w_n$ for some $n \ge 0$
		with $w_0, w_1, \ldots, w_n \in (\Sigma \setminus \Gamma)^+$
		and $u_1, \ldots, u_n \in \Gamma^*$. For $n = 0$,
		the claim follows by an easy induction from the defining equation of~$\delta'$.
		Next, suppose $n \ge 1$. Since $|w_{n-1}| \ge 1$,
		inductively, we obtain
		\[
		\delta'((q,p), w') 
		= (\delta(p, w'), \delta(p, w'))
		\]
		for $w' = w_0 u_1 w_1 \cdots w_{n-1}$. By Property~\ref{constr_sync:item:word_in_Gamma}
		we have
		\[
		\delta'((\delta(p, w'), \delta(p, w')), u_n) = (\delta(p, w'), \delta(p, w'u_n)).
		\]
		Then applying the claim for $n = 0$ again, we find 
		\[ 
		\delta'((\delta(p, w'), \delta(p, w'u_n)), w_n) = (\delta(p, w'u_n w_n), \delta(p, w'u_n w_n)).
		\]
		Thus we have
		$
		\applytrans{\delta'}{(q,p)}{w} = (\applytrans{\delta}{p}{w}, \applytrans{\delta}{p}{w}).
		$		
		\end{proofitem}		
	
		\item \label{constr_sync:item:QxQ}  For $v \in L \cap \Sigma^+$, 
		we have  $\delta'(Q', v) = \{(q,q) \mid q \in \delta(Q, v) \}$.
		
		\begin{proofitem}{\ref{constr_sync:item:QxQ}}
		By assumption on~$L$,
		we have $L \cap \Sigma^+ \subseteq \Sigma^*( \Sigma \setminus \Gamma )$
		and so $v \in \Sigma^* (\Sigma \setminus \Gamma)$.
		Thus, the claim follows from Property~\ref{constr_sync:item:Gamma_suffix}.		
		\end{proofitem}
	\end{enumerate}	
	
	\medskip 
	\noindent\underline{Verification:} We can assume $|Q| \ge 2$.
	First, suppose $|\delta(Q, v) | = 1$ for some $v \in L$.
	Then, by Property~\ref{constr_sync:item:QxQ}, $|\delta'(Q \times Q, v)| = 1$.
	So, $|\delta'(Q', vu)| = 1$ for $u \in U$.	
	
	Conversely, assume $|\delta'(Q', vu)| = 1$
	for $v \in L$ and $u \in U$. %, where we can assume $|v| > 0$.
%	By Property~\ref{constr_sync:item:Gamma_suffix}, as $v \notin \Sigma^*\Gamma$,
%	we have $\delta'(Q\times Q, v) = \{(q,q) \mid q \in \delta(Q, v) \}$.
	By Property~\ref{constr_sync:item:word_in_Gamma}, for $p,q \in Q$ and $u \in \Gamma^*$,
	we have $\delta'((q,q), u) = (q, \delta(q, u))$ and $\delta'((p,p), u) = (p, \delta(p, u))$.
	So,
	\[
	\delta'((q,q), u) = \delta'((p,p), u)
	\Leftrightarrow 
	q = p.
	\]
	Hence, with Property~\ref{constr_sync:item:QxQ}, $|\delta'(Q', vu)| = 1$
	implies $|\delta'(Q', v)| = 1$,
	which, again by Property~\ref{constr_sync:item:QxQ}, then implies $|\delta(Q, v)| = 1$.	
%	
%	
%	
%	
%	
%	
%	\bigskip 
%	
%	Let $v \in L$ and $u \in U$. We can assume $|v| > 0$.
%	By Property~\ref{constr_sync:item:Gamma_suffix}, as $v \notin \Sigma^*\Gamma$,
%	we have $\delta'(Q\times Q, v) = \{(q,q) \mid q \in \delta(Q, v) \}$.
%	So,  $|\delta(Q, v) | = 1$,
%	implies $|\delta'(Q \times Q, vu)| = 1$.
%	
%	%By Property~\ref{} for $v \in \Gamma^*$ 
%	%we have $\delta'(S \times T, v) = \bigcup_{q \in S} \{q\} \times \delta(T, v)$.
%	
%	Furthermore, by Property~\ref{constr_sync:item:word_in_Gamma}, as $u \in \Gamma^*$, we have $(q, \delta(q, u)) = \delta'((q,q), u)$ and $(p, \delta(p, u)) = \delta'((p,p), u)$,
%	and so $\delta'((q,q), u) = \delta'((p,p), u) = (p, \delta(u, p))$
%	if and only if $q = p$ for $q,p \in Q$. This yields
%	that $|\delta'(Q \times Q, vu)| = 1$
%	implies $|\delta'(Q \times Q, v)| = 1$,
%	which then implies $|\delta(Q, v)| = 1$.
\end{proof} 

Lastly, we give a condition for preservation of the complexity w.r.t.~concatenation.

\begin{proposition} \label{constrsync:genresults:prop:concat}
	Let $L_1, L_2, L_3 \subseteq \Sigma^*$ be regular languages.
	If $L_1 L_2 L_3 \subseteq \fact{L_2}$, then $
	\constrsync{L_1 L_2 L_3} \eqACZEROm \constrsync{L_2}.
	$
\end{proposition}
\begin{proof}
	In general, it is $\fact{L_2} 
	\subseteq 
	\fact{L_1 L_2 L_3}$.
	By assumption, $\fact{L_1 L_2 L_3}
	\subseteq \fact{L_2}$.
	So, the claim follows from~\cref{gen_results:cor:factors_etc_complexity}.
\end{proof}

\begin{corollary}
 	For $u,v \in \Sigma^*$ and $L \subseteq \Sigma^*$ with $uLv \subseteq \fact{L}$
 	we have $\constrsync{uLv} \eqACZEROm \constrsync{L}$.
\end{corollary}

\section{Hardness Conditions}
\label{constr_sync:sec:hardness}

%!TEX root = ./../../springer-mono-format/book.tex
%!TEX root = ./../../thesis-trier.tex

\todo[inline]{subsection teile auch in eigene dateien}

Here we present \NL-, \PSPACE-, and \NP-hardness conditions in 
\cref{constr_sync:sec:NL-hardness,constr_sync:sec:PSPACE-hardness,constr_sync:sec:NP-hardness}.
Some of these conditions use the notion of auxiliary state avoiding words, which
we introduce in \cref{constrsync:sec:aux}.
\Cref{constrsync:sec:auxcodes} investigates the relation of auxiliary state avoiding
words to codes. The proofs also uses equivalence
relations presented in \cref{constrsync:sec:eqLn} 
that are derived from the Myhill-Nerode relation from \cref{prelim:subsubsec:min_aut}.


\subsection{Auxiliary State Avoiding Words}
\label{constrsync:sec:aux}

In our \PSPACE-hardness reductions in \cref{constr_sync:sec:NL-hardness,constr_sync:sec:PSPACE-hardness}
we construct, from a given automaton, a new automaton
having additional ``auxiliary'' states. For this new automaton
to be synchronizable, those additional states must be mapped to a synchronizing state as well.
Existence of certain (weakly/strongly) auxiliary state avoiding words 
guarantees that this is possible, which also explains why we call them auxiliary state avoiding.

Here, we introduce (weakly/strongly)
auxiliary state avoiding words
and collect various statements
related to closure under concatenation (\cref{constrsync:prop:auxstate-concat}),
weak auxiliary state avoidability w.r.t.~different languages (\cref{constrsync:prop:auxLn}),
the relation between strong and weak auxiliary state avoidability (\cref{constrsync:cor:weakly-vs-strongly,constrsync:cor:weakly_suffix_strongly}), 
partitions (\cref{constrsync:eqn:partitions}),
and a sufficient condition for weak auxiliary state avoidability (\cref{constr_sync:prop:aux_state_sigmaL}).

\IN{auxiliary state avoiding word!weakly}
\IN{auxiliary state avoiding word!strongly}
Let $L \subseteq \Sigma^*$. We call a word $w \in \Sigma^*$
\begin{enumerate}
	\item \emph{weakly auxiliary state avoiding (w.r.t.~$L$)}
	if $w \in (  \pref{L^*}\Sigma  \setminus \pref{L^*} )^+$,
	\item \emph{strongly auxiliary state avoiding (w.r.t.~$L$)}
	if $w \in ( ( \pref{L} \setminus L )\Sigma  \setminus \pref{L} )^+$.
\end{enumerate}

When the context is clear, we omit mentioning the language with respect
to which we consider a word weakly or strongly auxiliary state avoiding.

\begin{example}
\label{constrsync:ex:aux}
\begin{enumerate}
\item  Let $\Gamma \subseteq \Sigma$ be a proper subset.
Then $(\Sigma \setminus \Gamma)^+$ 
is the set of all strongly auxiliary state avoiding
words w.r.t.~$\Gamma$
and $(\Gamma^* ( \Sigma \setminus \Gamma ))^+$
the set of weakly auxiliary state avoiding words w.r.t.~$\Gamma$.

\item Every non-empty word is weakly auxiliary state avoiding w.r.t.~$\{\varepsilon\}$
and, as  $\emptyset^* = \{\varepsilon\}$, also w.r.t.~$\emptyset$.
However, no word is strongly auxiliary state avoiding w.r.t.~$\{\varepsilon\}$
nor w.r.t.~$\emptyset$.

\item If $L = \pref{L}$, then no word is strongly auxiliary state avoiding w.r.t.~$L$.	

\item By \cref{constrsync:prop:prefix-maximal} no word is weakly or strongly auxiliary state avoiding w.r.t.~a maximal
prefix code. (in fact, for prefix codes this property characterizes maximality.)
\end{enumerate} 
\end{example} 

The empty word is never weakly nor strongly auxiliary state avoiding for any language.

Weak (resp.~strong) auxiliary state avoidability w.r.t.~to a common language is preserved
under concatenation.

\begin{proposition}
	\label{constrsync:prop:auxstate-concat}
	Let $L \subseteq \Sigma^*$ and $u,v \in \Sigma^*$ be two weakly (resp.~strongly) auxiliary
	state avoiding words w.r.t.~$L$.
	Then $uv$ is weakly (resp.~strongly) auxiliary state avoiding w.r.t.~$L$.
\end{proposition}
\begin{proof}
	This follows 
	as the set of weakly auxiliary state avoiding words $(\pref{L^*} \Sigma \setminus \pref{L^*})^+$ (resp.~the set of strongly auxiliary state avoiding words $((\pref{L} \setminus L)\Sigma \setminus \pref{L})^+$)
	forms a semigroup, i.e., is closed for concatenation.
\end{proof}

Next, we show that being weakly auxiliary state avoiding
w.r.t.~$L, L^n$ for $n \ge 1$, or $L^*$ are all equivalent properties.

\begin{proposition}
	\label{constrsync:prop:auxLn}
	\label{constrsync:prop:weakly-n-strongly}
	
	Let $L \subseteq \Sigma^*$, $u \in \Sigma^*$, and $n \ge 1$. 
	Then the following are equivalent:
	\begin{enumerate}
		\item $u$ is weakly auxiliary state avoiding w.r.t.~$L$,
		
		\item $u$ is weakly auxiliary state avoiding w.r.t.~$L^n$,
		
		\item $u$ is weakly auxiliary state avoiding w.r.t.~$L^*$.		
	\end{enumerate}
\end{proposition} 
\begin{proof}
    All equivalences follow as $\pref{L^*} = \pref{(L^n)^*}$ and $(L^{*})^* = L^*$.
\end{proof}

Next, we give a statement that yields, in \cref{constrsync:cor:weakly-vs-strongly}
a relation between weak and strong auxiliary state avoidability
w.r.t.~$L$ and $L^n$, respectively, for some $n \ge 1$.

\begin{proposition} \label{constrsync:prop:weakly-vs-strongly} \label{constrsync:item:Ln1}
	Let $L \subseteq \Sigma^+$ be some language and $n \ge 1$.  Then we have
	$
	( \pref{L^*}\Sigma \setminus \pref{L^*} ) \cap \Sigma^{\le n} 
	= ( ( \pref{L^n} \setminus L^n ) \Sigma \setminus \pref{L^n} ) \cap \Sigma^{\le n}.
	$
\end{proposition}
\begin{proof}
	Let $u \in \Sigma^{\le n}$. By \cref{constrsync:eqn:Lm}
	we have $v \in \pref{L^*}$ if and only if $v\in \pref{L^n}$ for $v \in \Sigma^{\le n}$.
	Thus, $u \in \pref{L^*}\Sigma$
	and $u \notin \pref{L^*}$ 
	are equivalent to $u \in \pref{L^n}\Sigma$
	and $u\notin \pref{L^n}$.
	Since $\varepsilon \notin L$,
	we have that $v \in L^m$ implies $|v| \ge m$ for $m \ge 1$.
	So from $u \in \Sigma^{\le n}$
	we get that for $v \in \Sigma^*$ with $u \in v\Sigma$
	we have $|v| < n$ 
	and so $v \notin L^n$.
	Thus we find $u \in \pref{L^n}\Sigma$
	if and only if $u \in ( \pref{L^n} \setminus L^n )\Sigma$.
	This gives the claim.
\end{proof}

\begin{corollary} \label{constrsync:cor:weakly-vs-strongly} \label{constrsync:item:Ln2}
	Let $L \subseteq \Sigma^+$ be some language and $n \ge 1$.
	A word $u \in \Sigma^{\le n}$ is weakly auxiliary state avoiding w.r.t.~$L$
	if and only if it is strongly auxiliary state avoiding w.r.t.~$L^n$.	
\end{corollary}
\begin{proof}
	By~\ref{constrsync:prop:weakly-vs-strongly}
	we obtain $( \pref{L^*}\Sigma \setminus \pref{L^*} )^+ \cap \Sigma^{\le n}
	= (( \pref{L^n} \Sigma \setminus L^n ) \setminus \pref{L^n} )^+ \cap \Sigma^{\le n}$,
	which gives the claim.	
\end{proof}

As a consequence of the next statement, we obtain that for languages with $\varepsilon \notin L$
every weakly auxiliary state avoiding word w.r.t~$L$
has a suffix that is strongly auxiliary state avoiding w.r.t.~$L$.

\begin{proposition} \label{constrsync:prop:weakly_suffix_strongly}  \label{item:suffix1}
	Let $L \subseteq \Sigma^+$ be a language not containing the empty word. Then we have
	 \[
	 \pref{L^*}\Sigma  \setminus \pref{L^*} \subseteq  L^* ( ( \pref{L} \setminus L )\Sigma \setminus \pref{L} ).
	 \]
\end{proposition}
\begin{proof}
Let $u \in \pref{L^*}\Sigma  \setminus \pref{L^*}$.
Write $u = u_1 \cdots u_n x$ with
\begin{enumerate}
	\item $u_1, \ldots, u_{n-1} \in L$, 
	
	\item $x \in \Sigma$, and
	
	\item $u_n = \varepsilon$ or $u_n \in \pref{L} \setminus L$.
\end{enumerate}
We have $\varepsilon \in \pref{L} \setminus L$ since $\varepsilon \notin L$.
Thus, even for $u_n = \varepsilon$, we have $u_n \in \pref{L} \setminus L$.
Since $u \notin \pref{L^*}$ we have $u_nx \notin \pref{L}$.
Hence, $u \in L^* ( ( \pref{L} \setminus L )\Sigma \setminus \pref{L} )$.	
\end{proof}

\Cref{constrsync:prop:weakly_suffix_strongly} might fail
for languages containing the empty word as shown next.

\begin{example}
	Let $L = \{\varepsilon, ac\}$.
	Then $\pref{L} \setminus L = \{a\}$
	and $acb \in \pref{L^*}\Sigma \setminus \pref{L^*}$
	but $acb \notin L^* ( ( \pref{L} \setminus L )\Sigma \setminus \pref{L} )$.
\end{example} 

The inclusion from \cref{constrsync:prop:weakly_suffix_strongly} can be proper,
as the next example shows.

\begin{example}
	Let $L = \{b, babb, aba\}$.
	Then $babb \in L^* ( ( \pref{L} \setminus L )\Sigma \setminus \pref{L} ) \setminus ( \pref{L^*}\Sigma  \setminus \pref{L^*} )$.
\end{example} 

\begin{corollary} \label{constrsync:cor:weakly_suffix_strongly}  \label{item:suffix2}
	Let $L \subseteq \Sigma^+$. Every weakly auxiliary state avoiding word w.r.t.~$L$
	has a suffix that is strongly auxiliary state avoiding w.r.t.~$L$.
\end{corollary}
\begin{proof}
	By~\cref{item:suffix1}, we have
	$(\pref{L^*}\Sigma  \setminus \pref{L^*})^+ \subseteq ( L^* ( ( \pref{L} \setminus L )\Sigma \setminus \pref{L} ) )^+$.	
\end{proof}

In \cref{constrsync:prop:Cn-weakly-aux-state}
we will show that 
we have equality in \cref{constrsync:prop:weakly_suffix_strongly}
for prefix codes.



Next, in \cref{constrsync:eqn:partitions}, we give various partitions of $\Sigma^*$.
\Cref{constrsync:eqn:partitions:item:weak1,constrsync:eqn:partitions:strong}
involve weakly and strongly auxiliary state avoiding words, respectively.
The proof of \cref{constrsync:eqn:partitions:strong} is based on the following lemma.

\begin{lemma} \label{constrsync:lem:properpref}
	Let $L \subseteq \Sigma^*$.
	Set $L_1 = (\pref{L} \setminus L)\Sigma \setminus \pref{L}$
	and $L_2 =  L\Sigma  \setminus \pref{L}$.
	Then we have $\pref{L}\Sigma \setminus \pref{L} = L_1 \cup L_2$
	and the right ideals $L_1 \Sigma^*$
	and $L_2 \Sigma^*$ are disjoint.
\end{lemma}
\begin{proof}
	Clearly, we have $\pref{L}\Sigma = (\pref{L}\setminus L)\Sigma \cup L\Sigma$,
	which gives the first equation.		
	
	Next, let $u_1 x v = u_2 y w$
	with $u_1x \in L_1$, $u_2y \in L_2$, $x,y \in \Sigma$, and $v,w \in \Sigma^*$.
	Since $u_1, u_2 \in \pref{L}$ while $u_1x, u_2y \notin \pref{L}$,
	neither $u_1$ nor $u_2$ can be a proper prefix of the other.
	Indeed, if $|u_1| < |u_2|$, then $u_1x$ would be a prefix of $u_2$,
	and hence $u_1 x \in \pref{L}$, a contradiction.
	Similarly, $|u_2| < |u_1|$ leads to a contradiction. Therefore $|u_1| = |u_2|$,
	which gives $u_1 = u_2$ since $u_1 x v = u_2 y w$.
	However, by the definitions of~$L_1$ and~$L_2$,
	we have $u_1 \notin L$ and $u_2 \in L$, contradicting $u_1 = u_2$.	
\end{proof}

\begin{proposition} \label{constrsync:eqn:partitions}
	Let $L \subseteq \Sigma^*$. Then the following holds:
	\begin{align}
		\Sigma^* & = \pref{L} \uplus (  \pref{L} \Sigma  \setminus \pref{L} ) \Sigma^* \label{prelim:eqn:prefcomplement} \\			
		\Sigma^* & = \pref{L} \uplus (\pref{L}\Sigma \setminus \pref{L} )^+ \pref{L} \label{constrsync:eqn:partitions:item:strong0} \\
		\Sigma^* & = \pref{L^*} \uplus (\pref{L^*}\Sigma \setminus \pref{L^*})^+ \pref{L^*} \label{constrsync:eqn:partitions:item:weak1} \\
		\Sigma^* & = U \uplus (  ( \pref{L} \setminus L )\Sigma \setminus \pref{L} ) )^+ U	\label{constrsync:eqn:partitions:strong}
	\end{align}	
	where $U = \pref{L} \uplus ( L\Sigma \setminus \pref{L} ) \Sigma^*$.
\end{proposition}
\begin{proof}
\underline{\Cref{prelim:eqn:prefcomplement}}: Let $u \in \Sigma^*$ and 
let $v \in \Sigma^*$ be a longest prefix of~$u$
in $\pref{L}$. Then we either have $u = v$, and so $u \in \pref{L}$,
or $u = vxw$ with $x \in \Sigma$ and $w \in \Sigma^*$
such that $vx \notin \pref{L}$, i.e.,
$u \in ( \pref{L}\Sigma \setminus \pref{L} ) \Sigma^*$.

Since the prefix of every word in $\pref{L}$
is in $\pref{L}$, and each word in $( \pref{L}\Sigma \setminus \pref{L} ) \Sigma^*$
has a prefix not in $\pref{L}$, both sets are disjoint.

\underline{\Cref{constrsync:eqn:partitions:item:strong0}}:
The language $( \pref{L}\Sigma \setminus \pref{L} ) \Sigma^*$
does not contain the empty word.
Then Arden's lemma (\cref{prelim:lem:arden}) applied to \cref{prelim:eqn:prefcomplement}
gives 
\[
\Sigma^* = ( \pref{L}\Sigma \setminus \pref{L} )^* \pref{L}
= \pref{L} \cup ( \pref{L}\Sigma \setminus \pref{L} )^+ \pref{L},
\]
where disjointness of the union also follows since the union from \cref{prelim:eqn:prefcomplement} 
is disjoint.\footnote{The claim also follows from \cref{prelim:eqn:Sigmastarpartition}
applied to $\pref{L}\Sigma \setminus \pref{L}$
since $\pref{L} = \Sigma^* \setminus ( \pref{L}\Sigma \setminus \pref{L} )\Sigma^*$.}

\underline{\Cref{constrsync:eqn:partitions:item:weak1}}: This follows
from \cref{constrsync:eqn:partitions:item:strong0} applied to $L^*$.

\underline{\Cref{constrsync:eqn:partitions:strong}}: 
By \cref{constrsync:lem:properpref} we find
\[
\Sigma^* = U  \uplus ( (\pref{L} \setminus L)\Sigma \setminus \pref{L} )\Sigma^*.
\]
Then \cref{prelim:lem:arden} gives $\Sigma^* = (  ( \pref{L} \setminus L )\Sigma \setminus \pref{L} ) )^* U$.
The sets $\pref{L}$ and $(  ( \pref{L} \setminus L )\Sigma \setminus \pref{L} ) ) \Sigma^*$
are disjoint. By \cref{constrsync:lem:properpref}, 
$(L\Sigma \setminus \pref{L})\Sigma^*$
and $(\pref{L} \setminus L)\Sigma \setminus \pref{L} )\Sigma^*$
are disjoint. These sets are also pairwise disjoint. This gives the claim.\footnote{The claim also follows by applying~\cref{prelim:eqn:Sigmastarpartition} to $( \pref{L} \setminus L )\Sigma \setminus \pref{L}$,
	which gives
	\[
	\Sigma^* = U \cup ( \pref{L} \setminus L )\Sigma \setminus \pref{L} )^+  U
	\]
	with $U = \Sigma^* \setminus ( ( \pref{L} \setminus L )\Sigma \setminus \pref{L} )\Sigma^* ) = \pref{L} \cup ( L\Sigma \setminus \pref{L} ) \Sigma^*$.}
\end{proof}

Lastly, we give a sufficient condition to identify weakly auxiliary state
avoiding words in terms of the Levenshtein-Romanov mapping.

\begin{proposition} \label{constr_sync:prop:aux_state_sigmaL}
	Let $L \subseteq \Sigma^*$ and $u \in \Sigma^+$.
	If $\sigma_L(u) = \varepsilon$, then $u$ is weakly
	auxiliary state avoiding w.r.t.~$L$.	
\end{proposition} 
\begin{proof}
	As $\sigma_{L}(u) = \varepsilon$, $u$
	has no non-empty suffix in $\pref{L^*}$.
	So, by~\cref{constrsync:eqn:partitions:item:weak1}
	we find that  $u \in (\pref{L^*}\Sigma \setminus \pref{L^*})^+$.
\end{proof} 

The next example gives $L \subseteq \Sigma^*$ and weakly auxiliary state avoiding $u \in \Sigma^*$
with $\sigma_L(u) \ne \varepsilon$.

\begin{example}
	Let $L = \{bc, a\}$.
	Then~$ba$ is strongly and weakly auxiliary state avoiding w.r.t.~$L$, but $\sigma_L(ba) = a$.
	For another example, the word $aa$ is strongly and weakly auxiliary state
	avoiding w.r.t.~$\{ab,b\}$ and we have $\sigma_{\{ab,b\}}(aa) = a$.	
\end{example}

\subsection{Auxiliary State Avoiding Words and Prefix Codes}
\label{constrsync:sec:auxcodes}

In our applications, the notions of strongly or weakly auxiliary state avoiding words
occur in conjunction with prefix codes. Here,
we investigate relations to prefix codes.

\Cref{constrsync:prop:prefix-weakly} states that the sets of weakly and strongly auxiliary 
state avoiding words w.r.t.~some language
are generated by prefix codes. This follows from \cref{constrsync:prop:prefix-weakly0}.
In \cref{constrsync:prop:prefix-weakly-implies-strongly}
we prove that w.r.t.~a prefix code, every strongly auxiliary state avoiding
word is weakly auxiliary state avoiding.
The property of weak auxiliary state avoidability
is preserved w.r.t.~codes when concatenating 
code words to the front, as proved in \cref{constrsync:cor:Cn-weakly-aux-state}.
\Cref{constrsync:cor:Cn-weakly-aux-state} follows from \cref{constrsync:prop:Cn-weakly-aux-state}.
\cref{constrsync:prop:vu} collects 
equivalent conditions related to concatenating
words to the front while preserving weak auxiliary state avoidability.

\begin{proposition} \label{constrsync:prop:prefix-weakly0} 
	Let $L \subseteq \Sigma^*$.
	Then $\pref{L}\Sigma \setminus \pref{L}$
	is a prefix code if it is non-empty.
\end{proposition}
\begin{proof}
	This follows since every proper prefix of a word from $\pref{L}\Sigma$
	is contained in $\pref{L}$
	and therefore cannot belong to $\pref{L}\Sigma \setminus \pref{L}$.
\end{proof}
	
\begin{corollary}\label{constrsync:prop:prefix-weakly} 
	Let $L \subseteq \Sigma^*$.
	If the language $\pref{L^*}\Sigma \setminus \pref{L^*}$ (resp.~$(\pref{L} \setminus L)\Sigma \setminus \pref{L}$) is non-empty,
	then it is a prefix codes.
\end{corollary}
\begin{proof} 
	For the first language the claim
	follows from \cref{constrsync:prop:prefix-weakly0} 
	applied to the language~$L^*$.
	For the second language
	the claim follows since it is a subset
	of $\pref{L}\Sigma \setminus \pref{L}$
	and thus inherits the property of being a prefix code.
\end{proof}

By \cref{constrsync:prop:prefix-weakly} a factorization
of weakly (resp.~strongly) auxiliary state avoiding
words into words from $\pref{L^*}\Sigma \setminus \pref{L^*}$ 
(resp.~$( \pref{L} \setminus L )\Sigma \setminus \pref{L}$) is unique.

Let $C \subseteq \Sigma^+$ be a prefix code.
In this case, we have $C\Sigma \cap \pref{C} = \emptyset$.
Thus $(C\Sigma \setminus \pref{C}) \Sigma^* = C \Sigma^+$.
Using this in \cref{constrsync:eqn:partitions:strong} then gives
the following partition involving strongly state avoiding words:
\begin{equation} \label{constrsync:eqn:partitions:strong2}
	\Sigma^* = \pref{C} \uplus C\Sigma^+ \uplus ( ( \pref{C} \setminus C )\Sigma \setminus \pref{C})^+ ( \pref{C} \uplus C\Sigma^+ ).
\end{equation}

Next, we show that every strong auxiliary state avoiding word
w.r.t.~some prefix code is also weakly auxiliary state avoiding w.r.t.~the same prefix code.

\begin{proposition}
\label{constrsync:prop:prefix-weakly-implies-strongly} 
Suppose $C \subseteq \Sigma^*$ is a prefix code. 
Then every strongly auxiliary state avoiding
word w.r.t.~$C$ is also weakly auxiliary state avoiding w.r.t.~$C$.
\end{proposition}
\begin{proof}
By~\cref{constrsync:eqn:capprefCstar} we have
$
 ( ( \pref{C} \setminus C)\Sigma ) \setminus \pref{C} 
 = ( ( \pref{C} \setminus C)\Sigma ) \setminus \pref{C^*}. 
$
As $\pref{C} \setminus C)\Sigma \subseteq \pref{C^*}\Sigma$
we find that $(\pref{C} \setminus C)\Sigma) \setminus \pref{C} )$
is contained in $\pref{C^*} \Sigma \setminus \pref{C^*}$,
which gives the claim.
\end{proof}

The following example shows that \cref{constrsync:prop:prefix-weakly-implies-strongly} does not hold for arbitrary codes.

\begin{example}
	Consider the code $C = \{a,aab,c\}$.
	Then $aac \in (\pref{C} \setminus C)\Sigma \setminus \pref{C}$
	but $aac \in C^*$.
	By \cref{constrsync:eqn:partitions:item:weak1}
	we have $aac \notin ( \pref{C^*}\Sigma \setminus \pref{C^*} )^+$
	and so $aac$ is not weakly auxiliary state avoiding w.r.t.~$C$.
	For another example (that is not a code),
	consider $L = \{a,b,aa\}$.
	Then $ab \in (\pref{L}\setminus L)\Sigma \setminus \pref{L}$
	but $ab \in \pref{L^*}$.
\end{example}

\Cref{constr_sync:thm:uC} involved weakly auxiliary state avoiding words of the form
$wu$ for $w \in \Sigma^*$ and a weakly auxiliary state avoiding $u \in \Sigma^*$.
Next, we investigate conditions when weak auxiliary state avoidability
is preserved by concatenating words to the front.

\begin{proposition} \label{constrsync:prop:Cn-weakly-aux-state}
	Let $C \subseteq \Sigma^+$ be a prefix code.
	Then the following holds:
	\[
	\pref{C^*}\Sigma \setminus \pref{C^*} = C^* ( (\pref{C} \setminus C)\Sigma \setminus \pref{C} )
	\]	
\end{proposition}
\begin{proof}
	By \cref{constrsync:eqn:capprefCstar} we have 
	$(\pref{C} \setminus C)\Sigma \setminus \pref{C} = (\pref{C} \setminus C)\Sigma \setminus \pref{C^*}$.
	Since $\pref{C}\setminus C \subseteq \pref{C^*}$,
	we have
	\[
	(\pref{C} \setminus C)\Sigma \setminus \pref{C^*} \subseteq \pref{C^*}\Sigma \setminus \pref{C^*}.
	\]
	Thus $C^* ( (\pref{C} \setminus C)\Sigma \setminus \pref{C^*} )
	\subseteq C^* ( \pref{C^*}\Sigma \setminus \pref{C^*} )$.
	Since  $C^* \pref{C^*}\Sigma \subseteq  \pref{C^*}\Sigma$,
	by \cref{constrsync:eqn:complementprefCstar} 
	we have $C^* ( \pref{C^*}\Sigma \setminus \pref{C^*} ) 
	= \pref{C^*}\Sigma \setminus \pref{C^*}$.
	The other inclusion follows from \cref{constrsync:prop:weakly_suffix_strongly}.
\end{proof}

\begin{corollary} \label{constrsync:cor:Cweakly-aux-state0}
	Let $C \subseteq \Sigma^+$ be a prefix code. Then we have
	\begin{align*}
		(\pref{C}\setminus C)\Sigma \setminus \pref{C} & \subseteq \pref{C^*}\Sigma \setminus \pref{C^*} \\
			\pref{C^*}\Sigma \setminus \pref{C^*} & = C^* ( \pref{C^*}\Sigma \setminus \pref{C^*} )
	\end{align*}	
\end{corollary}
\begin{proof}
	The first equation clearly follows from \cref{constrsync:prop:Cn-weakly-aux-state}.
	We have $C^* C^* ( (\pref{C} \setminus C)\Sigma \setminus \pref{C} )
	= C^* ( (\pref{C} \setminus C)\Sigma \setminus \pref{C} )$, which gives
	the second equation with \cref{constrsync:prop:Cn-weakly-aux-state}.	
\end{proof}

\begin{corollary}  \label{constrsync:cor:Cn-weakly-aux-state}
	Let $C \subseteq \Sigma^+$ be a prefix code
	and $u$ weakly or strongly state avoiding w.r.t.~$C$.
	Then $vu$ is weakly state avoiding w.r.t.~$C$ for each $v \in C^*$.
\end{corollary}
\begin{proof}
	We can write $u = u_1 \cdots u_n$
	where, for $i \in \{1,\ldots,n\}$,
	we have $u_i \in \pref{C^*}\Sigma \setminus \pref{C^*}$ 
	if $u$ is weakly auxiliary state avoiding
	and $u_i \in (\pref{C}\setminus C)\Sigma \setminus \pref{C}$
	if $u$ is strongly auxiliary state avoiding. 
	By \cref{constrsync:prop:Cn-weakly-aux-state,constrsync:cor:Cweakly-aux-state0} we have
	\[
	wu_1 \in \pref{C^*}\Sigma \setminus \pref{C^*}
	\]
	and
	\[
	u_i \in \pref{C^*}\Sigma \setminus \pref{C^*}
	\]
	for $i \in \{2,\ldots,n\}$ for both cases that $u$ is strongly or weakly
	auxiliary state avoiding.
	Thus $wu_1 \cdots u_n$ is weakly auxiliary state avoiding.
\end{proof}

\begin{example}
	\Cref{constrsync:prop:Cn-weakly-aux-state} 
	is not valid for arbitrary languages.
	For example consider $L = \{aa,bc,aabd\}$.
	Then $bd$ is weakly auxiliary state avoiding w.r.t.~$L$,
	but $aabd \in L$ is not weakly auxiliary state avoiding 
	w.r.t.~$L$.
\end{example}

\begin{corollary}
		Let $C \subseteq \Sigma^+$ be a prefix code.
		Then $( C \cup \pref{C^*}\Sigma \setminus \pref{C^*} )^* = (\pref{C^*}\Sigma \setminus \pref{C^*} )^*C^*$.
\end{corollary}
\begin{proof} 
	By \crefitem{prelim:langeqn}{prelim:langeqn:starcup}
	we have 
	$( C \cup \pref{C^*}\Sigma \setminus \pref{C^*} )^* = (C^* ( \pref{C^*}\Sigma \setminus \pref{C^*} ) )^* C^*$.
	By \cref{constrsync:eqn:complementprefCstar} 
	the last language equals 
	$( \pref{C^*}\Sigma \setminus \pref{C^*} )^* C^*$.
\end{proof} 

The following examples show that concatenating a word
from the set of prefixes to a weakly auxiliary state avoiding
word does not, in general, yield a weakly auxiliary state
avoiding word even w.r.t.~prefix codes.

\begin{example}
\begin{enumerate}
\item Consider the prefix code $C = \{aa,bac\}$.
Then $ac$ is weakly (and strongly) auxiliary state avoiding w.r.t.~$C$,
but $bac \in C$ is not weakly auxiliary state avoiding w.r.t.~$C$.

\item Consider the prefix code $C = aaa^*b$. 
Here,  $\pref{C^*}$
is a submonoid of $\Sigma^*$ 
and $ab$ is strongly and weakly 
auxiliary state avoiding w.r.t.~$C$. However, for  $a \in \pref{C^*}$ we have $aab \in C$
and so $aab$ is neither weakly nor strongly auxiliary state avoiding w.r.t.~$C$.
\end{enumerate}
\end{example}

Next, we show that the property of preservation of weak auxiliary state avoidability (for a fixed word)
is equivalent for certain types of prefixes.

\begin{proposition}	\label{constrsync:prop:vu}
	Let $C \subseteq \Sigma^*$ be a prefix code, $u \in \Sigma^*$ be weakly auxiliary state avoiding w.r.t.~$C$,
	and $n \ge 1$ be some fixed number.
	Then the following are equivalent.	 
	\begin{enumerate}
	\item \label{constrsync:aux:v} $vu$ is weakly auxiliary state avoiding  w.r.t.~$C$ 
	for each $v \in \pref{C} \setminus C$.		
		
	\item \label{constrsync:aux:vCstar} $vu$ is weakly auxiliary state avoiding  w.r.t.~$C$ for each $v \in \pref{C^*}$.
		
	\item \label{constrsync:aux:vCn} $vu$ is weakly auxiliary state avoiding  w.r.t.~$C$ for each $v \in \pref{C^n}$.
	\end{enumerate} 
\end{proposition} 
\begin{proof}
	$\ref{constrsync:aux:v} \Rightarrow~\ref{constrsync:aux:vCstar}$: 
	Set $L = \pref{C^*}\Sigma \setminus \pref{C^*}$ so that $L^+$ is the set of weakly auxiliary state avoiding words w.r.t.~$C$. 
	By \cref{prelim:eqn:prefLstarnoeps} 
	we have $\pref{C^*} = C^* (\pref{C} \setminus C)$.
	By \cref{constrsync:aux:v}
	we have $(\pref{C} \setminus C) L^+ \subseteq L^+$.
	By \cref{constrsync:cor:Cn-weakly-aux-state} 
	we have $C^* L^+ \subseteq L^+$.
	Thus we obtain  $\pref{C^*}L^+ \subseteq L^+$,
	which gives the claim.
	
	$\ref{constrsync:aux:vCstar} \Rightarrow~\ref{constrsync:aux:vCn}$: This is clear
	since $\pref{C^n} \subseteq \pref{C^*}$.
	
	$\ref{constrsync:aux:vCn} \Rightarrow~\ref{constrsync:aux:v}$: This is clear
	since $\pref{C} \setminus C \subseteq \pref{C^n}$.
\end{proof}

\subsection{Maximal Prefix Codes}
\label{constrsync:sec:max-prefix-codes}

\IN{maximal prefix code}
\IN{prefix code!maximal}
We call~$C \subseteq \Sigma^+$ \emph{maximal} if $C \cup \{u\}$ is not a prefix code for each $u \in \Sigma^* \setminus C$.
%Maximal codes are related to the defining conditions of weak and strong auxiliary state avoidability in the following way.

Auxiliary state avoiding words are related to maximal codes. Here, we demonstrate this by deriving
characterizations of maximal codes from the results of \cref{constrsync:sec:auxcodes,constrsync:sec:aux}.

\begin{proposition} \label{constrsync:prop:prefix-maximal}
	Let $C \subseteq \Sigma^+$ be a prefix code. Then the following are equivalent:
	\begin{enumerate}
		\item \label{constrsync:prop:prefix-maximal:1} $C$ is maximal.
		
		\item \label{constrsync:prop:prefix-maximal:2} $\Sigma^* = \pref{C} \uplus C\Sigma^+$.
		
		\item \label{constrsync:prop:prefix-maximal:3} $(\pref{C} \setminus C)\Sigma \setminus \pref{C} = \emptyset$.
		
		\item \label{constrsync:prop:prefix-maximal:4} $\pref{C^*}\Sigma  \setminus \pref{C^*} = \emptyset$.
		
		\item \label{constrsync:prop:prefix-maximal:5} $\pref{C^*} = \Sigma^*$.
	\end{enumerate}	
\end{proposition}
\begin{proof}
	$\ref{constrsync:prop:prefix-maximal:1} \Rightarrow \ref{constrsync:prop:prefix-maximal:2}$: Let $u \in \Sigma^*$.
	For $u \in C$, we have $u \in \pref{C}$.
	For $u \notin C$,
	since~$C$ is maximal,
	$C \cup \{u\}$ is not a prefix code.
	This yields that some word in~$C$
	has~$u$ as a prefix, i.e., $u \in \pref{C} \setminus C$,
	or that $u$ has some word from~$C$
	as a prefix, i.e., $u \in C\Sigma^+$.
	This yields $\Sigma^* = \pref{C} \uplus C\Sigma^+$.
	
	$\ref{constrsync:prop:prefix-maximal:2} \Rightarrow \ref{constrsync:prop:prefix-maximal:3}$: Since $\Sigma^* = \pref{C} \uplus C\Sigma^+$
	and the union from \cref{constrsync:eqn:partitions:strong2}
	is disjoint,
	we must have $( ( \pref{C} \setminus C )\Sigma \setminus \pref{C})^+ ( \pref{C} \uplus C\Sigma^+ ) = \emptyset$,
	which gives $(\pref{C} \setminus C)\Sigma \setminus \pref{C} = \emptyset$.
	
	$\ref{constrsync:prop:prefix-maximal:3} \Rightarrow \ref{constrsync:prop:prefix-maximal:4}$: By \cref{constrsync:prop:weakly_suffix_strongly} 
	every word in $\pref{C^*}\Sigma  \setminus \pref{C^*}$
	has a suffix in $(\pref{C} \setminus C)\Sigma \setminus \pref{C}$.
	However, this latter set is empty by assumption.
	Thus, there can be no word in $\pref{C^*}\Sigma \setminus \pref{C^*}$, i.e.,
	this set is empty as well.
	
	$\ref{constrsync:prop:prefix-maximal:4} \Rightarrow \ref{constrsync:prop:prefix-maximal:5}$: This follows from the disjoint union given in
	\cref{constrsync:eqn:partitions:item:weak1} since the 
	set $(\pref{L^*}\Sigma \setminus \pref{L^*})^+ \pref{L^*}$
	is empty by assumption.
	
	$\ref{constrsync:prop:prefix-maximal:5} \Rightarrow \ref{constrsync:prop:prefix-maximal:1}$: Suppose~$C$ is not maximal.
	Then $C \cup \{u\}$ is a prefix code for some $u \in \Sigma^*$
	with $u \notin C$.
	Since~$u$ has no prefix in~$C$ it is not contained in $C\pref{C^*}$
	and since no word from~$C$ has it as a prefix,
	it is not contained in $\pref{C}$.
	We have $\pref{C^*} = \pref{C} \cup C\pref{C^*}$
	and thus $u \notin \pref{C^*}$.
\end{proof}

\begin{remark}
	Let $C \subseteq \Sigma^+$ be a prefix code.
	The set $( (\pref{C} \setminus C)\Sigma ) \setminus \pref{C}$
	is precisely the set of minimal words
	that are not prefixes of words from~$C$.
\end{remark}

\begin{proposition} \label{constrsync:prop:prefix-maximalD}
	Let $C \subseteq \Sigma^+$ be a prefix code.
	Then $D = C \cup ( ( \pref{C} \setminus C ) \Sigma \setminus \pref{C} )$
	is a maximal prefix code.
\end{proposition}
\begin{proof}
	Since $C$ and~$( \pref{C} \setminus C ) \Sigma \setminus \pref{C}$
	are prefix codes
	and no word from~$C$ can be a prefix of a word
	from $( \pref{C} \setminus C ) \Sigma \setminus \pref{C}$
	and vice versa, we find that $D$ is a prefix code.
	
	Next, by \cref{constrsync:eqn:partitions:strong2}, we have
	\begin{align*}
	\Sigma^* & \subseteq \pref{C} \cup C\Sigma^+ \cup 
	((\pref{C}\setminus C)\Sigma \setminus \pref{C})\Sigma^* \\
	& = \pref{C} \cup ((\pref{C}\setminus C)\Sigma \setminus \pref{C}) 
	\cup C\Sigma^+ \cup ((\pref{C}\setminus C)\Sigma \setminus \pref{C})\Sigma^+ \\
	& \subseteq \pref{D} \cup D\Sigma^+.
	\end{align*}
	Thus $\Sigma^* = \pref{D} \cup D\Sigma^+$.
	Since~$D$ is a prefix code, this union is disjoint.
	So $D$ is maximal by \cref{constrsync:prop:prefix-maximal:2} of \cref{constrsync:prop:prefix-maximal}.	
\end{proof}

\begin{example}
	Let $C = a^i a^* b$ for some $i \ge 0$.
	Then $\pref{C} \setminus C = a^*$,
	$(\pref{C} \setminus C )\{a,b\} = \{ b, ab, aab, \ldots, a^{i-1}b \} \cup C \cup aa^*$,
	and $(\pref{C} \setminus C )\{a,b\} \setminus \pref{C} = 
	\{ b, ab, aab, \ldots, a^{i-1}b \}$.	
\end{example}

\begin{example}
	Let $C = \{ab,aab,aaab\}$.
	Then $\pref{C} \setminus C = \{ \varepsilon,a,aa,aaa\}$,
	$( \pref{C} \setminus C )\{a,b\} = \{ a, b, aa, ab, aaa, aab, aaaa, aaab \}$,
	and $( \pref{C} \setminus C )\{a,b\} \setminus \pref{C}
	= \{ b, aaaa \}$.
\end{example}

\begin{remark}
	Note that, in general, there might be several possibilities to obtain a maximal
	prefix code by extending a given prefix code~$C$.
	For example, for $C = \{b, ab, aab\}$,
	$D_i = \{ aaa^jb \mid 1 \le j \le i - 1 \}\cup \{ a^{i+2} \}$ for $i \ge 1$,
	and $D = \{ aaab, aaaab, aaaaab, \ldots \}$
	the sets $C \cup D_i$ for $i \ge 1$
	and $C \cup D$ all give maximal prefix codes.
	Note that here $D_1 = (\pref{C} \setminus C)\Sigma \setminus \pref{C}$.	
\end{remark}

\subsection{Relations Derived From the Myhill-Nerode Relation} 
\label{constrsync:sec:eqLn}

Here, we introduce relations derived from the Myhill-Nerode relation $\equiv_L$,
which was introduced in \cref{prelim:subsubsec:min_aut}. We first abstractly describe 
a method for constructing a new relation that refines the Myhill-Nerode relation from a given relation $\sim$ 
 when certain assumptions are satisfied.
We then derive various properties of this new relation.
Via this abstract construction we obtain a relation~$\eqLL{C}{V}$,
which will be central to most of our hardness proofs
given in the remaining sections since it allows 
to  distinguish, or ``keep track'', of certain words in automata that we construct in our reductions.

Let $L \subseteq \Sigma^*$. Suppose we have an equivalence
relation $\sim$ on~$\Sigma^*$ and a language $U \subseteq \Sigma^*$
such that, for all $u,v \in \Sigma^*$ and $x \in \Sigma$,
the following properties hold:
\begin{enumerate}
\item \label{constrsync:rel:refinesNerode} If $u \sim v$, then $u \equiv_L v$, i.e., $\sim$ refines $\equiv_L$.

\item \label{constrsync:rel:saturateU} The relation~$\sim$ saturates~$U$.

\item \label{constrsync:rel:right-congr2} If $u \equiv_L v$ and $u,v \in U$, then $ux,vx \in U$ or $ux \sim vx$.

\item \label{constrsync:rel:right-congr1} If $u \sim v$, then $ux \sim vx$ or $ux,vx \in U$.
\end{enumerate}
Then for $u,v \in \Sigma^*$ define a relation $\eqLsimU{L}{\sim}{U}$ by
\[
u \eqLsimU{L}{\sim}{U} v
\Leftrightarrow 
( u \sim v ) \lor ( u \equiv_L v \land u,v \in U ).
\]
\label{def:Nerode-generalization}

Before establishing some properties of this relation, we present two examples to build intuition.

\begin{example} \label{constrsync:ex:refineNerode}
	\begin{enumerate} 		
		\item  \label{constrsync:ex:refineNerode:right-ideal} Properties~\ref{constrsync:rel:refinesNerode},
		\ref{constrsync:rel:saturateU},\ref{constrsync:rel:right-congr2}, and~\ref{constrsync:rel:right-congr1}
		are obviously fulfilled when $\sim$ equals the equality relation
		and $U$ is a right ideal.
		
		% In dem Fall das zeigen:
		% There are finitely many equivalences classes for $\equiv_{L,U}$
		% if and only if~$L$ is regular and~$\Sigma^* \setminus U$ is finite.
		
		\item For $U = \Sigma^*$, the above construction always yields
		the Myhill-Nerode relation $\equiv_L$, regardless of which relation $\sim$ (that refines $\equiv_L$) we choose.
	\end{enumerate}
\end{example}

\begin{remark} \label{constrsync:rem:refined-Nerode}
	Since~$\sim$ refines $\equiv_L$, for $u,v \in \Sigma^*$, we have
	\[
	u \eqLsimU{L}{\sim}{U} v
	\Leftrightarrow
	(  u \sim v  \land |\{u,v\} \cap U| \le 1 ) \lor
	( u \equiv_L v \land u,v \in U ).
	\]
	In the special case that~$\sim$ is the equality relation,
	we also have
	\[
	u \eqLsimU{L}{\sim}{U} v
	\Leftrightarrow
	( u = v  \land u,v \notin U ) \lor
	( u \equiv_L v \land u,v \in U ).
	\]	
\end{remark}

The relation $\eqLsimU{L}{\sim}{U}$ is an equivalence relation that is invariant under right concatenation
as stated next.

\begin{theorem} \label{constrsync:thm:relation-right-congruence}
The relation $\eqLsimU{L}{\sim}{U}$ is an equivalence
relation such that, for all $u,v \in \Sigma^*$,
if $u \eqLsimU{L}{\sim}{U} v$,
then $ux \eqLsimU{L}{\sim}{U} vx$ for all $x \in \Sigma$.
\end{theorem}
\begin{proof}
%Clearly, the relation is symmetric and reflexive.
Symmetry and reflexivity are clear.
For transitivity, consider $u,v,w \in \Sigma^*$
such that $u \eqLsimU{L}{\sim}{U} v$
and $v \eqLsimU{L}{\sim}{U} w$.
If $u \sim v$ and $v \sim w$, then $u \sim w$
and so $u \eqLsimU{L}{\sim}{U} w$.
For $u \equiv_L v$ and $v \equiv_L w$ with $u,v,w \in U$
we have $u \equiv_L w$ 
since $\equiv_L$ is transitive and so $u \eqLsimU{L}{\sim}{U} w$.
Now, suppose $u \sim v$ and $v \equiv_L w$ with $v,w \in U$.
By Property~\ref{constrsync:rel:refinesNerode},
$u \sim v$ implies $u \equiv_L v$.
Since~$\sim$ 
saturates~$U$ by Property~\ref{constrsync:rel:saturateU},
we also have $u \in U$. Thus we find $u \equiv_L w$
and $u,w \in U$, i.e., $u \eqLsimU{L}{\sim}{U} w$.
Next, suppose $u \equiv_L v$ with $u,v \in U$
and $v \sim w$. By Property~\ref{constrsync:rel:refinesNerode}
we have $v \equiv_L w$. Since~$\sim$ 
saturates~$U$ by Property~\ref{constrsync:rel:saturateU},
we also have $w \in U$.
Thus we find $u \equiv_L w$
and $u,w \in U$, i.e., $u \eqLsimU{L}{\sim}{U} w$.

Next, suppose $u \eqLsimU{L}{\sim}{U} v$ for $u,v \in \Sigma^*$
and let $x \in \Sigma$. We show that then we have $ux \eqLsimU{L}{\sim}{U} vx$.
By combining Property~\ref{constrsync:rel:refinesNerode}
with \cref{prelim:eqn:Nerodecongruence}
and using Property~\ref{constrsync:rel:right-congr1}, we get
\[
u \sim v 
\Rightarrow 
( ux \sim vx ) \lor ( ux \equiv_L vx \land ux, vx \in U )
\]
Thus, if $u \sim v$, we have $ux \eqLsimU{L}{\sim}{U} vx$.
Next, let $u \equiv_L v$ with $u,v \in U$.
If $ux \sim vx$, then we have $ux \eqLsimU{L}{\sim}{U} vx$ 
by the definition of $\eqLsimU{L}{\sim}{U}$.
Otherwise, by Property~\ref{constrsync:rel:right-congr2},
we have $ux, vx \in U$. Since, by \cref{prelim:eqn:Nerodecongruence}, 
from $u \equiv_L v$ we get $ux \equiv_L vx$,
we also have $ux \eqLsimU{L}{\sim}{U} vx$. This shows the claim.
\end{proof}

\begin{proposition} \label{constrsync:prop:equiv-classes-form}
For $u \in \Sigma^*$ we have
\[
[u]_{\eqLsimU{L}{\sim}{U}}
= 
\begin{cases}
	[u]_{\sim}               & \text{if } u \notin U, \\
	[u]_{\sim} \cup ( [u]_{\equiv_L} \cap U ) & \text{if } u \in U.
\end{cases}
\]
\end{proposition}
\begin{proof}
Let $v \in [u]_{\eqLsimU{L}{\sim}{U}}$.
First, suppose $u \notin U$.
Then, by definition,
we must have $u \sim v$, i.e., $v \in [u]_{\sim}$.
Next, suppose $u \in U$
and $v \notin [u]_{\sim}$.
In this case, by definition,
we have $u \equiv_L v$ and $v \in U$, i.e.,
$v \in [u]_{\equiv_L} \cap U$.

Conversely, by definition,
we always have $[u]_{\sim} \subseteq [u]_{\eqLsimU{L}{\sim}{U}}$.
If also $u \in U$ and $v \in [u]_{\equiv_L} \cap U$,
then $u \equiv_L v$ and $u,v \in U$
so that $v \in [u]_{\eqLsimU{L}{\sim}{U}}$.
\end{proof}

\begin{proposition} \label{constrsync:prop:relation-refines}
Let $V \subseteq \Sigma^*$ be such that $U \subseteq V$
and $\approx$ be an equivalence
relation that is coarser than~$\sim$ and that fulfills
Properties~\ref{constrsync:rel:refinesNerode},
\ref{constrsync:rel:saturateU},\ref{constrsync:rel:right-congr2}, and~\ref{constrsync:rel:right-congr1} w.r.t.~$L$ and~$V$. % vielleicht begriff \sim und U compatible mit L oder so einführen
Then for $u,v \in \Sigma^*$ we have that
$u \eqLsimU{L}{\sim}{U} v$
implies
$u \eqLsimU{L}{\approx}{V} v$, i.e., $\eqLsimU{L}{\sim}{U}$ refines $\eqLsimU{L}{\approx}{V}$.
\end{proposition}
\begin{proof}
Suppose $u \eqLsimU{L}{\sim}{U} v$ for $u,v \in \Sigma^*$.
Then we have $u \sim v$
or $u \equiv_L v$ with $u,v \in U$.
In the former case, since $u \sim v$
implies $u \approx v$,
we also have $u \eqLsimU{L}{\approx}{V} v$.
In the latter case, since $U \subseteq V$,
we find $u,v \in V$
and with $u\equiv_L v$
we also have $u \eqLsimU{L}{\approx}{V} v$ .
\end{proof}

\begin{corollary} \label{constrsync:cor:relation-refines_Nerode}
The relation $\eqLsimU{L}{\sim}{U}$
refines $\equiv_L$.
\end{corollary} 
\begin{proof}
The equivalence relation $\equiv_L$ and the set~$\Sigma^*$
fulfill all Properties~\ref{constrsync:rel:refinesNerode},
\ref{constrsync:rel:saturateU},\ref{constrsync:rel:right-congr2}, and~\ref{constrsync:rel:right-congr1}.
Now $\eqLsimU{L}{\equiv_L}{\Sigma^*}$
equals~$\equiv_L$ and so the claim follows from \cref{constrsync:prop:relation-refines}.
\end{proof}

\begin{corollary}  \label{constrsync:cor:both-prefix-L}
	The relation $\eqLsimU{L}{\sim}{U}$
	saturates~$L$ and $\pref{L}$.
\end{corollary}
\begin{proof}
This follows from \cref{constrsync:cor:relation-refines_Nerode} 	
since $\equiv_L$ saturates~$L$
and $\pref{L}$.
That $\equiv_L$ saturates $\pref{L}$
follows since if $u \equiv_L v$
and $u \in \pref{L}$,
then $uw \in L$ for some $w \in \Sigma^*$.
Now we have $uw \equiv_L vw$.
Thus, we also get $vw \in L$
and so $v \in \pref{L}$.
\end{proof}

\label{def:equivLn}
Let $n \ge 0$. For $L \subseteq \Sigma^*$ 
and $U = \Sigma^n \Sigma^*$
we now consider
the equivalence relation $\eqLsimU{L}{=}{U}$,
which we also denote by $\eqLn$ in the following.
With \cref{constrsync:rem:refined-Nerode}, for $u,v \in \Sigma^*$, we have 
\begin{align*} 
u \eqLn v
& \Leftrightarrow 
( u = v )
\lor 
( \min\{ |u|, |v| \} \ge n \land u \equiv_L v ) \\
& \Leftrightarrow 
( \max\{ |u|, |v| \} < n \land u = v )
\lor 
( \min\{ |u|, |v| \} \ge n \land u \equiv_L v ) \\
& \Leftrightarrow 
( \min\{ |u|, |v| \} < n \land u = v   )
\lor 
( \min\{ |u|, |v| \} \ge n \land u \equiv_L v )
\end{align*}

Clearly, for $m \ge n$ the relation 
$\eqLp{m}$ refines $\eqLp{n}$
and $\eqLp{0}$ equals $\equiv_L$.

\begin{proposition} \label{constrsync:prop:eqLn-classes}
Let $u \in \Sigma^*$.
Then we have
\[
[u]_{\eqLn}
=
\begin{cases}
	\{u\} & \text{if } |u| < n, \\
	\{ v \in \Sigma^n\Sigma^* \mid v \equiv_L u \} & \text{otherwise.}
\end{cases}
\]
Thus, we have 
at least $|\Sigma^{< n}| = ( |\Sigma|^n - 1 ) / (|\Sigma| - 1)$
many equivalence classes.
\end{proposition}
\begin{proof}
This follows directly from \cref{constrsync:prop:equiv-classes-form}
with $U = \Sigma^n \Sigma^*$ and the fact that $\sim$
equals the equality relation in the definition of~$\eqLn$.
\end{proof}

\begin{corollary}
There are finitely many equivalence classes for $\eqLn$
if and only if~$L$ is regular.
\end{corollary}
\begin{proof}
First, suppose~$L$ is regular.
Then, by \cref{prelim:min_aut:terminal_object},
we have finitely many equivalence classes for $\equiv_L$.
Thus there are only finitely many 
distinct sets of the form $\{ v \in \Sigma^n\Sigma^* \mid v \equiv_L u \}$
for $u \in \Sigma^*$ with $|u| \ge n$.
Since there are only finitely many words of length at most~$n$,
with \cref{constrsync:prop:eqLn-classes} 
we find that there are only finitely many equivalence classes 
for $\eqLn$.

Conversely, since $\eqLn$ refines $\equiv_L$,
the  $\equiv_L$-classes	are unions of $\eqLn$-classes. Thus,
if there are only finitely many $\eqLn$-classes,
then there are only finitely many $\equiv_L$-classes
and so~$L$ is regular by \cref{prelim:min_aut:terminal_object}.
\end{proof}

\label{def:preeqLL}
Next, let $C \subseteq \Sigma^+$ be a prefix code and $V \subseteq C$. 
Consider the relation $\preeqLL{C}{V}$ on~$\Sigma^*$
which is, for $u,v \in \Sigma^*$, given by
\[
u \preeqLL{C}{V} v \Leftrightarrow ( u = v ) \lor ( \exists w \in \pref{V} \setminus V \colon u,v \in C^*w ).
\]

\begin{proposition}
The relation~$\preeqLL{C}{V}$ given above is an equivalence relation.
\end{proposition}
\begin{proof} 
Clearly, the relation is reflexive and symmetric.
It is also transitive, for assume we have $u_1, u_2, u_3 \in \Sigma^*$
with $u_1 \preeqLL{C}{V} u_2$ and $u_2 \preeqLL{C}{V} u_3$, where we can suppose all three words are distinct.
Then we find $u_1, u_2 \in C^* w_1$ and $u_2, u_3 \in C^* w_2$
for some $w_1, w_2 \in \pref{V} \setminus V$.
This gives  $u_2 \in C^* w_1 \cap C^* w_2$.
We have $w_1, w_2 \notin C\Sigma^*$
as otherwise some word from~$V$ would have a word from~$C$
as a proper prefix, which is not possible since~$C$ is a prefix code.
As~$C$ is a prefix code,
by \cref{constrsync:prop:char-prefix-code,constrsync:prop:suffix-factorizations}
we find $w_1 = w_1$.
Thus $u_1, u_3 \in C^* w_1$, giving $u_1 \preeqLL{C}{V} u_3$.
\end{proof}

\begin{remark}
	Transitivity might fail if~$C$ is not a prefix
	code. For example, consider $C = \{aa,ba,baa\}$, which is only a code, but not a prefix code.
	Then $\pref{C} \setminus C = \{\varepsilon, a, b\}$.
	We have $aaa, baa \in C^*a$ and $baa, ba \in C^*$,
	but there exists no $w \in \{\varepsilon,a,b\}$
	such that $C^*w$ contains 
	$aaa$ and $ba$.
	Thus, the relation defined above 
	would not yield a transitive relation.
	However, there exist codes that are not prefix
	codes 
	such that
	$u_1, u_2 \in C^*w_1$ and $u_2, u_3 \in C^*w_2$
	for some $w_1, w_2 \in (\pref{C}\setminus C)$
	would imply $u_1, u_3 \in C^* w_3$ for some $w_3 \in \pref{C} \setminus C$,
	and so they would yield transitivity.
	For example, this holds for the code $C = \{a, ab\}$, 
	because $\pref{C} \setminus C = \{\varepsilon\}$.
	However, we only need this construction in the 
	context of prefix codes.	
\end{remark}

\begin{proposition} \label{constrsync:prop:Cstarw-relation-unique}
	Let $u,v \in \pref{V}\setminus V$.
	Then $u \preeqLL{C}{V} v$
	implies $u = v$.
\end{proposition}
\begin{proof}
	By the definition of~$\preeqLL{C}{V}$,
	from $u \preeqLL{C}{V} v$
	we get $u,v \in C^* w$ for some $w \in \pref{V}\setminus V$.
	Since~$C$ is a prefix code, which gives $u,v \notin C\Sigma^*$, 
	by \cref{constrsync:prop:char-prefix-code,constrsync:prop:suffix-factorizations},
	from $u \in C^* u \cap C^* w$ and $v \in C^* v \cap C^*w$
	we get $u = w = v$.
\end{proof}

\begin{proposition} \label{constrsync:prop:prefix-code-pre-equivalence} 
	Let $u \in \Sigma^*$.
	Then we have
	\[
	[u]_{\preeqLL{C}{V}} =
	\begin{cases}
		C^*w   & \text{if } u \in C^*w \text{ with } w \in \pref{V} \setminus V, \\
		\{u\}  & \text{otherwise.} 
	\end{cases}
	\]
\end{proposition}
\begin{proof}
	First, suppose $u \notin C^* (\pref{V}\setminus V)$.
	By reflexivity, we have $\{u\} \subseteq [u]_{\preeqLL{C}{V}}$.
	Next, let $v \in [u]_{\preeqLL{C}{V}}$.
	By the definition of~$\preeqLL{C}{V}$
	we have $u = v$. This gives $[u]_{\preeqLL{C}{V}} = \{u\}$.
	Now, let $u \in C^* w$ for some $w \in \pref{V}\setminus V$.
	In this case, by the definition of $\preeqLL{C}{V}$,
	we have $C^*w \subseteq [u]_{\preeqLL{C}{V}}$.
	Next, let $v \in [u]_{\preeqLL{C}{V}}$.	
	We can assume $u \ne v$. Then $u \preeqLL{C}{V} v$ implies $u,v \in C^* w'$ for some $w' \in \pref{V}\setminus V$.
	Since $u \in C^*w \cap C^*w'$ and $w' \notin C\Sigma^*$,
	by \cref{constrsync:prop:char-prefix-code,constrsync:prop:suffix-factorizations},
	we get $w = w'$ and so $v \in C^*w$.
	This gives $[u]_{\preeqLL{C}{V}} \subseteq C^* w$.
\end{proof}

\begin{proposition} \label{constrsync:prop:properties-simCV}
The relation~$\preeqLL{C}{V}$ fulfills Properties~\ref{constrsync:rel:refinesNerode},
\ref{constrsync:rel:saturateU}, \ref{constrsync:rel:right-congr2}, and~\ref{constrsync:rel:right-congr1}
for $L = C^*$ and $U = \Sigma^* \setminus ( C^* ( \pref{V} \setminus V) )$.
\end{proposition}
\begin{proof}
Observe that $u \equiv_{C^*} \varepsilon$ if and only if $u \in C^*$ (this in fact holds for the Kleene star of any language). We will use this equivalence in the following arguments.	

Property~\ref{constrsync:rel:refinesNerode}: Suppose $u \preeqLL{C}{V} v$ for $u,v \in \Sigma^*$.
For $u = v$ we have $u \equiv_L v$. So, now let $u \ne v$. Then $u = u'w$
and $v = v' w$ for some $u',v' \in C^*$ and $w \in \pref{V} \setminus V$.
We have $u' \equiv_{C^*} \varepsilon$ and $v' \equiv_{C^*} \varepsilon$.
With \cref{prelim:eqn:Nerodecongruence} this yields $u \equiv_{C^*} w \equiv_{C^*} v$.

Property~\ref{constrsync:rel:saturateU}: By \cref{constrsync:prop:prefix-code-pre-equivalence},
for $u \in U$, we have $[u]_{\preeqLL{C}{V}} = \{u\}$. Thus~$\preeqLL{C}{V}$ saturates~$U$.

Property~\ref{constrsync:rel:right-congr2}: Suppose $u \equiv_{C^*} v$ for $u,v \in U$ and let $x\in \Sigma$.
We have $ux, vx \in U$ or, without loss of generality, $ux \notin U$.
Then $ux \in C^*w$ for some $w \in \pref{V} \setminus V$.
Assume $w \ne \varepsilon$. Then we can write $w = w'x$  for some $w' \in \pref{V} \setminus V$.
This gives $u \in C^*w'$, which is not possible since $u \in U$.
Thus, we must have $w = \varepsilon$ and so $ux \in C^*$.
By \cref{prelim:eqn:Nerodecongruence}, we find $ux \equiv_{C^*} vx$.
This then gives $vx \in C^*$ and so $ux \preeqLL{C}{V} vx$.

Property~\ref{constrsync:rel:right-congr1}: 
Suppose $u \preeqLL{C}{V} v$ for $u,v \in \Sigma^*$ and let $x \in \Sigma$.
We can assume $u \ne v$.
Then we have $u,v \in C^*w$
for some $w \in \pref{V} \setminus V$.
The claim holds for $ux, vx \in U$.
Next, without loss of generality, let $ux \notin U$. Then
$ux \in C^*w'$ for some $w' \in \pref{V} \setminus V$.
Suppose $w' = \varepsilon$.
By Property~\ref{constrsync:rel:refinesNerode} we have $u \equiv_{C^*} v$.
This gives $ux \equiv_{C^*} vx$.
With $ux \in C^*$ we have $ux \equiv_{C^*} \varepsilon$
and so $vx \equiv_{C^*} \varepsilon$ and thus $vx \in C^*$.
This then gives $ux \preeqLL{C}{V} vx$.
Next, suppose $w' \ne \varepsilon$.
We have $ux \in C^*w' \cap C^* wx$. 
Since~$C$ is a prefix code, we have $w', wx \notin C\Sigma^+$ (note that the possibility $wx \in C$
has not been excluded yet).
By \cref{constrsync:prop:char-prefix-code,constrsync:prop:suffix-factorizations}
and as $|w'|, |wx| > 0$ we then get $w' = wx$.
This gives $ux, vx \in C^*w'$ and so $ux \preeqLL{C}{V} vx$.
\end{proof} 

Thus, by \cref{constrsync:prop:properties-simCV}, we can define an equivalence
relation in the following way.

\begin{definition} \label{def:eqLL}
	Let $C \subseteq \Sigma^+$ be a prefix code and $V \subseteq C$.
	The relation $\eqLL{C}{V}$ is defined
	to be $\eqLsimU{C^*}{\preeqLL{C}{V}}{U}$ with $U = \Sigma^* \setminus ( C^* ( \pref{V} \setminus V) )$.
\end{definition}

This relation fulfills
 \cref{constrsync:thm:relation-right-congruence,constrsync:prop:equiv-classes-form,constrsync:prop:relation-refines}.
Unraveling the definition, for $u,v \in \Sigma^*$ we have the following:
\begin{align*}
u \eqLL{C}{V} v
& \Leftrightarrow
( u = v ) \lor ( \exists w \in \pref{V} \setminus V \colon u,v \in C^*w ) \lor
( u \equiv_{C^*} v \land u,v \notin C^* (\pref{V} \setminus V) ) \\
& \Leftrightarrow
( \exists w \in \pref{V} \setminus V \colon u,v \in C^*w ) \lor
( u \equiv_{C^*} v \land u,v \notin C^* (\pref{V} \setminus V) )
\end{align*}
where $u = v$ can be omitted since each word is either in $C^*(\pref{V}\setminus V)$
or its complement. Then with \cref{constrsync:prop:equiv-classes-form,constrsync:prop:prefix-code-pre-equivalence}
we have
\begin{equation}\label{constrsync:eqn:eqLLCV}
u \eqLL{C}{V} v \Leftrightarrow
\begin{cases}
	v \in C^*w & \text{if } u \in C^*w \text{ with } w \in \pref{V} \setminus V, \\
	v \in [u]_{\equiv_{C^*}} \setminus ( C^* ( \pref{V} \setminus V ) ) & \text{if } u \notin C^* ( \pref{V} \setminus V ).
\end{cases}
\end{equation}
With \cref{constrsync:prop:Cstarw-relation-unique} we also obtain that for $u,v \in \pref{V}\setminus V$
we have 
\begin{equation} \label{constrsync:eqn:eqLLCV-unique-on-prefV}
	u \eqLL{C}{V} v \Leftrightarrow u = v.
\end{equation} 


\begin{example}\label{constrsync:ex:eqLL}
For $C = aa^*b$ and $V = \{ab,aab\}$
we have the distinct equivalence classes $[\varepsilon]_{\eqLL{C}{V}} = C^*$,
$[a]_{\eqLL{C}{V}} = C^*a$, $[aa]_{\eqLL{C}{V}} = C^*aa$, $[aaa]_{\eqLL{C}{V}} = C^*aaaa^*$,
and $[b]_{\eqLL{C}{V}} = \Sigma^* \setminus \pref{C^*}$.
In fact,
we can construct an automaton
$\automat{A_{\eqLL{C}{V}}}{\Sigma^* / {\eqLL{C}{V}}}{\Sigma}{\delta_{\eqLL{C}{V}}}{[\varepsilon]_{\eqLL{C}{V}}}{\{ [\varepsilon]_{\eqLL{C}{V}} \}}$
with
transition function given by $\delta_{\eqLL{C}{V}}([w]_{\eqLL{C}{V}}, x) = [wx]_{\eqLL{C}{V}}$ for $w \in \Sigma^*$ and $x \in \Sigma$, which
is well-defined by \cref{constrsync:thm:relation-right-congruence}.
This automaton accepts $C^*$
since $[\varepsilon]_{\eqLL{C}{V}} = C^*$. See~\cref{fig:eqLLaut}
for a depiction of $\mathcal A_{\eqLL{C}{V}}$
with the minimal automaton~$\minaut{C^*}$ for comparison.

\begin{figure}
	\centering
	\begin{subfigure}[b]{0.55\textwidth}
		\centering			
		\begin{tikzpicture}[shorten >=1pt,node distance=2.5cm and 3cm,on grid,auto]
			\tikzstyle{state}=[]
			
			\node[state] (eps)                  {$[\varepsilon]_{\eqLL{C}{V}}$};
			\node[state] (a)   [right = of eps] {$[a]_{\eqLL{C}{V}}$};
			\node[state] (aa)  [right = of a]   {$[aa]_{\eqLL{C}{V}}$};
			\node[state] (aaa) [below = of aa]  {$[aaa]_{\eqLL{C}{V}}$};
			\node[state] (b)  [below = of eps] {$[b]_{\eqLL{C}{V}}$};
			
			\path[->] (eps) edge node {$a$} (a)
			(a)	edge node {$a$} (aa)
			(aa)   edge node {$a$} (aaa)
			(aaa) edge [loop right, looseness=2.75] node {$a$} (aaa);
			
			
			\path[->] (a)   edge [bend left=20] node {$b$} (eps)
			(aa)  edge [bend left=50, above] node {$b$} (eps)
			(aaa) edge [bend left=30, above, pos=.3] node {$b$} (eps);
			
			\path[->] (b)  edge [loop left, looseness=2.75] node {$a,b$} (b);	
			\path[->] (eps) edge [left] node {$b$} (b);
		\end{tikzpicture}
		\caption{The automaton $\mathcal A_{\eqLL{C}{V}}$}
		\label{fig:AeqLL}
	\end{subfigure}
	\hfill
	\begin{subfigure}[b]{0.36\textwidth}
		\centering			
		\begin{tikzpicture}[shorten >=1pt,node distance=2.5cm,on grid,auto]
			\tikzstyle{state}=[]
			
			\node[state] (eps)                  {$[\varepsilon]_{\equiv_{C^*}}$};
			\node[state] (a)   [right = of eps] {$[a]_{\equiv_{C^*}}$};
			\node[state] (b)  [below = of eps] {$[b]_{\equiv_{C^*}}$};
			
			\path[->] (eps) edge node {$a$} (a)
			(a)   edge [loop right, looseness=2.75] node {$a$} (a);
			\path[->] (b)  edge [loop right, looseness=2.75] node {$a,b$} (b);			
			\path[->] (eps) edge [left] node {$b$} (b)
			(a)   edge [bend left] node {$b$} (eps);
		\end{tikzpicture}
		\caption{The automaton $\minaut{C^*}$}
		\label{fig:ACstar}
	\end{subfigure}
	\caption[Example automaton derived from $\eqLL{C}{V}$]{\Cref{fig:AeqLL} shows the automaton~$\mathcal A_{\eqLL{C}{V}}$ 
		from \cref{constrsync:ex:eqLL} for $C = aa^*b$ and $V = \{ab,aab\}$;
		\cref{fig:ACstar} shows $\minaut{C^*}$.}
	\label{fig:eqLLaut}
\end{figure}
\end{example}

\begin{example}
	Consider $C = \{aba,bba\}$ and $V = \{aba\}$.
	Then we have $ab \NeqLL{C}{V} bb$,
	but $ab \equiv_{C^*} bb$.   
\end{example}


\begin{remark} % gibt das überhaupt Rechtskongruenz?
Setting two words $u,v \in \Sigma^*$ in relation precisely if
\[
( \exists w \in \pref{V} \setminus V \colon u,v \in C^*w ) \lor
( u \equiv_{C^*} v \land \forall w \in \pref{V} \setminus V \colon |\{u,v\} \cap C^*w| \le 1 )
\]	
For the code 
and language~$V$ from \cref{constrsync:ex:eqLL}  
this relation would give the same classes.
However, for example, 
for $C = \{ ac, bc\}$
and $V = \{ac\}$ we have that $a$ and $b$ are in relation
whereas $a \NeqLL{C}{V} b$.
\end{remark}

In our constructions
we need this relation to have finitely many classes, a property characterized
in the next statement.

\begin{proposition} \label{constrsync:prop:eqCV-finite-index}
There are finitely many equivalence classes for $\eqLL{C}{V}$ if and only if $C^*$ is regular
and~$V$ is finite. 
\end{proposition}
\begin{proof}
	First, suppose~$C^*$ is regular and~$V$ is finite.
	By \cref{prelim:min_aut:terminal_object} 
	there are only finitely many equivalence classes for $\equiv_{C^*}$.
	By \cref{constrsync:eqn:eqLLCV}
	each equivalence class has the form $C^*w$
	for some $w \in \pref{V}\setminus V$
	or $[u]_{\equiv_{C^*}} \cap U$.
	Since $\pref{V}\setminus V$ is finite
	and there are only finitely many $\equiv_{C^*}$-equivalence
	classes, we obtain that there are only
	finitely many $\eqLL{C}{V}$-equivalence classes.
	
	Conversely, suppose there are finitely many equivalence classes for $\eqLL{C}{V}$.
	By \cref{constrsync:cor:relation-refines_Nerode}, 
	$\eqLL{C}{V}$ refines $\equiv_{C^*}$.
	Thus the $\equiv_{C^*}$-equivalence classes
	are unions of $\eqLL{C}{V}$-equivalence classes
	and we find that there are only finitely
	many $\equiv_{C^*}$-equivalence classes.
	So $C^*$ is regular by \cref{prelim:min_aut:terminal_object}.
	By \cref{constrsync:prop:equiv-classes-form,constrsync:prop:Cstarw-relation-unique,constrsync:prop:prefix-code-pre-equivalence},
	distinct words in $\pref{V}\setminus V$ belong to distinct $\eqLL{C}{V}$-equivalence classes.
	Thus $\pref{V}\setminus V$ is finite.
	Since~$V$ is a prefix code, this set equals the set of proper prefixes of words from~$V$.
	Now, finiteness of its set of proper prefixes yields that~$V$ is finite.
\end{proof}

\subsection{An \NL-Hardness Condition}
\label{constr_sync:sec:NL-hardness}

\begin{theorem} \label{constrsync:thm:NL-hardness}
	Let $C \subseteq \Sigma^+$ be a prefix code that is a regular language
	and $u_1, u_2 \in C$ be distinct.
	Suppose there exists $u \in C^*$ such that
	\[
	vu \in ( (\pref{C}\setminus C)\Sigma \setminus \pref{C} \cup C )^*	
	\]
	for all $v \in \pref{C} \setminus C$.
	Then $\constrsync{C^*}$ is \NL-hard.
\end{theorem} 
\begin{proof}
	We give a reduction from \DFAS
	for automata over a binary alphabet, which
	is \NL-hard by \cref{prelim:complexity:general_aut_sync}.
	Let $\semiautomat{A}{Q}{\{a_1, a_2\}}{\delta}$
	be an input automaton over the binary alphabet $\{a_1, a_2\}$.
	Let $v_1, v_2 \in \Sigma^*$ such that $u_1 \in v_1\Sigma$
	and $u_2 \in v_2\Sigma$.
	Set $D = (\pref{C}\setminus C)\Sigma \setminus \pref{C} \cup C$.
	For short, denote $[v]_{\eqLpp{D}{\{u_1, u_2\}}}$ 
	simply by $[v]$ for $v \in \Sigma^*$ in the following.	
	Let $\semiautomat{A'}{Q'}{\Sigma}{\delta}$
	with $Q' = Q \times \{ [v] \mid v \in \Sigma^* \}$ and
	where, for $q \in Q$, $v \in \Sigma^*$, and $x \Sigma$, we set
	\[
	\delta'((q, [v]), x)
	=
	\begin{cases}
	(\delta(q, a_i), [vx]) & \text{if } [v] = [v_i] \text{ and } v_i x = u_i, \\ 
	(q, [vx]) & \text{otherwise.}
	\end{cases}
	\]
	Let $\varphi \colon D^* \to \{a_1, a_2\}^*$
	be given by $\varphi(u_i) = a_i$ for $i \in \{1,2\}$
	and $\varphi(v) = \varepsilon$ for $v \in D \setminus \{u_1, u_2\}$
	with $\varphi(w_1 \dots w_m) = \varphi(w_1) \cdots \varphi(w_m)$,
	which is well-defined since~$D$ is a code.
	Inductively, with $[\varepsilon] = [w]$ if and only if $w \in D^*$,
	we obtain that for $w_1, \ldots, w_m \in D$
	we have
	\begin{equation}\label{constrsync:eqn:NL-hardness}
		\delta'((q, [\varepsilon], w_1 \cdots w_m) = (\delta(q, \varphi(w_1 \cdots w_m)), [\varepsilon]).
	\end{equation}	
	
	Suppose for the word $w_1 \cdots w_m$ with $w_1, \ldots, w_m \in C$
	we have $|\delta'(Q \times \{ [\varepsilon] \}, w_1 \cdots w_m)| = 1$.
	By \cref{constrsync:eqn:NL-hardness}
	we find 
	$\delta'(Q \times \{ [\varepsilon] \}, w_1 \cdots w_m) = \delta(Q, \varphi(w_1 \cdots w_m)) \times \{ [\varepsilon]) \}$.
	Thus we have $| \delta(Q, \varphi(w_1 \cdots w_m)) | = 1$.
	
	Conversely, suppose the word $x_1 \cdots x_m$ with $x_1,\ldots,x_m \in \{a_1, a_2\}$
	synchronizes $\mathcal A$.
	For each $i \in \{1,\ldots,m\}$
	let $w_i = u_1$ if $x_i = a_1$ and $w_i = u_2$ if $x_i = a_2$.
	Then for $q \in Q$
	we have
	\[
	\delta'((q, [\varepsilon]), w_1 \cdots w_m) = (\delta(q, x_1 \cdots x_m), [\varepsilon]).
	\]
	By assumption, for $v \in \pref{C} \setminus C$, we have $vu \in D^*$
	and thus
	\[
	\delta'((q, [v]), u) = \delta'((q, [\varepsilon]), vu) \in Q \times \{[\varepsilon]\}.
	\]
	Since $[\varepsilon] = [w]$ if and only if $w \in D^*$
	and with $\pref{C}\setminus C = \pref{D} \setminus D$
	this also gives $\{[v] \mid v \in \Sigma^* \} = \{[v] \mid v \in \pref{C}\setminus C\}$,
	we find $\delta'(Q', vu) \subseteq Q \times \{[\varepsilon]\}$.
	Thus, we find $\delta'(Q', vu w_1 \cdots w_m) = \delta(Q, x_1 \cdots x_m) \times \{[\varepsilon]\}$,
	and this set is a singleton set.
\end{proof}

\begin{example} \label{constrsync:ex:finite-prefix-code-with-letter}
	Let Let $C \subseteq \Sigma^+$ be a finite prefix code with $a \in C$
	for some $a \in \Sigma$.
	Set $n = \max\{ i : a^i \in \fact{L} \} + 1$.
	Then for $u = a^{n+1}$ we have
	\[
	vu \in ( (\pref{C}\setminus C)\Sigma \setminus \pref{C} \cup C )^*	
	\]
	for all $v \in \pref{C} \setminus C$.
\end{example}

\begin{corollary}
	Let $C \subseteq \Sigma^+$ be a  finite prefix code with $a \in C$
	for some $a \in \Sigma$.
	Then $\constrsync{C^*}$ is \NL-hard.
\end{corollary}
\begin{proof}
	The claim follows from
	\cref{constrsync:thm:NL-hardness,constrsync:ex:finite-prefix-code-with-letter}.
\end{proof}


\subsection{Two \PSPACE-Hardness Conditions}
\label{constr_sync:sec:PSPACE-hardness}

Here, we present two \PSPACE-hardness conditions. 
\Cref{constr_sync:thm:uC} is concerned with constraint languages
of the form $UuV$, $U \subseteq u^*$ non-empty, 
and \cref{constr_sync:thm:Cu} with constraint
languages of the form $C^*uL$
for some prefix code~$C$, word $u$
and $L \subseteq \Sigma^*$, where $U, V, L$, and $C$ are regular languages.
The proof of \cref{constr_sync:thm:Cu} uses \cref{constrsync:lem:aux-state-new-code} 
from the present section.
In \cref{constrsync:prop:PSPACE-hard-cases}, utilizing
\Cref{constr_sync:thm:uC,constr_sync:thm:Cu} in its proof,
we give a list of binary constraint languages giving \PSPACE-hard
problems which will be used later in \cref{sec:constr_sync:three_states}.

\begin{theorem} \label{constr_sync:thm:uC}
	Let $C \subseteq \Sigma^+$ be a regular language that is a prefix code
	with $u_1, u_2 \in C^+$ where $\{u_1, u_2\}$ is a code
	and $U,V \subseteq \Sigma^*$ be non-empty and regular such that
	\[
	\{ u_1, u_2 \}^* \subseteq V \subseteq \pref{C^*}
	\]
	and $U \subseteq u^*$ for some $u \in \Sigma^+$.
	%that is strongly auxiliary state avoiding w.r.t.~$C$. % braucheich gar nicht mehr, geht jetzt auch mit weakly
	Suppose $vu$ is weakly auxiliary state avoiding 
	w.r.t.~$C$ for each $v \in \pref{C}\setminus C$.
	Then $\constrsync{UuV}$ is \PSPACE-hard.
\end{theorem}
\begin{proof}
	We give a reduction from the problem \SFS over a binary alphabet
	which is \PSPACE-complete 
	by~\cref{prelim:thm:binary_pspace_problems}.
	Let $(\mathcal A, S)$ with $\semiautomat{A}{Q}{\{a_1, a_2\}}{\delta}$ and $S \subseteq Q$ 
	be an \SFS-instance.
	We can assume $S$ is non-empty.	
	
	Set $n = r\cdot s \cdot |u|$
	for $r, s \ge 1$
	such that $u_1 \in C^r$
	and $u_2 \in C^s$.
	Let $\overline{u}_1 = u_1^{s|u|}$
	and $\overline{u}_2 = u_2^{r|u|}$.
	We have $\overline{u}_1, \overline{u}_2 \in C^n$
	and since $\{u_1, u_2\}$
	is a code, we have $\overline{u}_1 \ne \overline{u}_2$
	by \cref{constr_sync:lem:two-words-code}.
	For readability, in this proof, for $v \in \Sigma^*$,
	we denote $[v]_{\eqLL{C^n}{\{ \overline{u}_1, \overline{u}_2 \}}}$
	simply by $[v]$ (confusion with the notation for transition matrices
	should not arise as these matrices are not used in the present proof).
	Since~$C$ is regular, by \cref{constrsync:prop:eqCV-finite-index},
	$\{ [v] \mid v \in \Sigma^* \}$ is finite. 
	Note that $\pref{C^*} = \pref{(C^n)^*}$,
	and that $[v] = [\varepsilon]$ if and only if $v \in (C^n)^*$
	and so $\{ [v] \mid v \in \pref{C^*} \} = \{ [v] \mid v \in \pref{C^n} \}$.
		
	
	\medskip
	\noindent\underline{Construction:} Pick some $s \in S$ and let $\overline{S} \colon Q \to S$
	be given by $\overline{S}(q) = s$ for $q \in Q \setminus S$
	and $\overline{S}(q) = q$ for $q \in S$.	
	Then set $\semiautomat{A'}{Q'}{\Sigma}{\delta}$
	where $Q' = Q \times \{ [v] \mid v \in \pref{C^*} \}$ and, 
	for $v \in \pref{C^*}$ and $x \in \Sigma$, we set
	\[
	\delta'((q, [v]), x)
	=
	\begin{cases}
		(\delta(q, a_i), [\varepsilon]) & \text{if } vx \in (C^n)^* \overline{u}_i \text{ for some } i \in \{1,2\}, \\
		(q, [vx]) & \text{if } vx \in \pref{C^*} \text{ and } vx \notin (C^n)^*\overline{u}_1 \cup (C^n)^*\overline{u}_2, \\
		(\overline{S}(q), [\varepsilon]) & \text{otherwise.}		
	\end{cases}
	\]
	Consider $v, v' \in \pref{C^*}$
	with $[v] = [v']$ and $x \in \Sigma$.
	Suppose $vx \in (C^n)^* \overline{u}_i$ for some $i \in \{1,2\}$.
	Then we have $v \in (C^n)^* w$
	for $w \in \pref{\overline{u}_i}\setminus \{u_i\}$
	with $|w| + 1 = |\overline{u}_i|$.
	By \cref{constrsync:eqn:eqLLCV} 
	we find $v' \in (C^n)^*w$.
	Thus $v'x \in (C^n)^* \overline{u}_i$.
	We have $[vx] = [v'x]$ by \cref{constrsync:thm:relation-right-congruence}.
	Thus, by \cref{constrsync:cor:both-prefix-L}, either
	both words $vx, v'x$ are contained in $\pref{C^*}$
	or both are not contained in this set.
	Combining these facts,	
	we obtain that if one of the defining cases of~$\delta'$
	applies to~$v$ it also applies to~$v'$ and vice versa,
	and so the function $\delta'$ is well-defined.
	
	
	Let $\varphi \colon (C^n)^* \to \{a_1, a_2\}^*$ be given
	by $\applyfunc{w_1 \cdots w_m}{\varphi} = \applyfunc{w_1}{\varphi} \cdots \applyfunc{w_m}{\varphi}$ 
	for $w_1, \ldots, w_m \in C^n$
	where $\varphi(\overline{u}_i) = a_i$ for $i \in \{1,2\}$
	and $\varphi(v) = \varepsilon$ for $v \in C^n \setminus \{\overline{u}_1, \overline{u}_2\}$.
%	\[
%	\applyfunc{w_i}{\varphi} = 
%	\left\{ \begin{array}{ll}
%		a_1 & \text{if } w_i =  \overline{u}_1, \\
%		a_2 & \text{if } w_i \ne \overline{u}_1,
%	\end{array} 
%	\right.
%	\]
	Since the factorization of $w$ into words from~$C^n$ is unique, $\varphi$
	is well-defined. Note that $\varphi$ is surjective and for each $w \in C^*$
	there exists $\hat{w} \in \{\overline{u}_1, \overline{u}_2\}^*$
	such that $\applyfunc{w}{\varphi} = \applyfunc{\hat{w}}{\varphi}$.
				
	\medskip
	\noindent\underline{Properties:} By an easy induction, for $w \in (C^n)^*$, we find that we have
	\begin{equation} \label{constrsync:eqn:uCvarphi}
		\delta'((q, [\varepsilon]), w) = (\delta(q, \varphi(w)), [\varepsilon]).
	\end{equation}
	Next, consider a word $v \in \pref{C^*} \Sigma \setminus \pref{C^*}$.
	We can write $v = v_1 \cdots v_{m-1} v_m x$ for $m \ge 1$
	with $v_1,\ldots, v_{m-1} \in C^n$, $v_m \in \pref{C^n}\setminus C^n$,
	and $x \in \Sigma$
	where we set $v_1 \cdots v_{0}$ 
	equal to $\varepsilon$.
	Then, by \cref{constrsync:eqn:uCvarphi} and the definition of~$\delta'$,
	for $q \in Q$ with $p = \delta(q, \varphi(v_1 \cdots v_{m-1}))$ 
	we have
	\begin{equation}\label{constrsync:eqn:uCweakly-in-S}
	\delta'((q, [\varepsilon]), v) 
	= \delta'((p, [\varepsilon]), v_m x) 
	= \delta'((p, [v_m]), x) 
	= (\overline{S}(p), [\varepsilon]).
	\end{equation}
	Let $v \in \pref{C^*}$. 
	We have $\pref{C^*} = C^* (\pref{C}\setminus C)$.
	By assumption, each word in $(\pref{C}\setminus C )u$
	is weakly auxiliary state avoiding w.r.t.~$C$.
	By \cref{constrsync:cor:Cn-weakly-aux-state} 
	we obtain that each word in
	$C^* (\pref{C}\setminus Cu$ is weakly auxiliary state avoiding w.r.t.~$C$,
	i.e.,  a concatenation of words from $\pref{C^*}\Sigma\setminus \pref{C^*}$.
	By induction, with \cref{constrsync:eqn:uCweakly-in-S} we then obtain 
	\begin{equation} \label{constrsync:eqn:uCinS}
		\delta'((q, [v]), u) = \delta'((q, [\varepsilon]), vu) \in S \times \{[\varepsilon]\}.
	\end{equation}
	More precisely, since $|u| \le n$ (and so it does not contain $\overline{u}_1$ or $\overline{u}_2$
	as a factor), we also get\footnote{Note that $|u| \le n$ implies, since $u$ is also weakly auxiliary state avoiding w.r.t.~$C^n$
		by \cref{constrsync:prop:auxLn}, that~$u$ is strongly auxiliary state avoiding w.r.t.~$C^n$
		by \crefitem{constrsync:prop:weakly_suffix_strongly}{item:suffix2}.}
	\begin{equation} \label{constrsync:eqn:uContoS}
		\delta'((q, [\varepsilon]), u)  = (\overline{S}(q), [\varepsilon]).
	\end{equation}
	This then gives $\delta'(Q \times [\varepsilon], u^i) = S \times \{[\varepsilon]\}$
	for all $i \ge 1$ 
	since $\overline{S}^j = \overline{S}$ for all $j \ge 1$.
	Thus, with \cref{constrsync:eqn:uCinS},
	we find $\delta'(Q', u^i) = S \times \{[\varepsilon]\}$ for all $i \ge 1$.
			
	\medskip
	\noindent\underline{Verification:} Suppose $u^i v$ with $i \ge 1$ where $u^i \in U$ and $v \in V$
	synchronizes $\mathcal A'$. By \cref{constrsync:eqn:uCinS,constrsync:eqn:uContoS} 
	we have $\delta'(Q', u^i) = S \times \{\varepsilon\}$.
	Since $v \in \pref{C^*}$, we can write $v = v_1 \cdots v_{m-1} v_m$
	with $v_1, \ldots, v_{m-1} \in C^n$ and $v_m \in \pref{C^n}\setminus C^n$.
	Then, for each $q \in S$, we have
	\[
	\delta'((q, [\varepsilon]), v_1 \cdots v_{m-1}v_m ) 
	= (\delta(q, \varphi(v_1\cdots v_{m-1})), [v_m]).
	\]
	By assumption, this equals the same state for each $q \in S$
	and so $|\delta(S, \varphi(v_1\cdots v_{m-1}))| = 1$.
		
	Conversely, let $v \in \{a_1, a_2\}^*$
	be such that $|\delta(S, v)| = 1$.
	Pick a word $u^i \in Uu$ with $i \ge 1$.
	Write $v = v_1 \cdots v_m$ with $v_1, \ldots, v_m \in \{a_1, a_2\}$
	and let $w_j$ equal $\overline{u}_1$ if $v_j = a_1$
	and $\overline{u}_2$ if $v_j = a_2$ for each $j \in \{1,\ldots,m\}$.
	We have $v = \varphi(w_1 \cdots w_m)$ and so, from \cref{constrsync:eqn:uCvarphi},
	for $q \in S$, we get
	\[
	\delta'((q, [\varepsilon]), w_1 \cdots w_m) = (\delta(q, v), [\varepsilon]),
	\]
	where $\delta(q, v)$ equals the same state for each $q \in S$ by assumption.
	Thus, since $\delta'(Q', u^i) = S \times \{\varepsilon\}$
	by \cref{constrsync:eqn:uCinS,constrsync:eqn:uContoS},
	the word $u^i w_1 \cdots w_m$ synchronizes $\mathcal A'$.
	Since $w_1 \cdots w_m \in \{u_1, u_2\}^* \subseteq V$, this word
	is a synchronizing word for~$\mathcal A'$
	from $UuV$.
\end{proof}

\begin{remark}
	By \crefitem{constrsync:ex:refineNerode}{constrsync:ex:refineNerode:right-ideal},
	we obtain the relation
	$\eqLsimU{C^n}{=}{\Sigma^* \setminus \pref{\{\overline{u}_1, \overline{u}_2\}}}$. 
	Let us also denote $[w]_{\eqLsimU{C^n}{=}{\Sigma^* \setminus \pref{\{\overline{u}_1, \overline{u}_2\}}}}$
	by $[w]$ for $w \in \Sigma^*$ in the following.
	This relation can also be used to give an alternative construction
	in the proof of \cref{constr_sync:thm:uC}
	by setting $\semiautomat{A'}{Q \times \{ [w] \mid w \in \pref{C^n} \setminus C^n\}}{\Sigma}{\delta'}$
	and, for $q \in Q$, $w \in \pref{C^n}\setminus C^n$ and $x \in \Sigma$, 
	\[
	\delta'((q, [w], x)
	= \begin{cases}
		(q, [\varepsilon]) &\text{if } wx \in C^n \setminus \{\overline{u}_1, \overline{u}_2\}, \\
		(\delta(q, a_i), [\varepsilon]) & \text{if } wx = \overline{u}_i \text{ for some } i \in \{1,2\}, \\
		(q, [wx]) & \text{if } wx \in \pref{C^n} \setminus C^n, \\
		(\overline{S}(q), [\varepsilon]) & \text{if } wx \notin \pref{C}.		 		
	\end{cases}
	\]
	Then the remaining argument can be continued in a similar way as the original proof.	
	Yet another similar proof can be given
	by using the classes 
	$[w]_{\eqLpp{C}{|\overline{u}_1|}}$
	for $w \in \Sigma^*$, in the following simply denoted by~$[w]$.
	Then 
	let $\semiautomat{A'}{Q'}{\Sigma}{\delta'}$
	with $Q' = Q \times \{ [w]  : w \in \pref{C^n} \setminus C^n \}$
	and, for $q \in Q$, $w \in \pref{C^n}\setminus C^n$ and $x \in \Sigma$, set
	\[
	\applytrans{\delta'}{(q, [w])}{x} = 
	\begin{cases}
		(\applytrans{\delta}{q}{a_1}, [\varepsilon]) & \text{if } wx = \overline{u}_1, \\
		(\applytrans{\delta}{q}{a_2}, [\varepsilon]) & \text{if } wx \in C^n \setminus \{ \overline{u}_1 \}, \\
		(q, [wx])                                    & \text{if } wx \in \pref{C^n} \setminus C^n \\
		(\overline{S}(q), [\varepsilon]) & \text{if } wx \notin \pref{C^n}.	
	\end{cases}
	\]
\end{remark}

\begin{lemma} \label{constrsync:lem:aux-state-new-code}
	Let $C \subseteq \Sigma^+$ be a prefix code
	and $w \in C^+$
	be such that we have
	\[
	( \pref{C} \setminus C) w \subseteq ( (\pref{C}\setminus C)\Sigma \setminus \pref{C} \cup C)^*.
	\]
	Then $D = ( C \setminus \{u\} ) \cup u(C \setminus \{u\}) \cup uuC$
	for $u \in C$ is a prefix code 
	such that for $\hat{u} \in C \setminus \{u\}$ we have
	\[
	( \pref{D} \setminus D ) w \hat{u} \in ( (\pref{D\setminus D}\Sigma \setminus \pref{D} \cup D)^*.
	\]
\end{lemma}
\begin{proof} 
	Set $U = (\pref{C}\setminus C)\Sigma \setminus \pref{C}$
	and $V = (\pref{D}\setminus D)\Sigma \setminus \pref{D}$.
	The following equations are easy to verify:
	\begin{align}
		\pref{D} \setminus D & = \{\varepsilon, u, uu\} ( \pref{C} \setminus C )  \label{constrsync:eqn:prefD} \\
		\{\varepsilon, u, uu\}(\{\varepsilon\} \cup C \cup C^2 \cup C^3)& = (\{\varepsilon\} \cup D)\{\varepsilon, u, uu\} \label{constrsync:eqn:CDcommute} \\
		V & = \{\varepsilon,u,uu\} U \label{constrsync:eqn:VuuuU}
	\end{align}
	Note that, from~\cref{constrsync:eqn:CDcommute}, 
	by induction, we obtain
	$
		\{\varepsilon, u, uu\}C^* = D^*\{\varepsilon, u, uu\}.
	$
	Since $C^* \subseteq \{\varepsilon, u, uu\} C^*$
	and $ D^*\{\varepsilon, u, uu\} \subseteq C^*$
	this gives $C^* = D^*\{\varepsilon, u, uu\}$.	
	Next, we show that
	\begin{equation} \label{constrsync:eqn:UCstar} 
	(U \cup C)^* \subseteq (V \cup D)^*\{\varepsilon, u, uu\}
	\end{equation}
	Let $w \in (U \cup C)^*$.
	For $w = \varepsilon$, this is clear.
	Otherwise, we have $w = v x$ with $v \in (U \cup C)^*$
	and $x \in U \cup C$.
	Inductively, we find $v \in (V \cup D)^* \{\varepsilon, u, uu\}$.
	By \cref{constrsync:eqn:CDcommute,constrsync:eqn:VuuuU}, we have
	\[
	\{\varepsilon, u, uu\} x \subseteq
	\begin{cases}
		V     & \text{if } x \in U, \\
		D^* \{\varepsilon, u, uu\} & \text{if } x \in C. 
	\end{cases} 
	\]
	Thus $\{\varepsilon, u, uu\} x \in (V \cup D)^* \{\varepsilon, u, uu\}$,
	which gives $w \in (V \cup D)^* \{\varepsilon, u, uu\}$.
	
	Next, since $\hat{u} \in D$ and $\{\varepsilon, u, uu\}\hat{u} \in D$,
	with \cref{constrsync:eqn:UCstar} we find 
	$(U \cup C)^* \hat{u} \subseteq (V \cup D)^* D$.
	Then, from \cref{constrsync:eqn:prefD} and the assumption
	$(\pref{C}\setminus C)w \subseteq (U \cup C)^*$, we get
	\[
	(\pref{D}\setminus D)w \subseteq \{\varepsilon,u,uu\}(\pref{C} \setminus C) w \subseteq \{\varepsilon, u, uu \} ( U \cup C )^*
	\subseteq ( U \cup C )^*
	\]
	Thus we have $(\pref{D}\setminus D)w\hat{u} \subseteq (V\cup D)^*D$
	since $(U \cup C)^* \hat{u} \subseteq (V \cup D)^* D$.
\end{proof}
	
\begin{theorem} \label{constr_sync:thm:Cu}
	Let $C \subseteq \Sigma^+$ be a 
	prefix code	that is a regular language
	and contains three distinct words $u_1, u_2, u_3 \in C$.
	Suppose we have a non-empty regular language $L \subseteq \Sigma^*$
	with $|\fact{L} \cap \{u_1, u_2, u_3\}| \le 2$
	and there exists $w \in (C \setminus \{u_3\} )^*$
	such that
	\[
	( \pref{C}\setminus C )w \subseteq ( ( ( ( \pref{C}\setminus C)\Sigma) \setminus \pref{C} ) \cup C )^*.
	\]
	Then $\constrsync{(C \setminus \{u_3\})^* u_3 L}$
	is \PSPACE-hard.
\end{theorem}
\begin{proof}
	Set $D = ( C \setminus \{u_1\} ) \cup u_1 ( C \setminus \{u_1\} ) \cup u_1 u_1 C$
	and $E = ( ( \pref{D} \setminus D )\Sigma \setminus \pref{D} ) \cup D$.
	Then~$E$ is a maximal prefix code by \cref{constrsync:prop:prefix-maximalD}
	and, by \cref{constrsync:lem:aux-state-new-code}, we have
	\begin{equation}\label{constrsync:eqn:reset-pspace-reduction}
		v w u_2 \in E^*
	\end{equation}
	for all $v \in \pref{D} \setminus D$.
	Observe that we have $\pref{D} \setminus D = \pref{E} \setminus E$.	
	
	Let $\semiautomat{A}{Q}{\{a_1, a_2, a_3\}}{\delta}$
	be an automaton.
	Set $v_1 = u_1 u_2$, $v_2 = u_1 u_1 u_2$, $v_3 = u_3$, $v_4 = u_1 u_3$, 
	and $v_5 = u_1 u_1 u_3$.
	Define morphisms $\varphi \colon E^* \to \{a_1, a_2, a_3\}$
	and $\psi \colon \{a_1, a_2, a_3\}^* \to E^*$
	by $\varphi(v_i) = a_{\min\{i,3\}}$ for $i \in \{1,2,3,4,5\}$
	and	$\varphi(v) = \varepsilon$ for $v \in E \setminus \{v_1, v_2, v_3, v_4, v_5\}$,
	and $\psi(a_i) = v_i$ for $i \in \{1,2,3\}$
	and then $\varphi(w_1 \cdots w_n) = \varphi(w_1) \cdots \varphi(w_n)$
	for $w_1,\ldots, w_n \in E$
	and $\psi(x_1 \cdots x_m) = \psi(x_1) \cdots \psi(x_m)$
	for $x_1, \ldots, x_m \in \{a_1, a_2, a_3\}$.
	These morphisms are well-defined since~$E$ and $\{a_1,a_2,a_3\}$ are codes.	
	Let $L' = \{ u \in E^* \mid uv \in L \text{ for some $v \in \Sigma^*$} \}$.
	Since $|\fact{L'} \cap \{u_1, u_2, u_3\}| \le 2$,
	at most one of the two sets
	$\fact{L'} \cap \{v_1, v_2\}$
	or
	$\fact{L'} \cap \{v_3, v_4, v_5\}$ is non-empty.
	This yields
	that at most one of the two sets
	$\fact{\varphi(L')} \cap \{a_1, a_2\}$
	or
	$\fact{\varphi(L')} \cap \{a_3\}$
	is non-empty.
	By \crefitem{constr_sync:thm:introductory_examples}{constrsync:ex:item:abcL},
	$\constrsync{\{a_1, a_2\}^*a_3 \varphi(L')}$ is \PSPACE-complete.	
	Next, we give a reduction from this problem. Consider~$\mathcal A$
	to be an input automaton for this problem.
	For simplicity, in this proof, for $v \in \Sigma^*$,
	we denote $[v]_{\eqLL{E}{\{v_1, \ldots, v_5\}}}$
	simply by $[v]$ (as in the proof of \cref{constr_sync:thm:uC}, confusion
	with the notation for transition matrices should not arise).
	By \cref{constrsync:prop:eqCV-finite-index},
	we have finitely many equivalence classes w.r.t.~$\eqLL{E}{\{v_1, \ldots, v_5\}}$.
		
	\medskip	
	\noindent\underline{Construction:} Let $w_i \in \Sigma^*$ be such that 
	$v_i \in w_i \Sigma$ for $i \in \{1,\ldots,5\}$.
	Define $\semiautomat{A'}{Q \times \{ [v] : v \in \Sigma^* \}}{\Sigma}{\delta'}$
	where, for $q \in Q$, $v \in \Sigma^*$ and $x \in \Sigma$,
	\[
	\delta'((q, [v]),x)
	=
	\begin{cases}
	(\delta(q, a_{\min\{3,i\}}, [vx]) & \text{if } [v] = [w_i] \text{ for } i \in \{1,\ldots,5\} \text{ with } w_ix = v_i, \\
	(q, [vx])) & \text{otherwise.}
	\end{cases}
	\]
	Consider $v,v' \in \Sigma^*$ with $[v] = [v']$.
	Suppose $[w_i] = [w_j]$ for $i,j \in \{1,\ldots,5\}$.
	By \cref{constrsync:eqn:eqLLCV}
	this gives $w_i, w_j \in C^*w$ for some $w \in \pref{\{v_1, \ldots, v_5\}} \setminus \{v_1, \ldots, v_5\}$.
	By \cref{contrsync:prop:prefix-code-suffix-unique} 
	we find $w_i = w_j = w$.
	Thus $[v] = [w_i]$ for some $i \in \{1,\ldots,j\}$
	implies $[v] \ne [w_j]$ for $j \ne i$,
	and this yields that~$\delta'$ is well-defined. 
	
	Note that for $v \in \Sigma^*$ we have $[v] = [\varepsilon]$ 
	if and only if $v \in E^*$,
	which also gives $\{ [v] : v \in \pref{E} \setminus E \} 
	= \{ [v] : v \in \Sigma^* \}$ since $\pref{E^*} = \Sigma^*$
	by \cref{constrsync:prop:prefix-maximal}.\footnote{We note that $[v] = [w_i]$ and $w_i x = v_i$
		is equivalent to $vx \in E^* v_i$.
		Thus, the transition function can also be defined as in the proof
		of \cref{constr_sync:thm:uC} 
		by having a case for $vx \in E^* v_i$,
		and also setting the second component equal to $[\varepsilon]$
		in this case (but $vx \in E^* v_i$ implies $[vx] = [\varepsilon]$ anyway).
		However, we defined it as above since
		then it is easier to see that $\delta'((q, [v]), u) \in Q \times \{[vu]\}$
		for $u,v \in \Sigma^*$.}	
		
		
	\medskip
	\noindent\underline{Properties:} For $w_1,\ldots, w_m \in E$, $x_1, \ldots, x_l \in \{a_1,a_2,a_3\}$,
	and $q \in Q$, we have
	\begin{align}
		\delta'((q, [\varepsilon]), w_1 \cdots w_m) & = (\delta(q, \varphi( w_1 \cdots w_m ) ), [\varepsilon]) \label{constrsync:eqn:CuLvarphi} \\ 
		\delta'((q, [\varepsilon]), \psi(x_1 \cdots x_l)) & = (\delta(q, x_1 \cdots x_l), [\varepsilon]) \label{constrsync:eqn:CuLpsi} \\
		\varphi(\psi(x_1 \cdots x_l)) & = x_1 \cdots x_l \label{constrsync:eqn:psivarphi}
	\end{align}
		
	\medskip
	\noindent\underline{Verification:} Suppose $\mathcal A'$ has a synchronizing word
	of the form $w_1 \cdots w_m u_3 v$
	with $w_1, \ldots, w_m \in C \setminus\{u_3\}$
	and $v \in L$.
	We have $w_1 \cdots w_m \in (D \setminus \{ v_3, v_4, v_5 \})^* \{ \varepsilon, u_1, u_1 u_1\}$.
	Thus, we find
	$w_1 \cdots w_m u_3 \in (D \setminus \{ v_3, v_4, v_5 \})^* \{ v_3, v_4, v_5 \} \subseteq E^*$.
	Since~$E$ is a maximal prefix code,
	by \cref{constrsync:prop:prefix-maximal} 
	we have $v \in \pref{E^*}$
	and thus we can write $v = v' v''$ with $v' \in E^*$ and $v'' \in \pref{E} \setminus E$.
	For $q \in Q$,
	we have 
	\begin{align*}
		\delta'((q, [\varepsilon]), w_1 \cdots w_m u_3 v)
		& = \delta'( (\delta(q, \varphi( w_1 \cdots w_m u_3 v' )), [\varepsilon]), v'') \tag{Eqn.~\eqref{constrsync:eqn:CuLvarphi}} \\
		& = (\delta(q, \varphi( w_1 \cdots w_m u_3 v' )), [v'']). \tag{Def.~$\delta'$}
	\end{align*}
	Since $\delta'((q, [\varepsilon]), w_1 \cdots w_m u_3 v)$
	equals the same state for each $q \in Q$,
	also $\delta(q, \varphi( w_1 \cdots w_m u_3 v' ))$
	equals the same state for each $q \in Q$
	and so $\varphi( w_1 \cdots w_m u_3 v' )$
	synchronizes $\mathcal A$.
	We have $\varphi(w_1 \cdots w_m) \in\{a_1, a_2\}^*$,
	$\varphi(u_3) = a_3$
	and $\varphi(v') \in \varphi(L')$
	since $v' \in L'$.	
		
	Conversely, suppose the word $x_1 \cdots x_m a_3 v$ synchronizes $\mathcal A$
	where $x_1, \ldots, x_m \in \{a_1, a_2\}$
	and $v = \varphi(v')$ for some $v' \in L'$.
	For $q \in Q$ and $w' \in \pref{D} \setminus D$
	we have $w'wu_2 \in E^*$ by \cref{constrsync:eqn:reset-pspace-reduction}.
	This gives
	\[
	\delta'((q, [w']), w u_2)
	= (q, [w'wu_2]) \in Q \times \{ [\varepsilon] \}.
	\]	
	By \cref{constrsync:eqn:CuLvarphi,constrsync:eqn:psivarphi}, for $q \in Q$, we
	have
	\[
	\delta'((q, [\varepsilon]), \psi(x_1 \cdots x_m a_3)v'))
	= (\delta(q, \varphi( \psi(x_1 \cdots x_m a_3)v' ) ), [\varepsilon])
	= (\delta(q, x_1 \cdots x_m a_3 v), [\varepsilon]).
	\]
	Since on the right hand side always the same state
	results, the word $wu_2 \psi(x_1 \cdots x_m a_3)v'$
	synchronizes $\mathcal A'$.\footnote{Note that
		with $\varphi(\psi(v)) = v = \varphi(v')$ we have  
	$
	\varphi(\psi(x_1 \cdots x_m a_3 v))
	= \varphi(\psi(x_1)\cdots \psi(x_m) u_3 \psi(v))
	= \varphi(\psi(x_1 \cdots x_m a_3)v').
	$	
	Then \cref{constrsync:eqn:CuLpsi}
	yields that also $wu_2\psi(x_1 \cdots x_m a_3 v)$ 
	synchronizes~$\mathcal A'$.
	}	
	By the definition of~$L'$
	we have $v' v'' \in L$ for some $v'' \in \Sigma^*$.
	Then $wu_2 \psi(x_1 \cdots x_m a_3) v' v''$
	synchronizes $\mathcal A'$
	where $\psi(x_1, \ldots, x_m) \in ( D \setminus \{ v_3, v_4, v_5 \} )^* \subseteq ( C \setminus \{u_3\} )^*$
	and $\psi(a_3) = u_3$.	
\end{proof}

\begin{remark}
	Using $[v]_{\eqLL{D}{\{v_1, \ldots, v_5\}}}$
	instead of $[v]_{\eqLL{E}{\{v_1, \ldots, v_5\}}}$
	in the proof of \cref{constr_sync:thm:Cu} (denoting it by $[v]$
	in the following) for $v \in \Sigma^*$,
	we can also, alternatively, use the following
	construction in its proof: $\semiautomat{A'}{Q \times \{ [v] : v \in \pref{D} \setminus D \}}{\Sigma}{\delta'}$
		where, for $v \in \pref{D} \setminus D$ and $x \in \Sigma$,
		\[
		\delta'((q, [v]), x)
		=
		\begin{cases} 
			((\delta(q, a_{\min\{i,3\}}), [\varepsilon]) & \text{if } vx =v_i \text{ for some } i \in \{1,2,3,4,5\}, \\
			% nicht notwendig, da element in D zu eps äquivalent
			%	(q, [\varepsilon])               & \text{if } vx \in D \setminus \{v_1, v_2, v_3, v_4, v_5\}, \\
			(q, [vx])                        & \text{if } vx \in \pref{D} \setminus \{v_1, \ldots, v_5\}, \\
			(q, [\varepsilon])               & \text{otherwise.}
		\end{cases}
		\]		
		Let $v, v' \in \pref{D} \setminus D$ with $[v] = [v']$ and $x \in \Sigma$.		
		Then either both $v, v'$ are contained in $\pref{\{v_1,\ldots,v_5\}}$
		or none is contained in this set
		by \cref{constrsync:eqn:eqLLCV} 
		and since~$D$ is a prefix code.		
		If $v,v' \in \pref{\{v_1,\ldots,v_5\}}$,
		then $v = v'$ by \cref{constrsync:eqn:eqLLCV-unique-on-prefV} 
		and so $vx = v'x$.
		Otherwise, we have $[vx] = [v'x]$.
		By \cref{constrsync:cor:both-prefix-L}
		we have that $vx \in D$ if and only if $v'x \in D$
		and $vx \in \pref{D} \setminus D$ if and only if $v'x \in \pref{D} \setminus D$.
		Combining these relations, we find that 
		one of the defining cases of~$\delta'$ applies to $vx$
		if and only if it applies to~$v'x$
		and with $[vx] = [v'x]$ we get that~$\delta'$ is well-defined.
		Since we pick $v \in \pref{D} \setminus D$ and 
		not $v \in \pref{D^*}$ we can test for equality with one of the $v_i$
		in the first case, and do not have to check for containment in $D^*v_i$, or equivalently
		whether $[v] = [w_i]$ with $w_ix = v_i$.
		Since $[\varepsilon] = [u]$ for $u \in \Sigma^*$ if and only if $u \in D^*$,
		this construction works in the same way as the original construction
		and it can be verified similarly that it yields a reduction.		
\end{remark}

Recall that, by~\cref{constr_sync:lem:two-words-code,constr_sync:cor:code_growth_of_star},
a code contains at least two words if and only if $C^*$ has exponential growth.

\begin{example}
In the proof of~\crefitem{constr_sync:thm:introductory_examples}{constrsync:ex:item:cab}
it was directly shown 
that \constrsync{c(a+b)^*} is \PSPACE-hard.
By considering the prefix code $C = \{a,b\}$ and $u = c$,
this also follows easily 
from~\cref{constr_sync:thm:uC}, since
we have $\pref{C} \setminus C = \{\varepsilon\}$
and the word $\varepsilon\cdot u$
is weakly auxiliary state avoiding
w.r.t.~$C$ by \cref{constrsync:ex:aux}.
\end{example}

\begin{example}
	\label{constrsync:pspace-hard-example}
	Let $m,n \ge 1$.
	By \cref{constr_sync:thm:Cu} with $C = \{a^m,b,c\}$
	it follows that the problem $\constrsync{\{b,c\}^*a^m}$
	is \PSPACE-hard.
	Further, as $\{b,c\}^*a^m \cap \Sigma^*\{b,c\} = \emptyset$,
	by \cref{constr_sync:prop:Gamma_suffix} 
	we find that 
	$
	\constrsync{\{b,c\}^*a^m b^n c^n \{b,c\}^*}
	$
	is \PSPACE-hard.	
\end{example}

Next, we list a few example constraint languages that yield \PSPACE-complete problems. 
We will refer to most of them later in \cref{sec:constr_sync:three_states}.
\PSPACE-hardness for the constraint language from \crefitem{constrsync:prop:PSPACE-hard-cases}{constrsync:Cu:1} has already been established in \cref{constrsync:rem:construction-ideal-intersection} as an example; nevertheless, we include another proof using \cref{constr_sync:thm:Cu}, since we believe it gives a good illustration of its application.

\begin{proposition}
	\label{constrsync:prop:PSPACE-hard-cases}
Suppose $\Sigma = \{a,b\}$
and $\Gamma \subseteq \Sigma$ with $|\Gamma| \le 1$.
The problem $\constrsync{L}$ is \PSPACE-hard
for the following constraint languages $L$:
\begin{enumerate}
    \item \label{constr_sync:aa:1} $aa(bb^*a)^*$ and $bb(aa^*b)^*$,
    \item \label{constr_sync:aa:2} $aa(bb+ba)^*$ and $bb(aa+ab)^*$,
	\item \label{constr_sync:23:3} $aba(a+bb)^*$ and $bab(b+aa)^*$,
	\item \label{constr_sync:23:4} $bba(aa^*b)^*$ and $aab(bb^*a)^*$,	
	
	\item \label{constr_sync:23:1} $aab(b+ab)^*$ and $bba(a+ba)^*$,
	\item \label{constr_sync:23:2} $aab(bb+ab)^*$ and $bba(aa+ba)^*$,
	
    \item \label{constr_sync:23:5} $baab(aa+ab)^*$ and $baab(bb+ba)^*$,
        
	\item \label{constrsync:Cu:1} $(bb+ba)^*a\Gamma^*$,
	\item \label{constrsync:Cu:2} $(a+bb)^*ba\Gamma^*$,
	\item \label{constrsync:Cu:3} $(b+ab)^*aa\Gamma^*$.
\end{enumerate}
\end{proposition} 
\begin{proof} 
	Two languages listed side by side
	are the same up to renaming of the letters. Hence, it is sufficient
	to show the claim for one language in each case.
	
	\ref{constr_sync:aa:1}:	All words of the form $b^iaa$ 
	with $i \ge 0$ are weakly auxiliary state avoiding w.r.t.~$C = bb^*a$.
	So, $\constrsync{aaC^*}$ is \PSPACE-hard by \cref{constr_sync:thm:uC}.
		
	\ref{constr_sync:aa:2}: 
	The words $aa$ and $baa$ are weakly auxiliary state avoiding w.r.t.~$\{bb,ba\}$.
	So, the problem $\constrsync{aa\{bb,ba\}^*}$ is \PSPACE-hard by \cref{constr_sync:thm:uC}.
	    
    \ref{constr_sync:23:3}: The words $aba$ and $baba$
    are weakly auxiliary state avoiding w.r.t.~$\{a,bb\}$.
    So, the problem \constrsync{aba(a+bb)^*} is \PSPACE-hard by \cref{constr_sync:thm:uC}.

  \ref{constr_sync:23:4}: 
  The word $a^ibb$ is weakly auxiliary state
  avoiding w.r.t.~$ aa^*b$ for $i \ge 0$.
  Set $C = (aa^*b)^*$ and $V = aC$. Then $\{aab,aaab\}^* \subseteq V \subseteq C^*$.
  So, by \cref{constr_sync:thm:uC}, $\constrsync{bba( aa^*b )^*}$
  is \PSPACE-hard.
  
\ref{constr_sync:23:1}: Follows by~\ref{constr_sync:aa:1} and \cref{gen_results:cor:factors_etc_complexity} 
 as $\fact{aab(b+ab)^*} = \fact{aa(bb^*a)^*}$.\footnote{Observe that $b(b+ab)^* = L(\mathcal B_{1,\{2\}})$
 	and $(bb^*a)^* = L(\mathcal B)$
 	for $\automat{B}{\{1,2\}}{\{a,b\}}{\mu}{1}{\{1\}}$
 	with $\mu(1, b) = \mu(2, b) = 2$
 	and $\mu(2, a) = 1$. Now we have
 	$L(\mathcal B) = \{\varepsilon\} \cup L(\mathcal B_{1,\{2\}})a$
 	and $L(\mathcal B_{1,\{2\}}) = L(\mathcal B)b$.
 	This gives the stronger claims $aab(b+ab)^*a \subseteq aa(bb^*a)$
 	and $aa(bb^*a)b \subseteq aab(b+ab)^*$
 	than mere equality of the factor sets.
 	Thus, both languages
 	arise from the same constraint PDFA by merely
 	redistributing the final states inside a sink component.
 	For Item~\ref{constr_sync:23:2} a similar PDFA
 	can be found with this property.}

\ref{constr_sync:23:2}: Follows by~\ref{constr_sync:aa:2} and \cref{gen_results:cor:factors_etc_complexity}
 as $\fact{aab(bb+ab)^*} = \fact{aa(bb+ba)^*}$.

\ref{constr_sync:23:5}: 
The two words $baab$ and $abaab$ are weakly auxiliary state avoiding w.r.t~$\{aa,ab\}$.
So, $\constrsync{baab\{aa,ab\}^*}$ is \PSPACE-hard by \cref{constr_sync:thm:uC}.

\ref{constrsync:Cu:1}: Set $C = \{bb,ba,a\}$. We have $(\pref{C}\setminus C)ba \subseteq C$ with $ba \in C$
and since $|\Gamma| \le 1$ we have $|\fact{\Gamma^*} \cap \{bb,ba,a\}| \le 1$.
By \cref{constr_sync:thm:Cu}, $\constrsync{\{bb,ba\}^*a\Gamma^*}$
is \PSPACE-hard.

\ref{constrsync:Cu:2}: Set $C = \{a,bb,ba\}$. We have $(\pref{C}\setminus C)a \subseteq C$ with $a \in C$
and since $|\Gamma| \le 1$ we have $|\fact{\Gamma^*} \cap \{a,bb,ba\}| \le 1$.
By \cref{constr_sync:thm:Cu}, $\constrsync{\{a,bb\}^*ba\Gamma^*}$
is \PSPACE-hard.

\ref{constrsync:Cu:3}: Set $C = \{b,ab,aa\}$. We have $(\pref{C}\setminus C)b \subseteq C$ with $b \in C$
and since $|\Gamma| \le 1$ we have $|\fact{\Gamma^*} \cap \{a,ab,aa\}| \le 1$.
By \cref{constr_sync:thm:Cu}, $\constrsync{\{b,ab\}^*aa\Gamma^*}$
is \PSPACE-hard.
\end{proof}

\subsection{Two \NP-Hardness Conditions}
\label{constr_sync:sec:NP-hardness}

%!TEX root = ../../thesis-trier.tex

Here, we give two sufficient condition to deduce \NP-hardness
in \cref{constrsync:thm:NP-hardness,constrsync:letterbounded:lem:NPhard}.
The latter uses two lemmas: \Cref{constrsync:left-ideal-uvw}
on non-equivalent words w.r.t.~the Myhill-Nerode relation
for special left ideal languages,
\cref{constrsync:letter-bounded:lem:uabi}
on the Levenshtein-Romanov mapping for a language.

%A condition using the notions of auxiliary state avoiding words similarly as for \NL\
%and \PSPACE in the previous sections, generalizing the \NP-hardness results
%for constraint languages of the form $ab^*a$, can also be developed, but we do not give
%such a result since the construction would be very similar to the ones given before.
%
% uv^*w und sigma_{\{v\}}(u) = eps

\begin{theorem} \label{constrsync:thm:NP-hardness}
	Let $u, v \in \Sigma^*$ and $L \subseteq \Sigma^+$ be non-empty and regular
	with $|v| > 0$ such that
	$
	u \notin \factor{v^*}$ and 
	$
	v \notin \factor{L}.
	$
	Then $\constrsync{uv^+L}$ and $\constrsync{uv^*L}$ are \NP-hard.
\end{theorem}
\begin{proof}
	Since $v \notin \factor{L}$
	we have $\Sigma^* v \Sigma^* \cap uv^*L = uv^+L$
	and thus, by \cref{constr_sync:gen_results:thm:ideal_separation},
	it is sufficient to show the claim for $\constrsync{uv^+L}$.
	
	We give a reduction from \DST
	over a unary alphabet, which is \NP-complete by Theorem~\ref{compl:thm:comm_ST}.
	Let $(\mathcal A, S, T)$ be an instance
	with $\semiautomat{A}{Q}{\{b\}}{\delta}$ and $S, T \subseteq Q$
	such that $S \cap T = \emptyset$.
	Pick some $\hat s \in S$
	and let $\overline{S} \colon ( Q \setminus T ) \to S$
	be given by $\overline{S}(q) = \hat{s}$ if $q \in Q \setminus (S \cup T)$
	and $\overline{S}(q) = q$ otherwise. 	
	Let $\semiautomat{A'}{\Sigma}{ ( Q \times \pref{uv'} ) \cup \{t\}}{\delta'}$
	where, for $w \in \pref{uv'}$ and $x \in \Sigma$ we set
	\[ 
	\delta'((q, w), x)
	= \begin{cases}
		(\delta(q, b), u)    & \text{if } wx = uv, \\
		t                    & \text{if } wx \notin \pref{uv'} \text{ and } q \in T, \\
		(\overline{S}(q), u) & \text{if } \sigma_{\{u\}}(wx) = u, \\
		(q, \sigma_{\{uv'\}}(wx)) & \text{otherwise,}
	\end{cases}
	\]
	for $q \in Q$ and $\delta'(t, x) = t$. Next, we show
	that there exists $i \ge 0$ such that $\delta(q, b^i) \in T$
	for all $q \in S$ if and only if $\mathcal A'$ admits a synchronizing word in $uv^+L$.
	
	First, suppose we have $\delta(q, b^i) \in T$ for all $q \in S$
	and for some $i \ge 0$. Since $S\cap T = \emptyset$ we must have $i > 0$.
	Let $(q, w)$ with $q \in Q$ and $w \in \pref{uv'}$.
	Then we have $\delta'((q, w), u) \in S \times \{u\} \cup \{t\}$.
	This follows by the definition of the third case,
	since if we still read a prefix of $uv'$, but have encountered
	the $u$ as a suffix along this prefix, we move into a state with~$u$
	as it second component. This also works if we read~$u$
	but encounter the suffix~$u$ before having finished with reading~$u$ as input 
	by \crefitem{constr_sync:prop:sigma_L_properties}{constr_sync:sigme_L:item:concatenation}.
	Next, let $q \in S$.
	Then $\delta'((q, u), v^i) = (\delta(q, b^i), u)$.
	Consider $w' \in L$.
	Since~$v \notin \factor{w'}$,
	there is a smallest prefix $w''$ of $w'$ and $x \in \Sigma$
	such that $uw'' \in \pref{uv'}$
	and $uw''x \notin \pref{uv'}$.
	Since $\delta(q, b^i) \in T$
	we then find $\delta'( (\delta(q, b^i), u), w'x ) = t$.
	Combining everything, the word $uv^iw'$ maps
	all states to~$t$.
	
	Conversely, let $uv^iw$ with $i \ge 1$ and $w' \in L$
	be a synchronizing word for $\mathcal A'$.
	Since~$t$ is a trap state, it must be the synchronizing state.
	Consider a state $q \in S$.
	We have $\delta'((q, \varepsilon), uv^i) = (\delta(q, b^i), uv^i)$.	
	Assume $\delta(q, b^i) \notin T$.
	Write $w = w_1 \cdots w_m$ with $w_1, \ldots, w_m \in \Sigma$.
	We have $\delta'((\delta(q, b^i), u), w_1 \cdots w_i) \in ( Q \setminus T ) \times \{ \sigma_{\{uv'\}}(w_1 \cdots w_i) \}$
	for $i \in \{0, 1,\ldots,m\}$ (we set $w_1 \cdots w_0$ equal to $\varepsilon$).
	This follows inductively. For $i = 0$ the claim clearly holds.
	Next, let $i \in \{1,\ldots,m\}$.
	For $\sigma_{\{uv'\}}(w_1 \cdots w_{i-1}) = uv'$
	we have $uv'w_{i} \ne uv$ since $v \notin \factor{w}$,
	and thus only the last two cases of the transition function can apply,
	which gives the claim.
	For $\sigma_{\{uv'\}}(w_1 \cdots w_{i-1}) \ne uv'$,
	we also have $uv'w_i \ne uv$
	and also only the last two cases of the transition function can apply.
	Thus, we do not end up in the state~$t$
	and so we must have $\delta(q, b^i) \in T$
	for all $q \in S$.
\end{proof}

\begin{remark}
	\Cref{constrsync:thm:NP-hardness} 
	was given in~\cite[Prop.~8]{DBLP:conf/ictcs/Hoffman20}
	but with the additional assumption $\pref{v^*}\cap L = \emptyset$
	(and the assumption $|v| > 0$ was missing, which was a mistake).
\end{remark}

\begin{lemma} \label{constrsync:left-ideal-uvw}
	Let $I = \Sigma^* u v^* w$ for $u,v,w \in \Sigma^*$
	and $x,y \in \pref{uw}$
	with $x \ne y$.
	Then we have $x \not\equiv_I y$.
\end{lemma}
\begin{proof}
	Without loss of generality, let $|x| < |y|$.
	There exists $z \in \suff{u}$
	such that $yz = uw$.
	Thus we have $yz \in I$.
	Since $|xz| < |uw|$ and~$uw$ is the shortest word in~$I$,
	we have $xz \notin I$.
	This gives $x \not\equiv_I y$.		
\end{proof}

\begin{lemma} \label{constrsync:letter-bounded:lem:uabi}
	Suppose that $a, b, c \in \Sigma$ with $a \ne b$ and $c \ne b$ (while $a = c$ is allowed). 
	Let $u,v \in\Sigma^*$, $i > |u| + |v|$, and $p \ge 0$.	
	Set $L = uab^i(b^p)^*cv$.
	Then for $w \in \Sigma^*$ we have $\sigma_L(wuab^i) = uab^i$.
\end{lemma}
\begin{proof}
	Since $uab^i \in \pref{L}$ for $w \in \Sigma^*$
	we have $\sigma_L(w uab^i) = y uab^i$ for some $y \in \Sigma^*$
	such that $y uab^i \in \pref{L}$.
	The factor $ab^i$ appears in a word from~$L$ precisely once
	after the prefix~$u$ as $i > |u| + |v|$. 
	This yields that the only prefix of a word from~$L$
	ending with $ab^i$ is $uab^i$.
	Thus $y uab^i \in \pref{L}$
	gives $y = \varepsilon$.
\end{proof}

\begin{proposition} \label{constrsync:letterbounded:lem:NPhard}
	Suppose that $a, b, c \in \Sigma$ with $a \ne b$ and $c \ne b$ (while $a = c$ is allowed). 
	Let $L \subseteq \Sigma^*$ be regular such that
	\[
	uab^i(b^p)^*cv
	\subseteq
	L
	\subseteq 
	x^* uab^i(b^p)^*cv y^*
	\]
	for $x, y \in \Sigma$
	with $i > |u| + |v|$ and $p > 0$.
	Then $\constrsync{L}$ is \NP-hard.
\end{proposition} % Todo: Bild zu der Konstruktion.
\begin{proof}
	We give a reduction from \ST on unary automata,
	which is \NP-complete by \cref{compl:thm:comm_ST}.
	Set $I = \Sigma^* U$ for $U = uab^i(b^p)^*cv$.
	For short, in this proof we write $[w]$ for $[w]_{\equiv_I}$ and $w \in \Sigma^*$.
	
	\noindent\underline{Construction:} 
	Let $(\mathcal A, S, T)$ be an instance of \ST 
	where $\semiautomat{A}{Q}{\{b\}}{\delta}$.
	Our reduction first checks whether $S \subseteq T$,
	and if this is the case, $(\mathcal A, S, T)$
	is a \rYES-instance.
	Next, assume $S \not\subseteq T$.	
	Then pick some state $\hat q \in S$
	and let $\overline S \colon Q \to S$
	be given by $\overline{S}(q) = q$ if $q \in Q$
	and $\overline{S}(q) = \hat{q}$ if $q \notin S$.	
	Set $\semiautomat{A'}{Q'}{\Sigma}{\delta'}$
	with $Q' = (Q \cup \{t\}) \times P$ for $P = \{ [w] \mid w \in \Sigma^* \}$ and
	\[
	\delta'((q, [w]), x)
	=
	\begin{cases}
		(\delta(q, b), [wx]) & \text{if } wx \equiv_I uab^i, \\ %\text{ and } w \equiv_I uab^{i + p - 1}, \\
		(t, [wx])            & \text{if } w  \equiv_I uab^i \text{ and } x = c \text{ and } q \in T, \\
		(\overline{S}(q), [wx])	& \text{otherwise.}
	\end{cases}
	\]
	By~\cref{prelim:eqn:Nerodecongruence}, $\delta'$ is well-defined.
	Note that $P = \{ [w] \mid w \in \pref{uab^icv} \}$
	by \cref{prelim:formlang:cor:left-ideal}
	and since $w \in \pref{U}$ with $w = uab^i w'$ for some $w' \in \Sigma^*$
	is $\equiv_I$-equivalent to each word from $uab^i (b^p)^*w'$.	
		
	\noindent\underline{Verification:} Next, we show that $\mathcal A'$
	has a synchronizing word in $L$ if and only if
	there exists $j \ge 0$ such that $\delta(S, b^j) \subseteq T$
	in~$\mathcal A$.
	
	First, suppose $x^r u a b^{i + l\cdot p} c v y^s \in L$ for $l,r, s \ge 0$
	is a word synchronizing~$\mathcal A'$.
	Since the states in $\{t\} \times P$ form an upper set,
	the synchronizing state must belong to this set.
	Let $q \in Q$.
	Consider a prefix $u' \in \pref{x^r ua b^{i-1}}$.
	Assume $u'cv \in I$.
	Then $u'$ must have a word from $uab^i(b^p)^*$ as a suffix,
	in particular, it must have $ab^i$ as a factor.
	This implies that $b^i$ must be a factor of $u$,
	which is not possible since $i > |u|$.
	Thus we have $u'cv \notin I$.
	As $uab^icv \in I$, we find $u' \not\equiv_I uab^i$.
	Next, write $x^r ua b^{i-1} = z u'$ for some $z \in \Sigma$ and $u' \in \Sigma^*$.
	By the definition of~$\delta'$,
	we then find $\delta'((q, [\varepsilon]), z) = (\overline{S}(q), [z])$.
	Since $\overline{S}^2 = \overline{S}$,
	and no prefix of $x^ruab^{i-1}$ is $\equiv_I$-equivalent to $uab^i$,
	we obtain $\delta'((\overline{S}(q), [z]), u') = (\overline{S}(q), [zu'])$.
	By \cref{constrsync:letter-bounded:lem:uabi} 
	we have $\sigma_U(zu'b) = uab^i$.
	With \cref{prelim:formlang:prop:left-ideal} this gives
	$zu'b \equiv_I uab^i$.
	By the definition of~$\delta'$, we then obtain
	\[
	\delta'((\overline{S}(q), [zu']), b)
	= (\delta(\overline{S}(q), b), [uab^i]).
	\]
	Thus, for $q \in S$,
	we find $\delta'((q, [\varepsilon]), z^r u a b^i)) = (\delta(q,b), [uab^i])$.
	Next, we have 
	\[
	\delta'((\delta(q,b), [uab^i]), b^{l\cdot p}) = (\delta(q, b^{l+1}), [uab^i]).
	\]
	Now, assume $\delta(q, b^{l+1}) \notin T$.
	Then, since $uab^ic \not\equiv_I uab^i$ and by the definition of $\delta'$,
	we find	
	\[ 
	\delta'((\delta(q, b^{l+1}), [uab^i]), c) = (\overline{S}(\delta(q, b^{l+1})), [uab^ic]).
	\]
	Since the state $(\overline{S}(\delta(q, b^l)), [uab^ic])$
	is mapped to a state in $\{t\} \times P$ by $vy^s$, by the definition of~$\delta'$
	we must have $uab^ic w \equiv_I uab^i$ for some $w \in \pref{vy^s}$.
	By \cref{constrsync:left-ideal-uvw}
	we have $uab^icw \not\equiv_I uab^i$ for $w \in \pref{v}$.
	Thus, there exists $j > 0$ such that $uab^icvy^j \equiv_I uab^i$ where we 
	choose~$j$ to be minimal.
	This gives $y = b$, since for $y \ne b$
	we have $uab^icvycv\notin I$ since the word has no suffix from~$U$, 
	but $uab^icv\in I$.	
	Next, by the definition of~$\delta'$,
	since $uab^ic w \not\equiv_I uab^i$ for $w \in \pref{v}$, we have
	\[
	\delta'((\overline{S}(\delta(q, b^{l+1})), [uab^ic]), v) = (p, [uab^icv]))
	\]
	for some $p \in S$.
	Then with $y = b$ we find
	\[
	\delta'((\overline{S}(\delta(q, b^{l+1})), [uab^ic]), vy^s) \in Q \times \{ [uab^icvy^s] \}.
	\]
	Thus, the state $(\overline{S}(\delta(q, b^l)), [uab^ic])$
	cannot be mapped into $\{t\} \times P$.
	Hence, we have $\delta(q, b^{l+1}) \in T$
	for all $q \in S$.	
	
	Conversely, suppose there exists $l \ge 0$ such that $\delta(q, b^l) \in T$
	for all $q \in S$. Since $S \not\subseteq T$, we have $l > 0$.
	Consider the word $u a b^{i + (l-1)\cdot p} c v$ from~$L$.
	Let $q \in Q$ and $[w] \in P$ with $w \in \pref{uab^icv}$.
	By \cref{constrsync:letter-bounded:lem:uabi},
	we have $\sigma_U(wuab^i) = uab^i$.
	With \cref{prelim:formlang:prop:left-ideal} 
	we get $wuab^i \equiv_I uab^i$.
	If we have
	$\delta'((q, [w], uab^i) = \delta'((q, [\varepsilon]), wuab^i) = (t, [uab^i])$,
	then $\delta'((q, [w], uab^icv) = (t, [uab^icv])$.
	Otherwise, similarly as above, we find
	\[
	\delta'((q, [w], uab^i) \in \delta(S, b) \times [uab^i].
	\]
	Next, we have	
	\[ 
	\delta'(\delta(S, b) \times [uab^i], b^{(l-1)\cdot p}) \in \delta(S, b^l) \times [uab^i]).
	\]
	Since $\delta(S, b^l) \subseteq T$,
	we find $\delta'( \delta(S, b^l) \times [uab^i]), c) \subseteq \{t\} \times \{[uab^ic]\}$.
	Since $\delta'((t, [uab^ic], v) = (t, [uab^icv])$,
	we find that $(q, [w])$ is mapped to $(t, [uab^icv])$
	by $u a b^{i + (l-1)\cdot p} c v$.
	Since we also have $\delta'((t, [w]), u a b^{i + (l-1)\cdot p} a v) = (t, [uab^icv])$
	for each $w \in \Sigma^*$,
	we find that $u a b^{i + l\cdot p} c v$
	maps all states of~$\mathcal A'$
	to $(t, [uab^iav])$.
\end{proof}



\section{Remarks on Constraint PFDAs}
\label{constr_sync:sec:aut}

%!TEX root = ../../thesis-trier.tex


In~\cref{constrsync:sec:constraut:startfinal} we show that, in a constraint automaton, 
we have some degree of freedom on how to distribute the start state and the final states. 
Later, in some proofs, this allows us to choose those states 
in a way that fits best in a given situation.
%This comes in handy, as
%it allows us to choose a set of final states that fits best in a given situation.
In~\cref{constrsync:sec:constraut:concatlowerbounds} 
we remark how to reformulate some results from~\cref{constr_sync:subsec:concat} 
in terms of PDFAs.
Various conditions on constraint PDFAs implying
that the problem is equivalent to $\constrsync{\Sigma^*}$
are collected in~\cref{constrsync:sec:constraut:unconstrained}. 


\subsection{Distributing the Start State and the Final States}
\label{constrsync:sec:constraut:startfinal}


By~\cref{constr_sync:cor:factors_sets_included,gen_results:cor:factors_etc_complexity} the complexity
of the constrained synchronization problem depends on the set of factors of the constraint 
language. \Cref{constr_sync:prop:final_states}
concerns how we can distribute a start state and a set
of final states in a trim PDFA without changing the factor set of the accepted language.
Various conclusions from \cref{constr_sync:prop:final_states} 
on altering the start state and set of final states
while preserving the computational complexity are collected in \cref{constr_sync:move_states}.
Lastly, in \cref{constr_sync:prop:start_states}, 
we prove that for trim PDFAs, considering multiple starting
states does not change the complexity.

%Recall the definition of the reachability quasiorder $\lereach$ and the induced partial order on the strongly connected components
%from~\cref{prelim:sec:reach} which extends without difficulty to PDFAs.

\begin{proposition}
	\label{constr_sync:prop:final_states}
	Suppose $\automat{B}{P}{\Sigma}{\mu}{p_0}{F}$ is a trim PDFA. Let $p \in P$
	and $E \subseteq P$ be such that
	$p_0$ is accessible\footnote{As $\mathcal B$
		is assumed to be trim,  $p \lereach p_0$, i.e., $p_0$ being accessible in $\mathcal B_{p, E}$, is equivalent ot $p_0 \eqreach p$.}
	in $\mathcal B_{p, E}$
	and every state in $F$ is coaccessible in $\mathcal B_{p, E}$.
	Then $\fact{L(\mathcal B)} = \fact{L(\mathcal B_{p, E})}$.
\end{proposition}
\begin{proof} % maximal bedingung weg, nur corollar; und trim annahme weg?
	Let $u \in L(\mathcal B)$. 
	As $p_0$ is accessible in $\mathcal B_{p, E}$,
	there exists $v \in \Sigma^*$ such that 	
	$\applytrans{\mu}{p}{v} = p_0$.
	Furthermore, as $\applytrans{\mu}{p_0}{u} \in F$,
	and each state in $F$ is coaccessible in $\mathcal B_{p,E}$ 
	there exists $w \in \Sigma^*$
	such that $\applytrans{\mu}{\applytrans{\mu}{p_0}{u}}{w} \in E$.
	Hence $vuw \in L(\mathcal B_{p, E})$.
	
	Now, suppose $u \in L(\mathcal B_{p, E})$.
	As $\mathcal B$ is trim,
	there exist $v, w \in \Sigma^*$
	such that $\applytrans{\mu}{p_0}{v} = p$
	and $\applytrans{\mu}{\applytrans{\mu}{p}{u}}{w} \in F$.
	So $vuw \in L(\mathcal B)$.
\end{proof}

Next, using \cref{constr_sync:cor:factors_sets_included,gen_results:cor:factors_etc_complexity}, we state some consequences of~\cref{constr_sync:prop:final_states} regarding changes to the start state or modifications to the final state set, while preserving the computational complexity of the associated constrained problem.

\begin{corollary}
	\label{constr_sync:move_states}
	Suppose $\automat{B}{P}{\Sigma}{\mu}{p_0}{F}$
	is a trim PDFA.
	\begin{enumerate} 		
		\item \label{constr_sync:move_states:1} $\constrsync{L(\mathcal B)} \eqACZEROm 
		\constrsync{L(\mathcal B_{s, F \cup \{q\}})}$
		for every $s \in [p_0]_{\eqreach}$
		and $q \in P$.\footnote{Note that, as $\mathcal B$ is trim, we have $F \ne \emptyset$ and~$p$ is coaccessible.}
		
		
		\item \label{constr_sync:move_states:2} For $q \in P$ coaccessible in $\mathcal B_{p_0, F\setminus\{q\}}$ 
		we have $\constrsync{L(\mathcal B)} \eqACZEROm 
		\constrsync{L(\mathcal B_{p_0, F \setminus \{q\}})}$.	
		
		\item \label{constr_sync:move_states:3} $\constrsync{L(\mathcal B)} \eqACZEROm  \constrsync{L(\mathcal B_{p_0, P})}$.
		
		\item \label{constr_sync:move_states:4} Let $E \subseteq P$ be a set such that $E \cap [q]_{\eqreach} \ne \emptyset$ for each
		maximal strongly connected component $[q]_{\eqreach}$.
		Then $\constrsync{L(\mathcal B)} \eqACZEROm 
		\constrsync{L(\mathcal B_{p_0, E})}$.
		
	\end{enumerate} 
\end{corollary}
\begin{proof} In the following we only show equality of the factor sets,
	the equivalence for $\leACZEROm$-reductions follows from~\cref{gen_results:cor:factors_etc_complexity}.	
		
\ref{constr_sync:move_states:1}: By~\cref{constr_sync:prop:final_states} applied to $p = p_0$
	and $E = F \cup \{q\}$.	
	
\ref{constr_sync:move_states:2}:   By~\cref{constr_sync:prop:final_states} applied to $p = p_0$	and $E = F \setminus \{q\}$.
	
	
\ref{constr_sync:move_states:3}:  By~\cref{constr_sync:prop:final_states} applied to $p = p_0$
	and $P$ where, as $F \subseteq P$, every state in $F$ is coaccessible in $L(\mathcal B_{p_0, P})$.
	
\ref{constr_sync:move_states:4}: By the assumption
 every state is coaccessible in $B_{p_0, E}$
 so the claim follows from~\cref{constr_sync:prop:final_states}.
\end{proof}

	\begin{figure}[htbp]
	\centering	
	\begin{minipage}{0.2\textwidth}
		\begin{tikzpicture}[>=stealth,shorten >=1pt,auto, node distance = 1.8cm]
			\tikzset{every state/.style={minimum size=1pt,fill=white}}
			
			\node[state, initial above] (q0) {};
			\node[state, right = of q0] (q1) {};
			\node[state, below = of q1] (q2) {};
			\node[state, accepting, left  = of q2] (q3) {};
			
			\path[->] (q0) edge[bend left] node {$a$} (q1)
			(q1) edge[bend left] node {$b$} (q0)
			(q1) edge node {$a$} (q2);
			
			\path[->] (q2) edge [bend left] node {$a$} (q3)
			(q3) edge [bend left] node {$b$} (q2);
		\end{tikzpicture}
	\end{minipage}\hfill
	\begin{minipage}{0.2\textwidth}
		\begin{tikzpicture}[>=stealth,shorten >=1pt,auto, node distance = 1.8cm]
			\tikzset{every state/.style={minimum size=1pt,fill=white}}
			
			\node[state] (q0) {};
			\node[state, initial above, accepting, right = of q0] (q1) {};
			\node[state, below = of q1] (q2) {};
			\node[state, accepting, left  = of q2] (q3) {};
			
			\path[->] (q0) edge[bend left] node {$a$} (q1)
			(q1) edge[bend left] node {$b$} (q0)
			(q1) edge node {$a$} (q2);
			
			\path[->] (q2) edge [bend left] node {$a$} (q3)
			(q3) edge [bend left] node {$b$} (q2);
		\end{tikzpicture}
	\end{minipage}\hfill
	\begin{minipage}{0.2\textwidth}
		\begin{tikzpicture}[>=stealth,shorten >=1pt,auto, node distance = 1.8cm]
			\tikzset{every state/.style={minimum size=1pt,fill=white}}
			
			\node[state, initial above] (q0) {};
			\node[state, right = of q0] (q1) {};
			\node[state, accepting, below = of q1] (q2) {};
			\node[state, left  = of q2] (q3) {};
			
			\path[->] (q0) edge[bend left] node {$a$} (q1)
			(q1) edge[bend left] node {$b$} (q0)
			(q1) edge node {$a$} (q2);
			
			\path[->] (q2) edge [bend left] node {$a$} (q3)
			(q3) edge [bend left] node {$b$} (q2);
		\end{tikzpicture}
	\end{minipage}\hfill
	\begin{minipage}{0.2\textwidth}
		\begin{tikzpicture}[>=stealth,shorten >=1pt,auto, node distance = 1.8cm]
			\tikzset{every state/.style={minimum size=1pt,fill=white}}
			
			\node[state, accepting, initial above] (q0) {};
			\node[state, accepting, right = of q0] (q1) {};
			\node[state, accepting, below = of q1] (q2) {};
			\node[state, accepting, left  = of q2] (q3) {};
			
			\path[->] (q0) edge[bend left] node {$a$} (q1)
			(q1) edge[bend left] node {$b$} (q0)
			(q1) edge node {$a$} (q2);
			
			\path[->] (q2) edge [bend left] node {$a$} (q3)
			(q3) edge [bend left] node {$b$} (q2);
		\end{tikzpicture}
	\end{minipage}\hfill
	\caption[Distributions of final states that give equivalent constrained problems]{The following distributions
		of final states (marked with a double circle) and start state
		all yield the same complexity in the constrained synchronization problem
		when the accepted languages are used as constraint languages.}
	\label{constrsync:fig:distributions-final-states}
\end{figure}	

To illustrate \cref{constr_sync:move_states}, 
\cref{constrsync:fig:distributions-final-states} gives different distributions of final states and a start state
for a PDFA that all yield the same computational complexity in the constrained synchronization problem
when the constraint language is the language accepted by the PDFA.

When we omit the assumption of $\mathcal B$ being trim
in \cref{constr_sync:move_states}, \ref{constr_sync:move_states:1} still holds
for every coaccessible $q \in P$, \ref{constr_sync:move_states:2} holds as well without modification, 
\ref{constr_sync:move_states:3} holds by taking
the set of all coaccessible states instead of $P$, and in \ref{constr_sync:move_states:4} we have to consider
only those maximal $[q]_{\eqreach}$ containing at least one final state.

Lastly, we state that for trim PDFAs, considering multiple starting
states does not change the complexity.

\begin{proposition} \label{constr_sync:prop:start_states}
	Let $\automat{B}{P}{\Sigma}{\mu}{p_0}{F}$ be a trim PDFA.	
	Then 
	\[ 
	\constrsync{L(\mathcal B)} \eqACZEROm \constrsync{L(\mathcal B) \cup \bigcup_{p \in I} L(\mathcal B_{p, F})}
	\]
	for all $I \subseteq P$.
\end{proposition}
\begin{proof}
	Let $u \in L(\mathcal B_{p, F})$ for some $p \in I$.
	As $\mathcal B$ is trim, there exists $v \in \Sigma^*$
	such that $\mu(p_0, v) = p$ and we get $vu \in L(\mathcal B)$.
	So $\bigcup_{p \in I} L(\mathcal B_{p, F}) \subseteq \suff{L(\mathcal B)}$
	and by \cref{constrsync:genresults:cor:subset} 
	the claim follows.
\end{proof}

\subsection{Concatenation and Complexity Lower Bounds}
\label{constrsync:sec:constraut:concatlowerbounds}

The following statement is a reformulation 
of \crefitem{constrsync:cor:UanL}{constrsync:cor:UanL:1} and 
\cref{constr_sync:prop:Gamma_suffix}
in terms of constraint PDFAs to obtain
complexity lower bounds.

\begin{proposition}
	Let $\mathcal B = (\Sigma, P, \mu, p_0, F)$
	be a PDFA and set
	\begin{enumerate}
		\item $\Gamma_1 = \{ x \in \Sigma \mid \applytrans{\mu}{p_0}{x}\text{ is undefined } \}$ and
		
		\item $\Gamma_2 = \{ x \in \Sigma \mid \forall p \in P \colon \applytrans{\mu}{p}{x} \notin F \}$.\footnote{We consider $\applytrans{\mu}{p}{x} \notin F$
			to be true also when $\applytrans{\mu}{p}{x}$ is undefined;
see also the convention from \cref{constr_sync:subsec:pfun} 
where for partial functions $f$, $f(x) \ne y$ included the case that $f(x)$ is undefined.}
	\end{enumerate}	
	If $U_1 \subseteq \Gamma_1^*$,
	$U_2 \subseteq \Gamma_2^*$ are regular,
	then $\constrsync{L(\mathcal B)} \leACZEROm \constrsync{U_1 L(\mathcal B) U_2}$.		
\end{proposition}
\begin{proof}
	As $\Gamma_1 \Sigma^* \cap L(\mathcal B) = \emptyset$
	and $\Sigma^* \Gamma_2 \cap L(\mathcal B) = \emptyset$
	the claim follows from  \crefitem{constrsync:cor:UanL}{constrsync:cor:UanL:1}
	and  \cref{constr_sync:prop:Gamma_suffix}.	
\end{proof}


\subsection{Equivalence to Unconstrained Synchronization} 
\label{constrsync:sec:constraut:unconstrained}

First, \cref{constrsync:prop:upperset} states
a rather general condition implying that the factor set
of the language accepted by a PDFA equals $\Sigma^*$.
Then, \cref{constrsync:cor:upperset}
derives several conditions implying equivalence
to $\constrsync{\Sigma^*}$.

\begin{proposition}\label{constrsync:prop:upperset}
	Let $\automat{B}{P}{\Sigma}{\mu}{p_0}{F}$ be a PDFA.
	Suppose there exists a set $S \subseteq P$ of coaccessible states
	containing an accessible state $s \in S$
	and	such that $\mu(p, x) \in S$ for each $p \in S$ and $x \in \Sigma$.
    Then $\fact{L(\mathcal B)} = \Sigma^*$.
\end{proposition} 
\begin{proof}
	Let $v \in \Sigma^*$. 
    As $s$ is accessible there exists $u \in \Sigma^*$
	such that $\mu(p_0, u) = s$.
	By assumption $\mu(p_0, uv) \in S$
	and there exists $w \in \Sigma^*$ such that $\mu(p_0, uvw) \in F$.
	So, $v \in \fact{L(\mathcal B)}$
	and we get $\fact{L(\mathcal B)} = \Sigma^*$.	
\end{proof} 

\begin{corollary}\label{constrsync:cor:upperset}
	Let $\automat{B}{P}{\Sigma}{\mu}{p_0}{F}$ be a PDFA.
	Then 
	$\constrsync{L(\mathcal B)} \eqACZEROm \constrsync{\Sigma^*}$
	and so $\constrsync{L(\mathcal B)}$ is \NL-complete
	for $|\Sigma| \ge 2$
	and \LOGSPACE-complete for $|\Sigma| = 1$
	in the following cases:
	\begin{enumerate}
		\item \label{constrsync:cor:upperset:item:S} $\mathcal B$ is trim and there exists a non-empty $S \subseteq P$ such that $\mu(p,x) \in S$ for all $p \in S$ and $x \in \Sigma$. 
		
		\item \label{constrsync:cor:upperset:item:complete} $\mathcal B$ is complete
		and every state is coaccessible.
		
		\item \label{constrsync:cor:upperset:item:trap} $\mathcal B$ contains an accessible trap state
		that is a final state, i.e., 
		a state $p \in F$ with $L(\mathcal B_{p,\{p\}}) = \Sigma^*$
		and $\mu(p_0, u) = p$ for some $u \in \Sigma^*$.
		
		\item \label{constrsync:cor:upperset:item:trimtrap} $\mathcal B$ is trim and has a trap state.     
	\end{enumerate}
\end{corollary} 
\begin{proof}
	In the following, we only point out how to obtain
	a set $S$ such that we can deduce $\fact{L(\mathcal B)} = \Sigma^*$
	from \cref{constrsync:prop:upperset}.
	Then, the claim follows
	from \cref{gen_results:cor:factors_etc_complexity}, 
    \cref{prelim:complexity:general_aut_sync,thm:SYNC-solv}.

\ref{constrsync:cor:upperset:item:S}: As $\mathcal B$ is trim, $S$ from the statement fulfills the conditions in \cref{constrsync:prop:upperset}.

\ref{constrsync:cor:upperset:item:complete}: Consider $S = P$ for \cref{constrsync:prop:upperset}.

\ref{constrsync:cor:upperset:item:trap}: Consider the set $S$ containing the accessible final trap state.
	
\ref{constrsync:cor:upperset:item:trimtrap}: A coaccessible trap state must be a final state.
 Thus, since~$\mathcal B$ is trim, each of its trap states is final
 and the claim follows from \cref{constrsync:cor:upperset:item:trap}.
\end{proof} 

\begin{corollary}
	Let $\automat{B}{P}{\Sigma}{\mu}{p_0}{F}$ be a trim PDFA.
    If for each $u\in \Sigma^*$ we have $\mu(p, u) = q$ for some $p,q \in P$,
    then  $\constrsync{L(\mathcal B)} \eqACZEROm \constrsync{\Sigma^*}$.
\end{corollary}
\begin{proof}
	Let $u \in \Sigma^*$. Then  $\mu(p, u) = q$ for some $p,q \in P$.
	Since~$\mathcal B$ is trim,
	we have $\mu(p_0, v) = p$
	and $\mu(q, w) \in F$ for some $v,w \in \Sigma^*$.
	Thus $u \in \factor{L(\mathcal B)}$,
	which gives $\factor{L(\mathcal B)} = \Sigma^*$.
	So, the claim follows
	with \cref{gen_results:cor:factors_etc_complexity}, 
	\cref{prelim:complexity:general_aut_sync,thm:SYNC-solv}.	
\end{proof}

\begin{remark}\label{rem:PDFA-for-constraint}
	For a trim automaton $\mathcal B$ with a totally defined transition function,
	by~\crefitem{constrsync:cor:upperset}{constrsync:cor:upperset:item:complete}
	we have $\constrsync{L(\mathcal B)} \eqACZEROm \constrsync{\Sigma^*}$.
	Hence, to describe constraint languages giving other complexities than the unconstrained problem,
	automata with totally defined transitions functions have useless states.
	So, out of this reason, we opted to use PDFAs to specify constraint languages
	as a more compact description than automata 
	with a totally defined transition function.
\end{remark}

\section{Discussion}
\label{constr_sync:sec:gen_results:discussion}
	
%!TEX root = ./../../thesis-trier.tex

Here, we discuss the papers~\cite{FernauGHHVW19,DBLP:conf/fct/Hoffmann21,DBLP:conf/cocoon/Hoffmann21,HoffmannCocoonExtended} 
and how some of the result stated therein relate to the general results stated here.

Much of the general results as complexity-preservation for solid codes, the ideal intersection theorem,
sufficient conditions for containment in $\AC^0$, \LOGSPACE, or \NL,
and general \PSPACE-hardness conditions were first presented in~\cite{HoffmannCocoonExtended}. The 
\PSPACE-hardness conditions were proved differently in \cite{HoffmannCocoonExtended}
by using the equivalence relation $\eqLn$ and a general automaton constructions called \emph{automata unfolding}.
The method of proof presented here is more general. We gave proofs using another
equivalence relation arising from an abstract construction, which included $\eqLn$ as a special case.

Fernau et al.~\cite{FernauGHHVW19} were the first that
 systematically investigated the constrained synchronization problem on its own (the paper~\cite{Gusev2012} that introduced the constrained synchronization problem only used it as an auxiliary tool).
Next, we discuss some general results from~\cite{FernauGHHVW19} and their relation to the results from the present section.
For the convenience of the reader, we first restate each result from~\cite{FernauGHHVW19} before discussing it here.\footnote{In~\cite{FernauGHHVW19} the results were stated with $\lelogm$-reductions but the proofs of the cited statements actually give $\leACZEROm$-reductions.}
All languages in the statements are assumed to be regular.

\begin{enumerate}
	\item \cite[Thm.~15]{FernauGHHVW19}: Let $L\subseteq L'\subseteq\Sigma^*$. If  $L'\subseteq\{\,v\in \Sigma^*\mid \exists u,w \in \Sigma^*: uvw\in L\,\}$, then $\constrsync{L} \equiv^{\log}_m \constrsync{L'}$.
	
	This statement was generalized in \cref{constr_sync:cor:factors_sets_included,gen_results:cor:factors_etc_complexity}.
		
	\item \cite[Thm.~27]{FernauGHHVW19}:
	Let $L\subseteq\Gamma^*$ and $\varphi : \Sigma^* \to \Gamma^*$
	be a morphism such that $\varphi(\varphi^{-1}(L))=L$. Then
	$\constrsync{L} \leq^{\log}_m \constrsync{\varphi^{-1}(L)}$.
	
	Theorem~\ref{thm:hom_lower_bound_complexity} of the present work is a slight (but proper) generalization,
	i.e., Theorem~\ref{thm:hom_lower_bound_complexity} implies \cite[Thm.~27]{FernauGHHVW19}
	but not vice versa.
	For suppose $L \subseteq \Gamma^*$ is given and
	a morphism $\varphi : \Sigma^* \to \Gamma^*$
	such that $\varphi(\varphi^{-1}(L)) = L$.
	Set $U = \varphi^{-1}(L)$.
	Theorem~\ref{thm:hom_lower_bound_complexity}
	then yields
	\[
	\constrsync{\varphi(U)} \le_m^{\log} \constrsync{U}.
	\]
	So, as $\applyfunc{U}{\varphi} = L$, it is
	$\constrsync{L}  \le_m^{\log}
	\constrsync{\varphi^{-1}(L)}$.
	However, \cref{thm:hom_lower_bound_complexity} 
	is more general than~\cite[Thm.~27]{FernauGHHVW19} (even 
	when stated for $\leACZEROm$-reductions).
	For $L \subseteq \Sigma^*$ and
	a morphism $\varphi \colon \Sigma^* \to \Gamma^*$, in general,
	we only have $\varphi(\varphi^{-1}(L)) \subseteq L$,
	but not equality.
	We can apply~\cite[Thm.~27]{FernauGHHVW19}
	to $\varphi(L)$ as $\varphi(\varphi^{-1}(\varphi(L))) = \varphi(L)$,
	but this just gives
	$
	\constrsync{\varphi(L)} \le_m^{\log} \constrsync{\varphi^{-1}(\varphi(L))}.
	$
	
	\todo[inline]{makros nichts benuttz!}
	\item \cite[Thm.~31]{FernauGHHVW19}: (We have adjusted this statement slightly to conform to our notations and conventions.) 
	Let $\mathcal{B} = (P^{\mathcal{B}}, \Sigma^{\mathcal{B}}, 
	\mu^{\mathcal{B}}, p_0^{\mathcal{B}}, F)$ and 
	$\mathcal{C}=(P^{\mathcal{C}}, \Sigma^{\mathcal{C}},
	\mu^{\mathcal{C}}, p_0^{\mathcal{C}}, \emptyset)$ be PDFAs with
	$P^{\mathcal{B}} \cap P^{\mathcal{C}} = \emptyset$.  For\footnote{Here, 
		a transition function $\nu$ is identified with the set of tuples $\{ (p, x, q) \mid \applytrans{\nu}{p}{x} = q \}$.} $p_x \in
	P^{\mathcal{C}}$ let $\nu \subseteq \{p_x\} \times
	(\Sigma^{\mathcal{B}} \cup \Sigma^{\mathcal{C}}) \times
	\{p_0^{\mathcal{B}}\}$
	be such that the following three conditions are satisfied:
	\begin{enumerate}	
		\item $(p_x, x, p) \in \mu^{\mathcal C}$ for some $x \in \Sigma^{\mathcal C}$ and $p \in P^{\mathcal C}$
		implies $(p_x, x, p_0^{\mathcal B}) \notin \nu$,\footnote{In the original formulation
			this assumption is missing and instead it is equivalently
			assumed that $\mathcal B'$ is deterministic, but we do not introduce nondeterministic automata here
			and so reformulated it.}
		
		\item the automaton $\mathcal{C}'= (P^{\mathcal{C}} \cup\{p_0^{\mathcal{B}}\}, \Sigma^{\mathcal{B}} \cup \Sigma^{\mathcal{C}}, \mu^{\mathcal{C}} \cup \nu \cup \{p_0^{\mathcal{B}}\} \times (\Sigma^{\mathcal{B}} \cup \Sigma^{\mathcal{C}}) \times \{p_0^{\mathcal{B}}\})$
		
		is carefully synchronizing, and 
		\item there exists a synchronizing\footnote{This is the original formulation, not stating
			whether an exactly synchronizing word or a carefully synchronizing word is meant
			here (given assumption~(a), it probably means carefully synchronizing word).}
		word $v = v_1 \dots v_n$ [where $v_i \in \Sigma$, $i \in \{1,\ldots,n\}$]
		for $\mathcal{C}'$ such that $v_1 \dots v_{n-1} \in L(\mathcal{C}_{p_0^{\mathcal C}, \{p_x\}})$.
	\end{enumerate}
	Then $\constrsync{L(\mathcal B)} \leq_m^{\log} \constrsync{L(\mathcal B')}$
	for the PDFA\footnote{As often done, functions are viewed here as sets of tuples.}		
	$\mathcal{B'}=(P^{\mathcal{B}} \cup P^{\mathcal{C}}, \Sigma^{\mathcal{B}}
	\cup \Sigma^{\mathcal{C}}, 
	\mu^{\mathcal{B}} \cup \mu^{\mathcal{C}} \cup \nu,
	p_0^{\mathcal{C}}, F)$.
		
	This statement is covered by~\cref{constr_sync:thm:prefix_before}, which can be seen 
	by considering the language accepted by $\mathcal B'$.
	Define 
	\[
	\overline{\mathcal C} = (P^{\mathcal C} \cup \{p_0^{\mathcal B}\}, \Sigma^{\mathcal B} \cup \Sigma^{\mathcal C}, 
	\mu^{\mathcal C} \cup \nu, p_0^{\mathcal C}, \{p_0^{\mathcal B}\}).
	\]
	Then $L(\mathcal B') = L(\overline{ \mathcal C})  L(\mathcal B)$
	as 
	every accepting word has to pass through the transitions in $\nu$ (they form a choke point).
	As $\overline{\mathcal C}$ has a unique final state with no outgoing transitions,
	by \crefitem{constr_sync:prop:aut_prefix_free}{constr_sync:prop:aut_prefix_free:item:trim} 
	it accepts a prefix code.
	Then \cref{constr_sync:thm:prefix_before} 
	yields $\constrsync{L(\mathcal B)} \leq_m^{\log} \constrsync{L(\overline{ \mathcal C})  L(\mathcal B)}$.
\end{enumerate} 

\Cref{constr_sync:thm:union} about the union of constraint languages
was proven for the special case of \PTIME in~\cite[Lem.~13]{FernauGHHVW19}.
The formulation for multiple complexity classes 
appeared in~\cite[Lem.~7]{DBLP:journals/tcs/Hoffmann21} (but without mentioning $\AC^0$, $\NC^1$, $\LOGSPACE$ and~$\NL$).

Hoffmann~\cite{DBLP:conf/cocoon/Hoffmann21} 
showed the Ideal Intersection Theorem (\cref{constr_sync:gen_results:thm:ideal_separation}),
the complexity reduction by morphic images (\cref{thm:hom_lower_bound_complexity}),
and \cref{constr_sync:cor:factors_sets_included}.
Furthermore, in~\cite[Thm.~4.5]{DBLP:conf/cocoon/Hoffmann21} also appeared  a general \PSPACE-hardness condition
(called the $uC$-Theorem).
However, the proof (available in the arXiv version~\cite{DBLP:journals/corr/abs-2005-05907}) 
contained a mistake (fixed in the extended version~\cite{HoffmannCocoonExtended}).
The \PSPACE-hardness conditions presented in \cref{constr_sync:sec:PSPACE-hardness}
extend the condition from~\cite{DBLP:conf/cocoon/Hoffmann21}.
Our conditions also work with infinite prefix codes, whereas
the \PSPACE-hardness condition from~\cite{DBLP:conf/cocoon/Hoffmann21} only holds for finite prefix codes.

\IN{strongly self-synchronizing code}
\label{discussion:strongly-self-sync-codes}
The fact that coding morphisms of solid codes preserve the computational complexity of constraint languages (Theorem~\ref{thm:constr_sync_hom_strongly_self_sync}) was first stated in~\cite[Thm.~17]{DBLP:conf/fct/Hoffmann21}. However, the connection to solid codes was not made in~\cite{DBLP:conf/fct/Hoffmann21}. Instead, \emph{strongly self-synchronizing codes} were introduced as comma-free codes $C$ satisfying $\pref{C} \setminus C) C \cap \Sigma^* C \Sigma^+ = \emptyset$.
In the present work, we discovered that strongly self-synchronizing codes and solid codes are in fact identical,
which follows from \cref{constr_sync:prop:selfsync_implies_avoiding,thm:solidcode}.
For this reason we do not use the term strongly self-synchronizing codes anymore.\footnote{Self-synchronizing
	code is another name for comma-free code,
	see Wikipedia~\cite{FujiwaraTonchev2013,wiki:Self-synchronizing_code}.
	Since originally introduced as self-synchronizing codes
	satisfying an additional condition,
	we called this class strongly self-synchronizing codes; without being aware that
	this class equals the class of solid codes.
	We note that the term self-synchronizing does not seem to be used consistently in the literature
	to denote the same class of codes. For instance, in~\cite{DBLP:journals/tit/FergusonR84}
	codes admitting a synchronizing word for a canonical decoder (automaton)
	are called self-synchronizing, which gives an additional reason to avoid the term strongly self-synchronizing codes.}
We note in passing that~\cite[Remark 2]{DBLP:conf/fct/Hoffmann21} giving a construction procedure for solid codes misses 
the condition $\Sigma^*\{c\} \cap X = \emptyset$, \crefitem{constrsync:ex:binarysolidcode}{constrsync:ex:constr-solid-code}
gives the correct formulation with a proof. 